\documentclass[10pt,a4paper]{article}

% ----------------- Paquetes -------------------%
\usepackage[T1]{fontenc}
\usepackage[utf8]{inputenc}
\usepackage[a4paper,margin=2.5cm]{geometry}
\usepackage{mathptmx}             
\usepackage{changepage} 
\usepackage{setspace}
\usepackage[spanish]{babel}
\usepackage[labelfont=bf,labelsep=space]{caption}
\usepackage{amssymb}
\usepackage{amsmath}
\usepackage{ebgaramond}
\usepackage{placeins}
\usepackage{etoolbox}
\usepackage{setspace}
\usepackage{parskip} 
\usepackage{tcolorbox}
\usepackage{needspace}
\usepackage{xparse}
\usepackage{ocgx2}
\usepackage{hyperref}
\usepackage{hypcap}


%%%%%%%%%%%%%%%%%%%%%%%%%%%%%%%%%%%%%%%%%%%%%%%%%%%%%%%%%%%%%%%%
% --- Fija primero tus valores globales "buenos" ---
\setstretch{1.2}                 % Interlineado global
\setlength{\parskip}{0.7\baselineskip}
\setlength{\parindent}{0pt}
\newlength{\GlobalParskip}
\setlength{\GlobalParskip}{\parskip}

\newcounter{eqnum}[subsection]
\renewcommand{\theeqnum}{\arabic{eqnum}}
\newcounter{alphaitem}
\newcounter{numitem}

%% ------------------- Recuadros----------------------%%
\newcommand{\puntosclave}[1]{%
  \begin{tcolorbox}[
    colframe=black,
    title={Puntos clave},
    before upper={\setlength{\parindent}{0.5cm}\setlength{\parskip}{0.5\baselineskip}}
  ]
  #1
  \end{tcolorbox}
}

\newcommand{\modificaciones}[1]{
  \begin{tcolorbox}[
    colback=white,
    colframe=red,
    title={Modificaciones a realizar},
    before upper={\setlength{\parindent}{0.5cm}\setlength{\parskip}{0.5\baselineskip}}
  ]
  #1
  \end{tcolorbox}
}

%% ------------------- Preguntas TEST----------------------%%
% Opciones: radio (single-choice)
\NewDocumentCommand{\OptRadio}{m m m}{%
  \noindent\CheckBox[radio,name=#1,value=#2,width=1.2ex,height=1.2ex]{}%
  \;\textbf{#2.}~#3\par
}

% Opciones: checkbox (multi-choice)
\NewDocumentCommand{\OptCheck}{m m}{%
  \noindent\CheckBox[name=#1,width=1.2ex,height=1.2ex,bordercolor={0 0 0}]{}%
  \;#2\par
}

% Pregunta single-choice (radio)
% Args: 1=id, 2=enunciado, 3=opciones, 4=solución+justificación, 5=imagen (o {}), 6=origen
\NewDocumentCommand{\PreguntaTestSingle}{m m m m m m}{%
  \par\medskip
  \noindent\textbf{#1.}~#2\par\smallskip

    % Imagen (si no es {})
    \IfBlankTF{#5}{}{%
      \begin{center}
        \includegraphics[width=0.75\linewidth]{#5}
      \end{center}
    }

    \begin{Form}
      #3
    \end{Form}

    \par\smallskip
    \switchocg{sol-#1}{\fbox{Mostrar / ocultar solución}}%
    \hfill{\footnotesize #6}\par\smallskip

    \begin{ocg}{Solución}{sol-#1}{0}
      \noindent #4
    \end{ocg}
  \par\medskip
}

% Pregunta multi-choice (checkboxes)
\NewDocumentCommand{\PreguntaTestMulti}{m m m m m m}{%
  \par\medskip
  \noindent\textbf{#1.}~#2\par\smallskip

    % Imagen (si no es {})
    \IfBlankTF{#5}{}{%
      \begin{center}
        \includegraphics[width=0.75\linewidth]{#5}
      \end{center}
    }
    \begin{Form}
      #3
    \end{Form}

    \par\smallskip
    \switchocg{sol-#1}{\fbox{Mostrar / ocultar solución}}%
    \hfill{\footnotesize #6}\par\smallskip

    \begin{ocg}{Solución}{sol-#1}{0}
      \noindent #4
    \end{ocg}
  \par\medskip
}

%% ------------------- Listas----------------------%%
    % --- Lista alfabética ---
    \newenvironment{la}
      {\begin{list}{\alph{alphaitem})}{%
          \usecounter{alphaitem}%
          \setlength{\leftmargin}{1cm}%
          \setlength{\parskip}{3pt}% 
          \setlength{\itemsep}{0pt}% ← espacio entre ítems
          \setlength{\parsep}{0pt}%  ← párrafos dentro de un ítem
      }}
      {\end{list}}
      
    % --- Lista numérica ---
    \newenvironment{lnum}
      {\begin{list}{\arabic{numitem}.}{%
          \usecounter{numitem}%
          \setlength{\leftmargin}{1cm}%
          \setlength{\parskip}{3pt}% 
          \setlength{\itemsep}{0pt}% ← espacio entre ítems
          \setlength{\parsep}{0pt}%  ← párrafos dentro de un ítem
      }}
      {\end{list}}
      
% ------------------- Imágenes----------------------%
    \graphicspath{{Img/}}% Rutas por defecto 
    % --- Imagen adaptativa ---
    % Registros de dimensión definidos una sola vez
    \newdimen\imgcap
    \newdimen\imgmin
    \newdimen\imgmax
    \newdimen\imgremain
    \newdimen\imgh
    \newdimen\imgneed
    
    \newcommand{\imagenfit}[3]{%
      \addvspace{10pt}% separación antes
      \par\begingroup
        % Parámetros de composición (asignaciones, no \newdimen)
        \imgcap=1\baselineskip            % alto estimado de título+fuente
        \imgmin=0.15\textheight
        \imgmax=0.15\textheight
        \imgremain=\pagegoal
        \ifdim\pagegoal=\maxdimen \imgremain=0pt\fi
        \advance\imgremain by -\pagetotal
        \advance\imgremain by -\imgcap
        % Decisión de altura destino
        \ifdim\imgremain<\imgmin
          \newpage
          \imgh=0.20\textheight
        \else
          \imgh=\imgremain
          \ifdim\imgh>\imgmax \imgh=\imgmax \fi
        \fi
        % Reservar espacio para evitar cortes del bloque completo
        \imgneed=\imgh \advance\imgneed by \imgcap
        \Needspace{\imgneed}
        % Cuerpo
        \noindent\begin{minipage}{\textwidth}
          \centering
          \captionof{figure}{\textemdash\ #2}\par\vspace{0.3em}% título arriba
          \includegraphics[width=\linewidth,height=\imgh,keepaspectratio]{\detokenize{#1}}\par
          \vspace{0.3em}\small\textit{Fuente: }#3% fuente abajo
        \end{minipage}%
      \endgroup
      \par\addvspace{10pt}%
    }

%------------ Bloque de RECUADRO ESPECIAL -------------%
    \newenvironment{specialblock}[1]{%
      \begin{quote} % crea sangría a izquierda y derecha
      \small % texto más pequeño
      \noindent\textbf{#1}\\[-0.3em] % título en negrita
      \rule{\linewidth}{0.4pt}\\[0.6em]}{
      \end{quote}}
      
%-------------------------- Portada -----------------------%
\newcommand{\oldtitle}[1]{\def\OldTitle{#1}}
\newcommand{\changerationale}[1]{\def\ChangeWhy{#1}}

\newcommand{\changebox}{%
\begin{tcolorbox}[title={Historial de cambios del tema}]
\textbf{Título anterior:} \OldTitle\par
\textbf{Motivo del cambio:} \ChangeWhy
\end{tcolorbox}
}

%-------------------------- Author Font -----------------------%
    \newcommand{\authorfont}[1]{{\fontfamily{EBGaramond-TLF}\selectfont #1}}

%-------------------------- Anexos -----------------------%
    \makeatletter
    \newcounter{anexo}
    \renewcommand{\theanexo}{\Roman{anexo}}
    \newcommand{\anexo}[1]{%
      \clearpage
      \phantomsection
      \refstepcounter{anexo}%
      \noindent\textbf{Anexo \theanexo.- }#1\par\medskip
      \addcontentsline{stc}{section}{Anexo \theanexo.- #1}%
    }
    \makeatother

    %-------------------------- Lista de figuras (personal) -----------------------%
\makeatletter
\newcommand{\MyListOfFigures}{%
  \clearpage
  \phantomsection
  {\centering\bfseries\Large Índice de figuras\par}%
  \medskip
  % Entrada en nuestro índice de contenido (stc)
  \addcontentsline{stc}{section}{Índice de figuras}%
  % Recuperación del listado (sin cabecera propia)
  \begingroup
    \parindent=0pt\parskip=2pt
    \@starttoc{lof}%
  \endgroup
}
\makeatother

%-------------------------- Bibliografía -----------------------%
\makeatletter
% Contador para las referencias
\newcounter{myref}
% Cabecera de la bibliografía, entra en el índice 'stc'
\newcommand{\MyBibliography}{%
  \clearpage
  \phantomsection
  {\centering\bfseries\Large Bibliografía y referencias\par}%
  \medskip
  \addcontentsline{stc}{section}{Bibliografía y referencias}%
  \setcounter{myref}{0}%
}
\newcommand{\MyBibItem}[4]{%
  \refstepcounter{myref}%
  \noindent[\themyref]\ #1.\ \emph{#2}. #3. #4\par
}
\makeatother

%--------- Establecer cuatro niveles de índice --------%
    \setcounter{tocdepth}{4}   
    \setcounter{secnumdepth}{4}% numera hasta \paragraph

%%%%%%%%%%%%%%%%%%%%%%%%%%%%%%%%

\makeatletter

    %--------- Interlineado Global --------%
    \edef\GlobalBaseLineStretch{\baselinestretch}
    
    %--------- Espaciado de seccinones --------%
    \renewcommand\section{\@startsection{section}
      {1}{\z@}
      {2.5ex \@plus 1ex \@minus .2ex} % espacio antes
      {1.5ex \@plus .2ex}             % espacio después
      {\normalfont\Large\bfseries}    % formato del título
    }
    \renewcommand\subsection{\@startsection{subsection}{2}{\z@}
      {3ex \@plus 0.8ex \@minus .2ex}
      {1.5ex \@plus .2ex}
      {\normalfont\large\bfseries}
    }
    \renewcommand\subsubsection{\@startsection{subsubsection}{3}{\z@}
      {3ex \@plus 0.5ex \@minus .2ex}
      {1ex \@plus .2ex}
      {\normalfont\normalsize\bfseries}
    }
    \renewcommand\paragraph{\@startsection{paragraph}{4}{\z@}%
    {-2ex \@plus -0.5ex \@minus -.2ex}% espacio antes
    {0.8ex \@plus .2ex}% espacio después
    {\normalfont\normalsize\itshape}}% ← cursiva en lugar de negrita

    %--------- Ecuaciones --------%   
    % Variables globales para almacenar títulos y notas
    \gdef\leq@current@section{}%
    \gdef\leq@current@subsection{}%
    \gdef\leq@last@printed@section{}%
    \gdef\leq@last@printed@subsection{}%
    
    % Comando para añadir nota (invisible en el cuerpo)
    \newcommand{\eqnote}[1]{%
      \addcontentsline{leq}{leqnote}{#1}%
    }
    
    % Comando para ecuaciones
    \newcommand{\eqblock}[2]{%
      % Verificar si necesitamos imprimir el título de sección
      \ifx\leq@current@section\leq@last@printed@section\else
        \addcontentsline{leq}{leqsection}{\leq@current@section}%
        \global\let\leq@last@printed@section\leq@current@section
        \gdef\leq@last@printed@subsection{}% Reset subsección al cambiar sección
      \fi
      % Verificar si necesitamos imprimir el título de subsección
      \ifx\leq@current@subsection\leq@last@printed@subsection\else
        \ifx\leq@current@subsection\@empty\else
          \addcontentsline{leq}{leqsubsection}{\leq@current@subsection}%
          \global\let\leq@last@printed@subsection\leq@current@subsection
        \fi
      \fi
      % Ecuación
      \refstepcounter{eqnum}%
      \begin{equation*}
        #1\tag{E.\theeqnum}
      \end{equation*}%
      \addcontentsline{leq}{leqt}{[E.\theeqnum] -- #2}%
      \addcontentsline{leq}{leqm}{#1}%
    }
    
    %%% ----- Índice de ecuaciones ----------%%%
    % Tamaño para []
    \newcommand{\leqIndexFont}{\normalsize}
    \newcommand{\leqTitleFont}{\tiny}
    \newcommand{\leqNoteFont}{\tiny}
    % Cabecera de sección 
    \newcommand*\l@leqsection[2]{%
      \addvspace{25pt}% Mayor separación antes de nueva sección
      \noindent\leqIndexFont
      \makebox[\linewidth]{\rule{.38\linewidth}{0.4pt}\ \ \textbf{  \MakeUppercase{#1}}\ \ \rule{.38\linewidth}{0.4pt}}%
      \par\vspace{-10pt}
    }
    % Cabecera de subsección 
    \newcommand*\l@leqsubsection[2]{%
      \par\addvspace{15pt}%
      \noindent\leqIndexFont #1
      \par\vspace{-7pt}
    }
    % Título de ecuación 
    \newcommand*\l@leqt[2]{%
    \par\addvspace{10pt}%
      \par\noindent\leqTitleFont #1\par
    }
    % Miniatura de ecuación
    \newcommand*\l@leqm[2]{%
      \vspace{0.3\baselineskip}%
      \par\scriptsize\noindent\hspace*{1.5em}\(#1\)\par%
    }
    % Nota de ecuación 
    \newcommand*\l@leqnote[2]{%
      \vspace{0.2\baselineskip}%
      \par\noindent\leqNoteFont\hspace*{1.5em}\textbf{Nota.--} #1\par\vspace{0.5\baselineskip}%
    }
    % Generación del índice
    \newcommand{\mylistofequations}{%
      \clearpage
      \phantomsection
      % Título visual de la lista (ajústalo a tu gusto):
      {\centering\bfseries\Large Índice de ecuaciones\par}
      % Entrada en nuestro índice 'stc':
      \addcontentsline{stc}{section}{Índice de ecuaciones}%
      % Llamada a tu lista real (usa el nombre exacto de tu macro):
      \@starttoc{leq}%
    }
    
    %--------- Captura de títulos de sección --------%
    \let\leq@original@section\section
    \let\leq@original@subsection\subsection
    % Flag para saber si estamos en una sección asterisco
    \newif\ifLeqStarSection
    \LeqStarSectionfalse
    
    % Redefinir \section CON CONTROL DE ASTERISCO
    \renewcommand{\section}{%
      \@ifstar{\leq@section@star}{\@dblarg\leq@section@nostar}%
    }
    \def\leq@section@star#1{%
      \global\LeqStarSectiontrue
      \gdef\leq@current@section{#1}%
      \gdef\leq@current@subsection{}%
      \leq@original@section*{#1}%
      \global\LeqStarSectionfalse
    }
    \def\leq@section@nostar[#1]#2{%
      \global\LeqStarSectionfalse
      \gdef\leq@current@section{#2}%
      \gdef\leq@current@subsection{}%
      \leq@original@section[#1]{#2}%
    }
    % Redefinir \subsection
    \renewcommand{\subsection}{%
      \@ifstar{\leq@subsection@star}{\@dblarg\leq@subsection@nostar}%
    }
    \def\leq@subsection@star#1{%
      \global\LeqStarSectiontrue
      \gdef\leq@current@subsection{#1}%
      \leq@original@subsection*{#1}%
      \global\LeqStarSectionfalse
    }
    \def\leq@subsection@nostar[#1]#2{%
      \global\LeqStarSectionfalse
      \gdef\leq@current@subsection{#2}%
      \leq@original@subsection[#1]{#2}%
    }

    % ----- Restauración de valores Globales de estilo ----------%
    % Flags
    \newif\ifInFootnote
    \InFootnotefalse
    \pretocmd{\@footnotetext}{\InFootnotetrue}{}{}
    \apptocmd{\@footnotetext}{\InFootnotefalse}{}{}
    % Profundidad de listas (soporta anidadas)
    \newcount\MyListDepth \MyListDepth=0
    \AtBeginEnvironment{la}{\advance\MyListDepth by 1}
    \AfterEndEnvironment{la}{\advance\MyListDepth by -1}
    \AtBeginEnvironment{lnum}{\advance\MyListDepth by 1}
    \AfterEndEnvironment{lnum}{\advance\MyListDepth by -1}
    \AtBeginEnvironment{itemize}{\advance\MyListDepth by 1}
    \AfterEndEnvironment{itemize}{\advance\MyListDepth by -1}
    \AtBeginEnvironment{enumerate}{\advance\MyListDepth by 1}
    \AfterEndEnvironment{enumerate}{\advance\MyListDepth by -1}
    % Restauración diferida: hazlo al INICIO del próximo párrafo real
    \newif\if@pendingRestore
    \@pendingRestorefalse
    \newcommand{\RequestRestore}{\global\@pendingRestoretrue}
    \everypar{%
      \if@pendingRestore
        \global\@pendingRestorefalse
        \ifInFootnote\else
          \ifnum\MyListDepth=0
            \setlength{\parskip}{\GlobalParskip}%
            \setstretch{\GlobalBaseLineStretch}%
          \fi
        \fi
      \fi
    }
    % Al final de bloques:
    \AfterEndEnvironment{equation}{\RequestRestore}
    \AfterEndEnvironment{equation*}{\RequestRestore}
    \AfterEndEnvironment{la}{\RequestRestore}
    \AfterEndEnvironment{lnum}{\RequestRestore}
    % Al final de macros:
    \apptocmd{\eqblock}{\RequestRestore}{}{}
    \apptocmd{\imagenfit}{\RequestRestore}{}{}
    
\makeatother

% ===================== ToC propio con 4 niveles (único bloque) =====================
\makeatletter

% --- Escritores: añadimos una línea al .stc en cada cabecera numerada ---
% Guardamos las versiones actuales para llamar al comportamiento normal
\let\MyTOC@orig@section\section
\let\MyTOC@orig@subsection\subsection
\let\MyTOC@orig@subsubsection\subsubsection
\let\MyTOC@orig@paragraph\paragraph

% SECTION
\renewcommand{\section}{%
  \@ifstar{\MyTOC@section@star}{\@dblarg\MyTOC@section@nostar}%
}
\def\MyTOC@section@star#1{%
  \MyTOC@orig@section*{#1}% sin entrada al .stc
}
\def\MyTOC@section@nostar[#1]#2{%
  \MyTOC@orig@section[{#1}]{#2}% comportamiento normal (escribe al .toc)
  % Escribimos también nuestra línea al .stc (texto con \numberline)
  \addcontentsline{stc}{section}{\protect\numberline{\thesection}#2}%
}

% SUBSECTION
\renewcommand{\subsection}{%
  \@ifstar{\MyTOC@subsection@star}{\@dblarg\MyTOC@subsection@nostar}%
}
\def\MyTOC@subsection@star#1{%
  \MyTOC@orig@subsection*{#1}%
}
\def\MyTOC@subsection@nostar[#1]#2{%
  \MyTOC@orig@subsection[{#1}]{#2}%
  \addcontentsline{stc}{subsection}{\protect\numberline{\thesubsection}#2}%
}

% SUBSUBSECTION
\renewcommand{\subsubsection}{%
  \@ifstar{\MyTOC@subsubsection@star}{\@dblarg\MyTOC@subsubsection@nostar}%
}
\def\MyTOC@subsubsection@star#1{%
  \MyTOC@orig@subsubsection*{#1}%
}
\def\MyTOC@subsubsection@nostar[#1]#2{%
  \MyTOC@orig@subsubsection[{#1}]{#2}%
  \addcontentsline{stc}{subsubsection}{\protect\numberline{\thesubsubsection}#2}%
}

% PARAGRAPH
\renewcommand{\paragraph}{%
  \@ifstar{\MyTOC@paragraph@star}{\@dblarg\MyTOC@paragraph@nostar}%
}
\def\MyTOC@paragraph@star#1{%
  \MyTOC@orig@paragraph*{#1}%
}
\def\MyTOC@paragraph@nostar[#1]#2{%
  \MyTOC@orig@paragraph[{#1}]{#2}%
  % Si no numeras los \paragraph, cambia \theparagraph por texto plano sin \numberline
  \addcontentsline{stc}{paragraph}{\protect\numberline{\theparagraph}#2}%
}

% --- Impresor del índice propio (lee el .stc) ---
\newcommand{\MyTOC}{%
  \par\bigskip
  {\centering\bfseries\Large Índice de contenido\par}%
  \medskip
  \begingroup
    \parindent=0pt\parskip=2pt
    \@starttoc{stc}%
  \endgroup
}

\makeatother
% ===================== fin ToC propio =====================


%%%%%%%%%%%%%%


% --- Datos del tema ---
\title{Tema 3.A.4 -- El pensamiento económico de \authorfont{KEYNES}\\
\large Formalización y comparación con el modelo neoclásico. Referencia a la economía post-keynesiana y el desequilibrio}
\author{Marcelo Ortega}
\date{Noviembre 2025}
\newcommand{\TemaCode}{3A04}

\oldtitle{Tema 3.A.4. El pensamiento económico de Keynes. Formalización y comparación con el modelo neoclásico. La síntesis neoclásica}
\changerationale{Se desplaza al siguiente tema la Síntesis Neoclásica, para que cuente con un mayor tiempo de exposición. Se remplaza esta escuela por la otra línea de autores que se reivindicó \textit{keynesiana}: los Post-Keynesianos, a los que se debería de dedicar menos espacio en la exposición que al propio Keynes. Es importante no reducir a los autores post-keynesianos a los keynesianos del desequilibrio (que, de hecho, no son especialmente relevantes dentro de esta corriente) y ser consciente de que autores especialmente relevantes son otros como Harrod, Robinson, Minsky o Kalecki.}


\begin{document}


\maketitle
\changebox

\newpage

% --- Página de resumen y modificaciones ---
\puntosclave{ %% Escribir puntos clave Apuntes CECO
    \textbf{IMPORTANTÍSIMO.-} 
    
    \begin{specialblock}{Qué es la Teoría Neoclásica y cómo se desarrolla la crítica de KEYNES}
        No existe como tal (en la época de KEYNES) un marco completo y unificado de lo que hoy conocemos como \textit{Teoría Neoclásica}. Esto conduce, inevitablemente, a que todo el desarrollo de KEYNES sea (y es) a la vez una labor reflexiva sobre lo que la Teoría Clásica (como el la llama) significa, principalmente en relación a los ciclos y los mecanismos de ajuste; y a la vez una crítica sobre esta misma reflexión elaborada. Esta aproximación es la que seguiremos aquí, por eso mismo es probable que sea difícil de seguir a veces ya que introducimos simultaneamente la reflexión de KEYNES sobre la Teoría clásica y su crítica. 
    
        Este hecho es completamente visible si uno ojea la Teoría General de KEYNES, ya que, en cada concepto introductorio que resulta relevante para explicar su propuesta en primer lugar hace un ejercicio de refexión (al que normalmente dedica uno o dos Capítulos) sobre \textit{¿Qué es ... en la Teoría Clásica?}. 
        
        Por tanto, podemos decir que KEYNES elabora una crítica comentada \textit{ad hoc} sobre la Teoría clásica, cuyos mecanismos de ajuste e interpretaciones unifica y a la vez critica en su libro. 

        Lo curioso es que ningún (o no muchos) autor se centra en criticar su interpretación acerca de la Teoría clásica (lo cual es muy interesante pues parece haberla expuesto entonces genuinamente bien; cosa que no es de extrañar teniendo en cuenta que él se formó en estas Teorías y las enseñío durante muchos años, como deja explícitamente claro) sino que centrarán su crítica sobre las propuestas/críticas que él mismo expone. Interesante.
    \end{specialblock}
    }
\vspace{1cm}

\modificaciones{
    Falta pulir la introducción y añadir la conclusión
    }

\newpage
\MyTOC
\newpage

\mylistofequations
\clearpage

\MyListOfFigures
\clearpage


\pagenumbering{arabic}
\setcounter{page}{1}


\section{Introducción}\label{intro}
En primer lugar, podríamos preguntarnos porque es importante estudiar \textbf{X} desde un punto de vista histórico. 

Para responder a esta pregunta puede resultar muy ilustrativo el Premio Nobel de Economía de 2024, en particular, el concepto de \textit{destrucción creativa}. Un autor, o autores, puede encuadrarse en una época, resultar ser el contorno entre épocas o precisamente las ideas germinales de cambio que motivan el nacimiento de otra época, pero lo que es innegable es que ese autor es fruto del conocimiento previo que la ciencia ha alcanzado.

\imagenfit{DestruccionCreativa.png}{Ilustración: Destrucción creativa}{\href{https://www.nobelprize.org/prizes/economic-sciences/2025/press-release/}{Nobel Prize Outreach 2025. Nov 2025.}}

Si visualizamos el conocemiento económico como una escalera que se ha ido desarrollando a lo largo de su breve historia, es imprescindible conocer en que momento (peldaño) nos encontramos, es decir, cuál es la frontera del conocimiento, para poder avanzar, pero no por ello es menos importante entender y conocer el camino que nos ha traído hasta aquí, puesto que este conocimiento previo es condición necesaria para la innovación.

\imagenfit{PrescriptiveKnowledgeVSPropositionalKnowledge.png}{Ilustración: Conocimiento Prescriptivo y Conocimiento Proposicional}{\href{https://www.nobelprize.org/prizes/economic-sciences/2025/press-release/}{Nobel Prize Outreach 2025. Nov 2025.}}

\subsection{Contextualización}\label{context}
En esta exposición, nos centraremos en las aportaciones de JOHN MAYNARD KEYNES (1883-1946). En particular, compararemos su pensamiento económico con la de autores contemporáneos y anteriores a él, enmarcados en la escuela neoclásica [ver tema 3.A.3]. 

También analizaremos las aportaciones de los economistas post-keynesianos (es decir, autores cuyas aportaciones se ven influenciadas por la obra de KEYNES). ‒ Para realizar este análisis necesitamos remontarnos a cómo era el estado de la ciencia económica por aquel entonces: o A finales del siglo XVIII, los economistas clásicos habían sistematizado una serie de proposiciones con sus aportaciones, y se considera, que gracias a estos autores, la economía adquiere el carácter de ciencia [ver tema 3.A.2]. o Un siglo más tarde, a finales del siglo XIX y principios del siglo XX, los economistas neoclásicos refinaron y formalizaron una serie de contribuciones (por ejemplo, en el marco de la teoría del valor) cambiando el enfoque de una aproximación de «economía política» a «economía» [ver tema 3.A.3]. ‒ Sin embargo, la Gran Depresión puso en jaque la defensa del laissez-faire (i.e. las virtudes de un sistema autorregulado defendido por estos autores). Todavía notándose los efectos de la Gran Depresión, KEYNES publica en 1936 su obra “The General Theory of Employment, Interest and Money”. Esta obra lo situará como uno de los economistas más relevantes de toda la historia del pensamiento económico. o En ella, KEYNES observa la presencia de un desempleo involuntario masivo y considera que los clásicos y neoclásicos habían sido demasiado optimistas respecto al laissez-faire y a las virtudes de un sistema autorregulado y reconoce que no se puede luchar contra la Gran Depresión con sus recetas de flexibilidad de precios, sino que es necesario estimular la demanda agregada. o El pensamiento de KEYNES supone un cambio radical en la forma de analizar la economía de los países. • Los neoclásicos, preocupados por la determinación de los precios de equilibrio tienen como punto de partida en su análisis el comportamiento maximizador de empresas e individuos. El equilibrio se logra en el precio y cantidad que iguala lo oferta con la demanda. No obstante, esta demanda tiene un papel secundario y está determinada por factores de oferta, en la medida en que el poder adquisitivo de los individuos depende principalmente de su oferta laboral, y el dinero se utiliza exclusivamente para las transacciones. • Para KEYNES, la economía se estudia desde un enfoque agregado con el consumo como variable clave en su análisis. Este autor engloba bajo el término “economistas clásicos” a la escuela de la teoría del valor-trabajo junto con las aportaciones de los neoclásicos [ver anexo A.4], y les critica precisamente la falta de interés en la demanda agregada como determinante fundamental de la situación económica de un país. Rescata así las ideas de MALTHUS y, aunque nunca formaliza sus ideas, defiende que la demanda efectiva determina la existencia de paro en un país; si es reducida, las empresas disminuirán su producción en consecuencia, y caerá su demanda de trabajo. Así, habrá desempleo involuntario por una insuficiente demanda efectiva. La demanda determina la oferta; muchos de los problemas de la economía desaparecen en la medida en que se incentive a la demanda. o En este sentido, la obra de KEYNES supone una revolución en términos de KUHN [ver tema 3.A.1], un ataque frontal al paradigma neoclásico que se materializaba en sus dificultades para ofrecer una explicación completa y convincente a los efectos de la Gran Depresión. • No obstante, debemos señalar que la revolución científica que predecía KUHN no se produjo del modo en que presagiaba dicho autor. En efecto, a pesar de las anomalías acumuladas por el paradigma neoclásico, éste no quedó reemplazado completamente por el keynesiano sino que ambos coexistieron y llegaron a quedar planteados bajo la síntesis neoclásica, un paradigma que se impuso hasta mediados de los 70. ‒ ‒ ¿Cómo se determinan las principales variables económicas de un país (PIB, empleo, etc.)? ¿Cómo podemos explicar el funcionamiento de una economía en su conjunto? ¿Cómo podemos analizar los efectos de diferentes políticas económicas? A éstas y a otras muchas cuestiones trata de responder la macroeconomía. ‒ La forma en que hoy estudiamos la macroeconomía2 viene en gran medida determinada por el pensamiento de JOHN MAYNARD KEYNES (1883-1946) y por su obra “The General Theory of Employment, Interest and Money” (1936). ‒ KEYNES aceptará algunos aspectos de la teoría neoclásica (p.ej. la demanda de trabajo), pero criticará muchos otros, algunos de ellos tan fundamentales como la premisa básica neoclásica de que los mercados se autorregulan. ‒ Así, KEYNES señaló que, ante ciertos desequilibrios, el mecanismo de precios podría no ser capaz de restablecer el equilibrio, abriendo así la puerta a la intervención del sector público, frente a la habitual desconfianza neoclásica ante tal intervención. ‒ La aparición de la obra de KEYNES marcará el inicio de una notable división en el pensamiento macroeconómico. No obstante, a mediados de los años 50, SAMUELSON realizará una famosa afirmación, al señalar que el 90 % de los economistas americanos habían dejado de ser pro- o anti- keynesianos y habían pasado a estar comprometidos con la Síntesis Neoclásica, que trataba de aunar ambos enfoques: el keynesiano para el corto plazo y el neoclásico para el largo.

\subsubsection{Relevancia}\label{importance}
SCHUMPETER (1954) aduce 4 motivos por los que considera fundamental el estudio de la historia del pensamiento económico:
1) Ventajas pedagógicas: Resulta muy instructivo no contentarse con las enseñanzas del último tratado científico sobre una materia, puesto que éste se encuentra condicionado por el autor (por su método, su personalidad, su ideología y sus preferencias), especialmente en las denominadas ciencias sociales. 2) Nuevas ideas: Este estudio constituye una fuente de inspiración al encontrarse antiguos principios que en su momento no fructificaron y que en el presente pueden orientar nuevas investigaciones y avances en el conocimiento científico. 3) Comprensión del proceder del espíritu humano: Gracias al estudio de la historia del pensamiento económico, se entiende mejor el marco institucional y el espíritu de los tiempos que explica el por qué se producen unos determinados logros en unas épocas y no en otras. 4) Desplazamiento progresivo de las fronteras de la ciencia económica:\footnote{La evolución del pensamiento económico es un proceso complejo que resulta de la interacción de varios factores. Entre ellos, destacan el contexto económico del momento, el pensamiento económico precedente y los avances en otras disciplinas como la filosofía, las matemáticas, la física o la biología. El conocimiento de la historia del pensamiento económico permite entender las raíces intelectuales del pensamiento actual, aproximarse al análisis de los fenómenos económicos pasados y actuales, y valorar la importancia de los diferentes programas de investigación.} En cada época aparecen hechos y problemas distintos que requieren sus propias soluciones y su peculiar modo de afrontarlos, y por lo tanto, el estudio de la historia del pensamiento económico nos permite entender cómo ha evolucionado el objeto y el método de la ciencia económica.

La obra de KEYNES es de gran relevancia en la literatura económica debido a múltiples factores: 
\begin{lnum}
    \item Supone el origen de la macroeconomía (p.ej. la contabilidad nacional es producto de la teoría keynesiana);
    \item Es pionero en introducir componentes psicológicos que guían a la acción humana (p.ej. ley psicológica del consumo y animal spirits); y,
    \item Supone una revolución en cuanto a las implicaciones de política económica (podemos considerar a KEYNES como el inventor de la política fiscal con finalidad de estabilización).
\end{lnum} 




\subsubsection{Historia}\label{history}


\subsubsection{Problemática}\label{problem}

¿Quién fue KEYNES? o ¿Qué aportaciones realizó? o ¿Qué diferenciaba el pensamiento de KEYNES del modelo clásico? o ¿Qué implicaciones ha tenido la obra de KEYNES? o ¿Qué formalizaciones del pensamiento de KEYNES se han consolidado? ‒ ¿Qué es la síntesis neoclásica?


\subsection{Estructura de la exposición}\label{structure}


\clearpage
\section{El pensamiento económico de KEYNES}

\puntosclave{
Keynes desarrolla su pensamiento económico en un contexto de fuerte inestabilidad,
marcado por la guerra mundial y el adevimiento de la segunda, la degradación del funcionamiento del patrón oro
y la Gran Depresión.

La Teoría General del Empleo, el Interés y el Dinero es la obra más importante de Keynes.
En ella, Keynes argumenta que la teoría económica neoclásica es incapaz de explicar el
comportamiento de una economía monetaria con incertidumbre. En contraposición,
argumenta que su Teoría General es válida para explicar el funcionamiento de la economía
en el mundo real, y que el equilibrio walrasiano con vaciamiento de mercados es un caso
excepcional.

Keynes afirma que los agentes económicos atesoran dinero para protegerse de la
incertidumbre. La preferencia por la liquidez determina el tipo de interés, que no tiene por
qué igualar el ahorro y la inversión. En consecuencia, nada garantiza el cumplimiento de
la Ley de Say.

El incumplimiento de la Ley de Say da lugar al cumplimiento del Principio de la Demanda
Efectiva, según el cual, la ocupación no viene determinada por la interacción de los
trabajadores y los empresarios en el mercado de trabajo, sino por las fuerzas que
determinan la demanda agregada de bienes y servicios.

El volumen de ocupación puede ser (y suele ser) inferior al pleno empleo, dando lugar a la
existencia de desempleo involuntario. La flexibilidad á la baja de los salarios nominales no
es un remedio adecuado para eliminar el desempleo. En efecto, este supuesto remedio
podría llegar a ser contraproducente.

La demanda agregada viene determinada, principalmente, por la propensión a consumir,
las expectativas de los hombres de negocios y el tipo de interés. Puede (y debe) utilizarse
la política económica para alterar estas variables y alcanzar el pleno empleo.

La publicación de la Teoría General tuvo una gran repercusion académica y política. En el
ámbito académico, propició la aparición de la Macroeconomía, disciplina dominada hasta
los aflos 70 por la denominada síntesis neoc/ásica-keynesiana. En el ámbito politico,
inspiró las políticas económicas contra-cíclicas, utilizadas hasta nuestros días.

}

\subsection{Contexto histórico} 
El 24 de octubre de 1929 quedó marcado en la historia de WALL STREET como Jueves Negro. Ese día, el índice Dow Janes cayó más de un 9\%, desatando el pánico entre los inversores y dando inicio al crash del 29. En poco tiempo, la crisis bursátil derivó en la quiebra del sistema bancario y el inicio de un largo periodo recesivo conocido como Gran Depresión. Las secuelas de esta crisis se notaron hasta poco antes del inicio de la Segunda Guerra Mundial. Según los datos de LEBERGOTT, la tasa de paro de Estados Unidos pasó del 3,2\% en 1929 al 23,6\% en 1932.

La teoría neoclásica del mercado de trabajo explica el desempleo como una consecuencia, principalmente, del poder de los trabajadores para fijar un salario real excesivo. Ciertamente, los salarios reales crecieron en Estados Unidos un 14,6\% entre 1929 y 1932 (Margo, 1993), algo que explica fácilmente el incremento del desempleo. Sin embargo, parece poco convicente que dicho incremento del salario real tuviese su origen en un aumento repentino del poder de negociación de los sindicatos. De hecho, los salarios nominales cayeron fuertemente en casi todas las industrias entre 1929 y 1932 (National Bureau of Economic Research, 1933) y los desempleados no contaban con seguro por desempleo. 

Cabe preguntarse, entonces, ¿por qué aumentaban los salarios reales cuando los trabajadores aceptaban reducciones de salarios nominales? La respuesta se encuentra en el desplome del nivel general de precios, superior en proporción a la reducción de los salarios nominales. Dicha reducción de precios, sin embargo, era difícil de explicar en el marco de la economía neoclásica. Bajo un sistema monetario de patrón oro, una caída drástica y simultánea de la producción y de los precios solo era posible a través de un desplome repentino de la velocidad de circulación del dinero\footnote{Décadas más tarde. Milton Friedman y Anna Schwartz argumentarían que la causa principal de la fuerte
caída de la producción y los precios no se encuentra en la velocidad de circulación del dinero. Según sus
investigaciones. la quiebra del sistema bancario desplomó el multiplicador mondario y -dado un stock de
oro- la oferta monetaria.}, que sin embargo, siempre había sido considerada una variable muy estable por los economistas neoclásicos. En este contexto, la teoría neoclásica encuentra dificultades para explicar de manera convincente las razones últimas por las qué subían los salarios reales, y por tanto, las causas finales del desempleo involuntario durante la Gran Depresión.

\subsubsection{Deficiencias en Política Económica}

\paragraph{Finales del s.XIX hasta 1929}
En las décadas anteriores, uno de los casos de atesoramiento más citados por economistas británicos era el de India bajo dominación colonial. Allí era frecuente que la expansión monetaria impulsada por las autoridades no se tradujera en un aumento del gasto/inversión, porque la población tendía a conservar moneda metálica como forma de ahorro seguro. Las \textit{currency reports} del gobierno británico en la India, desde finales del siglo XIX hasta la década de 1910, mencionaban repetidamente que el dinero adicional inyectado al sistema terminaba absorbido por el público en forma de moneda guardada y no movilizada. 

Del mismo modo, en las economías industriales también hubo episodios previos en los que la demanda de dinero ($M^d$) aumentó súbitamente. La \href{https://es.wikipedia.org/wiki/Pánico_financiero_de_1907}{crisis bancaria estadounidense de 1907} produjo retiradas masivas de depósitos y acumulación de efectivo por parte del público.\footnote{Los informes bancarios de la época registran cómo los ciudadanos y las empresas se lanzaron a guardar billetes y oro, lo que provocó una escasez de medios de pago y paralizó el crédito.\\
\href{https://elischolar.library.yale.edu/cgi/viewcontent.cgi?article=14123&context=ypfs-documents}{New York Clearing House Association, Annual Report (1907–1908)}; \href{https://fraser.stlouisfed.org/files/docs/historical/nmc/nmc_405_1910.pdf}{Interviews on the Bankinga and Currency Systems of England, Scotland,  France, Germany, Italy and Switzerland (1910)}} Sin embargo, la situación se resolvió relativamente rápido porque el sistema financiero, aunque fragmentado, no estaba combinado con un colapso general de la actividad real como el de 1929.

En el Reino Unido, la crisis de 1921, después de la Primera Guerra Mundial, también mostró patrones de atesoramiento. Tras la deflación fuerte aplicada por CUNLIFFE y la política de retorno al patrón oro, muchos agentes económicos redujeron drásticamente su gasto e incrementaron sus saldos líquidos. Los informes del Banco de Inglaterra de aquellos años hablan de un credit deadlock, una situación en la que los bancos estaban técnicamente capaces de prestar, pero la demanda de crédito era nula y las empresas preferían mantener liquidez por miedo a la inestabilidad de precios y empleo. Es importante entender que, antes de 1929, muchos responsables de política económica interpretaban estos episodios en clave casi psicológica o institucional, como comportamientos excepcionales vinculados al miedo o a la ausencia de un sistema bancario moderno. No existía todavía una teoría que permitiera ver el atesoramiento como una posibilidad estructural del capitalismo, capaz de bloquear el mecanismo clásico de ajuste. Keynes, observador atento de la historia monetaria y del sistema colonial británico, sabía que el fenómeno no era nuevo. Lo que cambió después de 1929 es que se manifestó en el corazón de las economías avanzadas, a gran escala y durante años, reduciendo a la impotencia los instrumentos de política monetaria. 

\imagenfit{Crack29.png}{Stock Market Index y Préstamos a los borkers durante el \textit{crack} de Nueva York}{WHITE, 1990}

\paragraph{Límites de la Teoría Neoclásica: Crack del 29 y década de los 30's}
A partir del colapso de 1929, el fenómeno del atesoramiento de dinero y otros fenómenos malentendidos se convirtieron en un problema práctico que los responsables de política económica observaban de forma directa y angustiosa. Las autoridades monetarias registraban retiradas masivas de depósitos, colas frente a los bancos, acumulación de efectivo en manos del público y una profunda parálisis del crédito. Lo que se debatía entonces, especialmente entre 1930 y 1933, era por qué las intervenciones monetarias tradicionales estaban fracasando. Los bancos centrales bajaban los tipos de interés, ampliaban el crédito y realizaban operaciones de mercado abierto, pero los préstamos no aumentaban, la actividad no se reactivaba y los bancos seguían acumulando reservas ociosas. 

Los responsables de política se preguntaban si estaban ante una crisis de confianza, una crisis bancaria o una crisis de exceso de ahorro, pero ninguno de esos diagnósticos encajaba del todo en la teoría dominante. Lo que estaban observando empíricamente era que, incluso con tipos de interés cercanos a cero, empresas y familias preferían conservar liquidez antes que invertir o consumir. De hecho, la Reserva Federal señalaba que, pese a las inyecciones de liquidez, buena parte del dinero salía del sistema bancario hacia cajas fuertes privadas, o quedaba como reservas inactivas en los propios bancos, sin generar crédito.\footnote{\href{https://fraser.stlouisfed.org/files/docs/publications/arfr/1930s/arfr_1931.pdf}{18th Annual Report of the Federal Reserve Board (1931)} y \href{https://fraser.stlouisfed.org/title/annual-report-board-governors-federal-reserve-system-117/1932-2490}{19th Annual Report of the Federal Reserve Board (1932)}} 

\imagenfit{EEUUInterés1929.png}{Tipos de interés en Estados Unidos durante el periodo 1926-1931}{WHITE, 1990}

En Reino Unido, el Banco de Inglaterra informaba al Tesoro de la incapacidad de la baja de tipos para detener la contracción de la actividad, porque la demanda de crédito prácticamente había desaparecido.

En varios artículos entre 1930 y 1931, KEYNES afirmaba que se enfrentaban a un proceso de parálisis monetaria: el público no gastaba, los bancos no prestaban y los empresarios no invertían. La política monetaria estaba perdiendo eficacia ante una preferencia por la liquidez que no respondía a la simple relación entre tipo de interés y demanda de dinero propuesta por la teoría neoclásica. Lo que Keynes observó fue que la economía podía entrar en un estado en el que el deseo de liquidez (deseo en buena medida motivado por la incertidumbre y el miedo) era suficientemente fuerte como para neutralizar cualquier estímulo convencional. En ese contexto, insistir en bajar los salarios o dejar que el mercado corrigiera la situación solo profundizaba la depresión, porque reducía aún más la confianza y reforzaba el impulso a guardar liquidez.

Como veremos, el tema del atesoramiento, las expectativas y la Demanda Efectiva se vuelve central desde el principio de su reforma teórica. Keynes estaba respondiendo a un problema concreto que los gobiernos veían cada día: la economía podía quedar atrapada en un equilibrio \textbf{depresivo} en el que \textbf{nadie gastaba} porque esperaba que la situación empeorara, y ese comportamiento colectivo hacía que efectivamente empeorara. \textbf{El dinero}, lejos de circular, \textbf{quedaba retenido}, y ese simple hecho rompía el mecanismo de ajuste automático defendido por los neoclásicos. 

\subsection{J.M. KEYNES}
Para facilitar la exposición podemos distinguir tres etapas en la vida y obra del autor: 
\begin{lnum}
    \item La 1ª etapa abarcaba las primeras décadas del siglo XX, hasta el Crash de 1929 (desde el Tratado sobre el sistema monetario y financiero indio (1913) hasta finales de la década de 1920);
    \item La 2ª etapa llegaba hasta el estallido de la Segunda Guerra Mundial, correspondiendo por lo tanto a los años de la Gran Depresión (desde el Treatise on money, publicado en 1930, a la General Theory (1936), incluyendo los escritos inmediatamente posteriores que defendían el enfoque propuesto en la última obra y destacaban su novedad);
    \item La 3ª etapa comenzó cuando los problemas económicos relacionados con la guerra atrajeron la atención general hacia el modelo de relaciones económicas internacionales (la última etapa era la de How to pay for the war (1940) y de las propuestas como la Clearing Union, para una nueva organización económica internacional).
\end{lnum} 

En general, los comentaristas concentran principalmente la atención en la segunda etapa, puesto que los escritos que pertenecen a este período son correctamente considerados como las principales contribuciones de KEYNES en el campo teórico. Sin embargo, de esta forma se pierde mucho, tanto en lo que se refiere a la interpretación del pensamiento de KEYNES como a su programa de acción. En particular, su contribución a la teoría de la probabilidad, con su noción de incertidumbre que merece un tratamiento específico, antes incluso de considerar sus contribuciones a la teoría del dinero y de la ocupación.

\subsubsection{Primera etapa: Inicios}
J.M. KEYNES (1883-1946) nació en Cambridge en el año 1883. Su padre J.N. KEYNES fue profesor de filosofía y lógica en Cambridge, y ampliamente conocido por ser autor de uno de los mejores tratados de metodología económica. Su madre, F. ADA BROWN, fue una escritora de éxito y reformista social, primera alcaldesa de Cambridge. Keynes tuvo dos hermanos. Se casó en 1925 con la bailarina Lydia Lopokova y murió de un ataque al corazón en 1946. 

En 1902 accedió al King's College de Cambridge y en 1905 se licenció en matemáticas. En Cambridge conoció a algunos de los economistas más influyentes de la época, como EDGEWORTH y MARSHALL. Durante su etapa universitaria, formó parte de los \textit{Apóstoles de Cambridge}, un club de debate conformado por la élite intelecutal de la Universidad de Cambridge, en el que conoció a personalidades como el filósofo SIDGWICK, el escritor B. RUSSELL, el politólogo L. DICKINSON y el filósofo G. E. MOORE. Posteriormente, y al igual que otros miembros de este club, participó muy activamente en el \textit{Grupo de Bloomsbury}, un influyente grupo de intelectuales que compartían ideas en temas sociales, políticos y filosóficos. Algunos de sus miembros más destacados fueron la escritora VIRGINIA WOOLF, el escritor E. M. FORSTER, el pintor DUNCAN GRANT y el biógrafo L. STRACHEY. Después de licenciarse en matemáticas, en 1906 KEYNES hizo los exámenes de ingreso en la Administración Pública (civil service), pero habiendo quedado en segundo lugar, tuvo que contentarse con un puesto en la India Office (el primero de la lista iba tradicionalmente al Tesoro).

Allí había poco trabajo, y KEYNES tuvo tiempo para escribir un tratado sobre el sistema monetario indio (publicado en 1913 con el título de Indian currency and finance), gracias al cual, después de un primer e infructuoso intento, obtuvo una beca en el King’s College de Cambridge donde desde aceptó el cargo de profesor ayudante (su modesto sueldo lo pagaba PIGOU de su propio bolsillo, continuando así una tradición iniciada por su maestro, MARSHALL) y, paralelamente, inició sus estudios en economía bajo la dirección de MARSHALL.\footnote{KEYNES iba a comprometerse activamente con su college durante toda su vida; elegido tesorero en 1924, llenó las arcas del college con una serie de inteligentes inversiones inmobiliarias y especulando arriesgadamente en la bolsa.} Fue también director de \textit{The Economic Journal}, desde 1911 hasta 1937. 

Durante la Primera Guerra Mundial, siguiendo el ejemplo de sus amigos de Bloomsbury, KEYNES se declaró objetor de conciencia, aunque en 1917, el ministro de finanzas ANDREW SONAR LAW le nombró jefe de la division del Tesoro británico encargada de la financiación exterior, trabajando en el negociado de exteriores del Consejo Financiero del Ministerio de Hacienda y del Tesoro\footnote{Se cuenta que fue reuniendo gran cantidad de pesetas, que vendió rápida y especulativamente consiguiendo un descenso en la cotización de la divisa española, muy favorable a los intereses económicos británicos, ya que se mantenía un intenso tráfico comercial con España por ser un país no beligerante (LEKACHMAN, 1966, p. 34).} sobre temas relacionados con la financiación del esfuerzo de guerra; con cierta falta de coherencia, que le atrajo la crítica de sus más intransigentes amigos literatos. En 1917, se le recompensó confiriéndole el título de caballero de la Orden del Baño.

En 1919 fue miembro de la delegación inglesa en la Conferencia de Paz de Versalles, pero se opuso a las reparaciones impuestas a Alemania, considerándolas una carga insostenible para la economía y la sociedad alemanas: así que dimitió y, de vuelta a Cambridge, abordó la cuestión en su exitoso libro sobre \textit{The economic consequences of the peace}.\footnote{Obra en la que argumentaba que las reparaciones de Guerra impuestas a Alemania tras perder la Primera Guerra Mundial eran impagables. Los hechos parecieron darle la razón en 1923, cuando la economía alemana se hundió en un escenario de hiperinflación, tras pagar tan solo una pequeña parte de las sanciones impuestas} Tras la Primera Guerra Mundial, a sus diversas responsabilidades académicas añadió entonces la de presidente de una compañía de seguros y, en asociación con otros, se embarcó en la especulación en bolsa, por su cuenta y en beneficio de parientes y amigos (aunque los resultados no fueron siempre felices). En 1921 Keynes publica un libro titulado \textit{Tratado sobre la Probabilidad}, en el que desarrolla su Teoría logica de la probabilidad, según la cual, la probabilidad no es un fenómeno estadístico, sino lógico.\footnote{Keynes distingue tres tipos de probabilidad:
\begin{lnum}
    \item La primera de ellas es la probabilidad cardinal, que aparece cuando conocemos la función de probabilidad o densidad de una variable aleatoria;
    \item La segunda es la probabilidad ordinal, que se da cuando no conocemos la función de probabilidad o densidad de una variable aleatoria, pero podemos ordenar los posibles valores de dicha variable según su probabilidad; y,
    \item La tercera es la probabilidad desconocida o incertidumbre irreducible, que aparece en aquellas circunstancias en las que no solo no conocemos la función de probabilidad o densidad de una variable aleatoria, sino que tampoco podemos comparar la probabilidad de dos posibles valores.
\end{lnum} \href{https://www.economicas.uba.ar/wp-content/uploads/2016/04/8-El-concepto-de-probabilidad-en-la-obra-de-Lord-Keynes.-Alberto-H.-Landro.pdf}{\textit{El concepto de probabilidad en la obra de Lord KEYNES}. A.H. LANDRO (s/n)}\\
La existencia de este tercer tipo de probabilidad y su impacto en el comportamiento económico de los agentes será un elemento clave de su Teoría General, publicada quince años después.}

En 1923, Keynes publica su "\textit{Tratado sobre la Reforma Monetaria}", en el que rechaza el restablecimiento del patron oro tras su supensión en la Primera Guerra Mundial, afirmando que se trataba de un sistema demasiado rígido. En efecto, el patron oro exigía el sacrificio de los objetivos de equilibrio interno (esto es, estabilidad de precios) para poder mantener la paridad de la moneda nacional con el oro y con el resto de monedas.\footnote{Bajo un sistema de patrón oro, el mantenimiento de la paridad de la moneda nacional requiere que la oferta monetaria varíe en sentido contrario al saldo exterior (entendido como tal la suma de las exportaciones netas de mercancías más las entradas netas de capital). En consecuencia, la oferta monetaria puede oscilar considerablemente y de manera independiente al volumen de producción de un país.\\
La teoría cuantitativa del dinero nos muestra que tal situación nos puede llevar a episodios de inestabilidad de precios.\\
Ejemplo: si el Reino Unido registraba un saldo exterior negativo de $X$ libras, ello significaría que ahora esas $X$ libras están en manos de extranjeros, que acudirían a sus respectivos bancos centrales para cambiar tales libras por su equivalente en su propia moneda nacional. A continuación, dichos bancos centrales acudirían al banco central del Reino Unido para intercambiar las $X$ libras por su respectivo valor en oro (nótese que dichos bancos centrales extranjeros necesitan ese oro para respaldar la moneda nacional que entregaron a sus residentes a cambio de las $X$ libras). El resultado es que el banco central del Reino Unido ha recibido $X$ libras (entregadas por los bancos centrales extranjeros) que no puede volver a poner en circulación (pues ha entregado el oro que las respaldaba). El resultado es una reducción de la oferta monetaria que, para un mismo nivel de producción, supondría una caída generalizada de precios.}

En 1925 el ministro de Hacienda del gobierno conservador, W. CHURCHILL restablecía el patron oro con la misma paridad de la libra existente antes de la guerra. Dicha decision fue duramente criticada por Keynes en su obra \textit{Las Consecuencias Económicas de Mr. Churchill}, (1925). Keynes se da cuenta de que la inflación sufrida por el Reino Unido durante la guerra fue muy superior a la registrada por sus principales socios comerciales. En consecuencia, el restablecimiento de la paridad al nivel anterior al de la guerra supondría una pérdida de competitividad muy significativa para el Reino Unido.\footnote{El mantenimiento de la paridad con el oro al nivel anterior a la guerra es equivalente a mantener inalterado el tipo de cambio nominal. En consecuencia, dada la mayor inflación registrada en Reino Unido durante la guerra, esta política implica una apreciación real de la economía británica y, por ende, una pérdida de competitividad-precio en el exterior.\\
Los hechos parecieron ser favorables de nuevo a Keynes cuando la caída de las ventas exteriores del Reino Unido provocó una reducción considerable de las reservas de oro en manos del Banco de Inglaterra, forzando al Reino Unido a abandonar de nuevo el patron oro en 1931. } También en 1925, habiendo dedicado una gran parte de su vida al cultivo de relaciones homosexuales, KEYNES se casó con una famosa bailarina rusa, LYDIA LOPOKOVA, a pesar de la mal disimulada oposición de sus amigos de Bloomsbury. 

\paragraph{Breve aventura política}
En el ámbito político e ideológico, Keynes estuvo muy implicado con el Partido Liberal entre 1924 y 1929. En 1925, dicta una conferencia bajo el título \textit{¿Soy un liberal?} en la escuela de verano del partido. En ella, Keynes rechaza tanto el principio hereditario\footnote{
\begin{adjustwidth}{2cm}{0pt}
\raggedright
\textit{El principio hereditario en la transmisión de riqueza y el control de los negocios es la razón por la cual el liderazgo de la causa capitalista es débil y estúpido. Está demasiado dominado por hombres de tercera generación. Nada hará que una institución social decaiga con más certeza que su apego al principio hereditario.}\\
\raggedleft{\authorfont{KEYNES}, Conferencia: “\textit{¿Soy un liberal?}”, 1925}
\end{adjustwidth}}, que asocia al Partido Conservador, como la lucha de clases\footnote{
\begin{adjustwidth}{2cm}{0pt}
\raggedright
\textit{Debería, entonces, unirme al Partido Laborista? Superficialmente eso es más atractivo. Pero mirando más de cerca, hay grandes dificultades. Para empezar, es un partido de clase, y la clase no es mi clase... Me puede influir lo que me parezca justo y de buen sentido, pero la guerra de clases me encontrará del lado de la burguesía educada}\\
\raggedleft{\authorfont{KEYNES}, Conferencia: “\textit{¿Soy un liberal?}”, 1925}
\end{adjustwidth}}, que asocia al Partido Laborista. En 1927, participa en la redacción del \textit{Libro amarillo del partido}, y en 1929 redacta el ambicioso programa de estímulo estatal titulado \textit{¡Podemos superar el desempleo!}, con el que el candidato liberal LLOYD GEORGE concurrió a las elecciones de ese mismo año. 
Tras la derrota electoral de 1929\footnote{El Partido Liberal ganó terreno en esas elecciones, pero siguió en tercera posición, por detrás de los conservadores y los laboristas.}, el Partido Liberal pierde protagonismo y Keynes se aparta de la esfera política, centrándose enteramente en el de la investigación académica. 

\subsubsection{Segunda etapa: Década de los 30's}
Como se viene diciendo, la crisis que asolaba el Reino Unido desde 1921 y que llega con fuerza a Estados Unidos con el crack del 1929 hace pensar si los pronósticos de KARL MARX sobre el capitalismo (que parecían haber sido aparentemente desmentidos tanto por la escuela neoclásica como por la experiencia histórica) podrían, a fin de cuentas, no estar tan equivocados;\footnote{“Keynes’s aim in writing The General Theory was to identify the causes of the mass unemployment that affected all developed economies in the Great Depression years. The 1930s were also a time during which Russia was witnessing strong economic results to the effect that a possible electoral victory of parties leaning toward communism (or taking power in more unorthodox ways) was a possibility that could not be discarded. In short, capitalism was in peril, both economically and politically, and Keynes realized that its survival implied important changes in its functioning.”
DE VROEY (2016)} ya que, con todo su esplendor, la teoría ortodoxa no estaba preparada para explicar la longitud y profundidad del declive económico y el desempleo masivo que trajo consigo, pues el modelo neoclásico sostenía que las desviaciones del equilibrio de pleno empleo serían temporales, y que el mercado contaba con los mecanismos para autorregularse, algo que parecía cada vez menos cierto. A pesar de esto, los neoclásicos contaban, y persistían en posibles explicaciones para esta situación de persistencia del desempleo: (i) rigideces en el sistema económico que paralizaban el mecanismo de ajuste hacia el pleno empleo; (ii) rigidez salarial por la influencia de los sindicatos, estudiada por autores como HICKS [Ver Tema 3.A.26]; (iii) poder de mercado de las grandes corporaciones, con capacidad de influir en los precios.\footnote{El modelo de competencia monopolística de CHAMBERLAIN (1933) arrojaba luz sobre este fenómeno [Ver Tema 3.A.18].} En este contexto surgen las principales obras de KEYNES, quien ofrece una perspectiva novedosa sobre el comportamiento agregado de la economía y proporciona apoyo teórico a un programa de acción pública que promueva el pleno empleo. 

En 1930, se publica su ensayo "\textit{Las Posibilidades Económicas de Nuestros Nietos}", un breve texto que ya había presentado en una conferencia dos años antes. En este trabajo imaginaba cómo serían los estándares de vida en Europa y Estados Unidos pasados cien años.\footnote{Según Keynes, para entonces la productividad sería tan elevada que bastaría una jornada laboral de tres horas diarias para mantener satisfechas las necesidades económicas elementales de la sociedad, dando lugar a un aumento considerable del ocio. En esta ocasion, Keynes parece haber acertado en pronosticar el elevado crecimiento de la renta per cápita vivida desde entonces, pero falló por completo al pronosticar una reducción tan considerable del tiempo de trabajo.} 

Durante la década de los años 30's, en un contexto marcado por el crash del 29\footnote{El crash del 1929 hace referencia a la catástrofe bursátil iniciada el 24 de octubre de 1929 (Jueves Negro) en EEUU. La situación se agravó considerable en los días 28 y 29 de octubre (Lunes Negro y Martes Negro), derivando en una crisis financiera cuyas secuelas provocaron la Gran Depresión.} y la Gran Depresión\footnote{La Gran Depresión hace referencia a la fuerte crisis económica y financiera mundial iniciada en Estados Unidos en 1929. Se trata de la depresión económica más larga, intensa y de mayor alcance geográfico del siglo XX.\\
“The main diagnosis about the crisis available to economists at the time was of “Austrian” inspiration. The crisis, the story ran, signaled a situation of over-investment and misallocation of resources, a state of affairs that required for its solution a process of ‘liquidation’, a real wage deflation, on the one hand, and some sanctioning of the firms that had engaged in wrong investment decisions, on the other. Flexibility was thus the motto. […] However, when the depression kept its course without wages deflation exerting its proclaimed effect, economists started to waver about the virtues of laissez-faire and to wonder whether, this doctrine to the contrary notwithstanding, governments should engage more actively in the economy. Thus, economists were torn between the policy conclusions following from accepted theory and gut feeling that another path should be taken. Keynes’s project was to remove this contradiction by providing a theoretical argument in favor of the gut feeling. The General Theory ensued.”
DE VROEY (2016)}, se dedica con gran intensidad a la teoría económica. En 1930 publica su \textit{Tratado sobre el Dinero}, que contiene un primer volumen dedicado a la teoría pura del dinero y otro dedicado a la teoría aplicada. En esta obra Keynes ya introduce algunos elementos que serán clave en su Teoría General, como la preferencia por la liquidez y la teoría del tipo de interés, si bien la teoría monetaria de Keynes se mantiene esencialmente enmarcada dentro de la escuela neoclásica de Cambridge. En 1932, comienza a redactar la que será su obra más importante, la \textit{Teoría General del Empleo, el Interés y el Dinero}. Los sucesivos borradores de la Teoría General fueron discutidos en el seno del \textit{Cambridge Circus}, un grupo de jóvenes economistas de Cambridge que comienza a reunirse para leer y discutir las nuevas ideas de KEYNES. Algunos miembros destacados de este grupo fueron R. KHAN, JOAN ROBINSON, AUSTIN ROBINSON, JAMES MEADE y PIERO SRAFFA. 

Finalmente, la Teoría General fue publicada en 1936, ocasionando una gran expectación y un intenso debate entre sus detractores y partidarios.\footnote{\textit{Treatise on Money} (1931) y la \textit{Teoría General} (1936), serán las dos obras más importantes, a las que se debe su fama como economista teórico. KEYNES dedicó la Teoría General a MALTHUS, autor que ya hablaba de la importancia de la demanda efectiva.\\Otras contribuciones importantes fueron los vivos y provocadores ensayos reunidos en \textit{Essays in persuasion}, y las bien documentadas e incisivas biografías reunidas en \textit{Essays in biography}. El mismo año en que se publicó la General Theory, KEYNES inauguraba en Cambridge el Arts Theatre, construido casi enteramente a sus expensas. Al año siguiente sufrió un ataque al corazón y se vio obligado a reducir su carga de trabajo.} En el Capítulo I de dicha publicación, expone claramente el objetivo fundamental de esta publicación:

\begin{adjustwidth}{2cm}{0pt}
\raggedright 
\textit{Sostendré que los postulados de la teoría clásica\footnote{Keynes no distingue entre economistas clásicos y neoclásicos. Utiliza el primero de estos términos para referirse a todos ellos.} sólo son aplicables a un caso especial, y no en general, porque las condiciones que supone son un caso extremo de todas las posiciones posibles de equilibrio. Más aún, las características del caso especialsupuesto por la teoría clásica no son las de la sociedad económica en que hoy vivimos,razón por la que sus enseñanzas engañan y son desastrosas si intentamos aplicarlas alos hechos reales.}\\
\raggedleft{\authorfont{KEYNES}, “\textit{Teoría General del Empleo, el Interés y el Dinero}”, 1936}
\end{adjustwidth}

En mayo de 1937 es víctima de una enfermedad cardíaca que lo mantiene inactivo hasta 1940, año en el que publica "\textit{Cómo financiar la guerra}". En dicha obra, Keynes estudia como financiar la guerra sin crear inflación.

\subsubsection{Tercera etapa: Segunda Guerra Mundial y proyección internacional}
Aunque, su dimisión tras el conflicto en las negociaciones de París tras la Primera Guerra Mundial le habían apartado casi definitivamente de la administración pública, al declararse la Segunda Guerra Mundial el gobierno requirió nuevamente sus servicios como asesor del ministro de Hacienda y del Tesoro. En este cargo se sumergió de nuevo en los problemas de financiación de la guerra, negociando préstamos de los Estados Unidos. 

En 1941 se incorporó al consejo de administración del Banco de Inglaterra.\footnote{También en esta ocasión se le recompensó por los servicios prestados con el título de barón Keynes de Tilton en 1942.} Durante la guerra, KEYNES ya había comenzado a redactar una serie de planes para reformar el orden económico internacional de la posguerra y antes de finalizar la Segunda Guerra Mundial se le nombró representante británico en la conferencia de Bretton Woods para el restablecimiento de un orden monetario internacional. 

KEYNES presentó el denominado \textit{Plan Keynes}\footnote{
Podemos resumir el Plan Keynes en dos ideas:
\begin{lnum}
    \item \textit{Vertiente financiera}. KEYNES propuso la creación de un nuevo banco central mundial que fuera capaz de crear liquidez. Así, este banco podría emitir moneda (bancor), que actuaría como moneda de reserva internacional. KEYNES propuso que el stock de bancor aumentara en línea con el comercio mundial. El objetivo era reducir la dependencia del dólar (dotando así de más facilidad a los países para obtener medios de pago).
    \item \textit{Vertiente monetaria.} KEYNES propuso un sistema de tipos de cambio fijos pero ajustables con bandas de ±5\% con ajuste simétrico.
    \begin{la}
        \item \textit{Tipos de cambio fijos}. La idea subyacente es que el equilibrio interno es una condición necesaria pero no suficiente para alcanzar pleno empleo. En otras palabras, se necesita equilibrio externo y para ello se necesita la estabilidad monetaria garantizada por el tipo de cambio fijo.
        \item \textit{Bandas de ±5\%}. Se necesita flexibilidad, y de este modo se permite un ajuste del tipo de cambio evitando que el ajuste sea costoso (es decir, que se realice un ajuste en la demanda interna). De hecho, KEYNES propone también utilizar el control de movimientos de capitales para evitar el ajuste vía demanda interna.
        \item \textit{Ajuste simétrico}. Persigue que no todo el ajuste recaiga en los países deudores.
    \end{la}
\end{lnum}}, por parte del gobierno británico, y los americanos el \textit{Plan White} con su propuesta del Fondo Monetario Internacional y la oferta a KEYNES de ser su primer gobernador [Ver Tema 3.B.21]. Evidentemente, cada plan defendía más los intereses de su nación y el peso específico de Estados Unidos inclinó la balanza hacia la aceptación básica del \textit{Plan White}, con las pertinentes concesiones propias de una negociación. Las dificultades financieras del Reino Unido obligaron a KEYNES a viajar varias veces a Estados Unidos para negociar préstamos. Un último viaje a ese país que tuvo que hacer poco después de la finalización de la Segunda Guerra Mundial lo realizó para asistir, como gobernador del FMI, a las inauguraciones de las instituciones acordadas en las conferencias de Bretton Woods. Unos meses más tarde, en abril de 1946, moría KEYNES en su casa de Tilton (Sussex) debido a un ataque al corazón. 

\subsubsection{Legado}
Existe una inmensa literatura sobre el pensamiento de KEYNES. Podríamos decir que no hay ningún aspecto de sus teorías que no haya experimentado una diversidad de interpretaciones; argumento también válido para la interpretación de la relación entre su análisis y la época en la que vivió. 

Muchos (de hecho, una abrumadora mayoría) coinciden en trazar un paralelismo entre KEYNES y la Gran Crisis de 1929 (o, dicho con más precisión, con la recesión en el Reino Unido después de la vuelta al patrón oro en 1925). Sin embargo, no cabe ninguna duda de que, en cualquier caso, las condiciones del paro elevado y persistente en la década de 1930, por lo menos favorecieron la difusión de las ideas keynesianas.\footnote{Algunos comentaristas destacan otros aspectos, tales como la óptica decididamente británica de KEYNES, que veía como su país retrocedía ante los Estados Unidos, tanto en los mercados de manufacturas como en los financieros, pero con mayor rapidez en los primeros, lo que parece explicar su mayor interés por estos últimos. }

En lo referente a los problemas de reconstrucción del sistema monetario internacional, KEYNES esbozó planes que tuvieran en cuenta los intereses de las monedas menos fuertes, como previsiblemente iba a ser el caso de la libra esterlina en un mundo dominado por el dólar americano. En efecto, al establecer los fundamentos de su análisis teórico, KEYNES, cuyo excepcional pragmatismo destacaron repetidamente sus contemporáneos, reflejaba los principales acontecimientos económicos de un período de tiempo que incluía la Gran Crisis, pero que no se circunscribía únicamente a la década de 1930. El compromiso básico de KEYNES era el de un economista cuya profesión coincidía con su sentido cívico. Su gran propósito era el de contribuir a un sistema reformado de capitalismo, que garantizara una justicia creciente junto con un amplio grado de libertad y eficiencia.\footnote{En realidad, el objetivo de KEYNES es teorizar acerca de la existencia de paro involuntario, un fenómeno cuya existencia es evidente, pero que no tiene cabida en teorías precedentes y, en consecuencia se necesitaba una nueva.\\
Por añadidura, KEYNES subrayó que su teoría tenía un claro fin moral: mejorar la vida de las personas. Para ello era indispensable un mejor funcionamiento de las economías de mercado, en un momento en que empezaba a cobrar fuerza el bloque comunista representado por la URSS y además en los países occidentales el marxismo estaba ideológicamente, casi, ganando la batalla. En pocas palabras, el capitalismo estaba en peligro y KEYNES se preocupaba por que las economías de mercado se repensaran para garantizar su funcionamiento adecuado y evitar crisis prolongadas como la Gran Depresión.} De hecho, tal propósito adoptaba una configuración cada vez más precisa frente al ascenso de los sistemas totalitarios como el fascismo italiano (posteriormente en España), el nacionalsocialimso alemán y el estalinismo en la Unión Soviética.\footnote{Para ver más sobre el planteamiento económico de dichos sistemas:\begin{lnum}
    \item \href{https://es.wikipedia.org/wiki/Estalinismo#En_lo_económico_y_en_lo_político}{Estalinismo económico}. Especial referencia a la planificación quinquenal.
    \item \href{https://es.wikipedia.org/wiki/Economía_en_la_Italia_fascista}{Economía en la Italia Fascista}. Hay que tener muy en cuenta que el fascismo no se puede caracterizar por tener un profundo pensamiento económico, se debe entender más desde una perspectiva socio-político con una tendencia a la autarquización y, sobretodo, al \href{https://es.wikipedia.org/wiki/Corporativismo}{Corporativismo} (sistema económico fundamentado sobre la economía pre-industrial que trata de adaptar dicho régimen al contexto industrial, con fuertes bases colusorias y oligarquicas)
    \item \href{https://es.wikipedia.org/wiki/Economía_de_la_Alemania_nazi}{Economía de la Alemania Nazi}. Resulta especialmente interesante comprender las políticas económicas de la Alemania Nazi ya que, muchos historiadores coinciden en que fue una de las causas principales del conflicto. 
\end{lnum}} De ahí su reconocimiento del fin de la ideología del liberalismo económico en su forma más simple y extrema: \textit{The end of laissez-faire}, como tituló un opúsculo publicado en 1926. Reconsiderando desde esta perspectiva el papel desempeñado por el Estado en la economía, surgía una doble necesidad: 
\begin{lnum}
    \item Por una parte, suministrar una crítica de la teoría entonces dominante, mostrando la insuficiencia de los mecanismos equilibradores del mercado libre; y,
    \item Por otra construir una teoría de la intervención del Estado y, por lo tanto, de la política económica.
\end{lnum} 

Éste es el proyecto que, con su actividad científica y su implicación directa, KEYNES se esforzó en llevar a cabo.  La publicación de la Teoría General transformó para siempre la ciencia económica y convirtió a
Keynes en uno de los economistas más influyentes del siglo XX. Podemos destacar la influencia
de Keynes en una doble vertiente:
\begin{lnum}
    \item En el ámbito académico, la publicación de la Teoría General supone la división de la ciencia
económica en dos disciplinas: la microeconomía, centrada en estudiar el comportamiento
individual de los agentes económicos, y la macroeconomía, orientada al estudio de los
agregados macroeconómicos. Keynes es considerado comunmente como el padre de la
macroeconomía.
    \item En el ámbito de la política económica, la publicación de la Teoría General supuso el abandono
del laissez:faire en lo referido a la gestión del ciclo económico. Desde entonces, las
autoridades de todo el mundo han utilizado -en mayor o menor medida- la política fiscal y la
política monetaria para suavizar los ciclos económicos.
\end{lnum}

\paragraph{Metodología}
Las contribuciones de Keynes a la ciencia económica no se limitan a sus aportaciones teóricas, Keynes también contribuye al debate metodológico en economía, rechazando los métodos empleados en las ciencias naturales, pues considera que la economía es una ciencia moral, que debe lidiar necesariamente con elementos que no son constantes ni homogéneos, como los  objetivos, las expectativas, los valores y las incerdumbres psicológicas de los agentes económicos. El comportamiento observado de los agentes económicos depende de elementos difícilmente cuantificables, comparables u observables, como las intenciones profundas o su psicología. En consecuencia, la modelización de este comportamiento requiere de introspección y juicios de valor, a diferencia de las ciencias naturales. Pese a su formación matemática, Keynes tampoco será especialmente partidario del empleo de herramientas estadísticas o econométricas para la obtención de conclusiones científicas en economía. Consideraba que las relaciones estadísticas entre variables económicas no solían ser constante en el tiempo, dudando así de la utilidad del análisis de regresión y renunciando, implícitamente, a la hipotésis ergódica en el estudio de la economía (distanciándola de las ciencias puras).

\paragraph{Objeto}
Por otro lado, se puede considerar que la escuela tiene un cuerpo doctrinal propio, que se podría sintetizar en un objeto común o parecido. 

En relación al objeto, siguiendo a DE VROEY (2016), el principal objetivo de KEYNES en la Teoría General es demostrar la existencia teórica del desempleo involuntario en un sistema de economía de mercado,\footnote{KEYNES distingue entre desempleo friccional y desempleo involuntario. Considera que el desempleo friccional estaba bien entendido, por lo que se centró en estudiar el desempleo involuntario.} cuya existencia en la realidad es innegable (en especial en los años de la Gran Depresión) y para el cual la teoría económica de la época no tiene explicación. La existencia de desempleados involuntarios implica que algunos individuos no pueden cumplir su plan de optimización (es decir, quieren trabajar por el salario del mercado, pero no pueden porque existe racionamiento en el mercado de trabajo). 

\begin{specialblock}{Nota personal.- Economía de Mercado y desmpleo involuntario}
    Es probable que la persistente creencia o fe en la inexistencia del desempleo involuntario de la escuela neoclásica venga en gran medida determinada por su contexto histórico. Recordemos que durante el siglo XIX, mediados principalmente, se está experimentando el verdadero cambio de modelo en el sistema económico, habiendo situaciones muy diversas en el Viejo Continente, desde la Inglaterra pionera en la industrialización hasta la Rusia zarista totalmente opuesta a las ideas industrializadoras con el afán de preservar los grandes privilegios de la clase aristocrática, estrechamente ligada a la propiedad de la tierra. Estos casi 150 años de intervalo que van desde la publicación de la obra de SMITH hasta la publicación de la obra de KEYNES pueden suponer, sin ánimo de aventurarme, el tiempo necesario y suficiente para entender verdaderamente las dinámicas que caracterizan al nuevo sistema económico: el capitalismo.
    
    Tal y como parece que discurrió una reunión entre ERIC WILLIAMS y su director de tésis: 
    
    \textit{(Sobre la abolición de la esclavitud en Gran Bretaña, finalmente impuesta a mediados del s. XIX..) Mucha gente piensa que los británicos mostraron buena voluntad. Había mucha gente que tenía buena voluntad, pero fue la base, la base económica, la que permitió que la buena voluntad funcionara.}
    
    ¿Qué quiero decir? En el sistema económico anterior al surgimiento del sistema capitalista, el desempleo no era, en ninguna medida, una cuestión involuntaria. La estructura socio-política debía permitir, y permitía, el desarrollo económico mínimo y suficiente como garantía de subsistencia a cualquier persona en busca de ésta. El sistema económico capitalista no tiene, ni puede tener por su condición, este tipo de ligadura socio-político. El sistema provee de empleo, pero solo en la medida en la que es rentable. No hay ninguna posibilidad consciente (más allá del altruismo voluntario y racional o las limitaciones legales) de que un empleador mantenga a su empleado en el puesto de trabajo a sabiendas de que este le genera únicamente pérdidas salariales, pues no hay ninguna manera de que este sea capaz de recuperarlas más allá del empleo de su capital. Es decir, el riesgo recae sobre el empresario. No así en el sistema económico inmediatamente anterior, donde el riesgo era repartido entre toda la sociedad puesto que la garantía de subsistencia era el seguro del empleado y el riesgo del empleador, y la superestructura social era el riesgo del empleado y el seguro del empleador.
    
    Véase: \textit{Capitalismo y esclavitud.}, ERIC WILLIAMS (1947, Ed. 2022)
    
\end{specialblock}

Con la teoría económica de la época, los economistas justificaban la existencia de desempleo mediante rigideces salariales frente a lo que KEYNES se opone argumentando que existe un fallo sistémico que afecta a las economías descentralizadas. 

En concreto, vincula el desempleo involuntario a una deficiente demanda agregada, que a su vez estaba asociada con alguna fuga del sector productivo hacia el sector financiero. Como resultado, era necesario matizar la interpretación optimista de la economía de mercado defendida por los economistas desde ADAM SMITH (más precisamente, desde RICARDO).

\paragraph{Estructura de la Teoría General}
A lo largo de la Teoría General, va levantando los supuestos que se esconden detrás del caso especial de la economía neoclásica:

\imagenfit{Neoclasica-Keynes_Objeto.png}{Comparativa sobre el objeto}{Apuntes Víctor Gutiérrez. Repositorio TCEE}

\begin{specialblock}{Estructura de la Teoría General del Empleo, el Interés y el Dinero}
    Esta obra cuenta con veinticuatro capítulos repartidos en seis libros. 
    
    El primer libro (Introducción) está compuesto por tres capítulos en los que critica los postulados de la economía clásica y desarrolla el Principio de la Demanda Efectiva. 
    
    El segundo libro \textit{Definiciones e Ideas} versa sobre las unidades de medida utilizadas para medir las variables económicas\footnote{Precisa que las variables económicas las considerará medidas en unidades de salarios, es decir, en volumen de empleo.} y aclara algunas definiciones relevantes para el resto de la obra, como la definición de ingreso. Este libro consta de cuatro capítulos. 
    
    En el tercer libro (La propensión a consumir) Keynes explica los determinantes de la función de consumo a lo largo de los tres capítulos que lo componen. 
    
    El cuarto libro (El incentiv o para invertir) es el más largo, con ocho capítulos. En este libro expone minuciosamente los determinantes de la inversión, desarrollando conceptos clave como la eficiencia marginal del capital, las propiedades esenciales del dinero y la teoría general del tipo de interés. 
    
    En el quinto libro (Salarios nominales y precios), que consta de tres capítulos, argumenta que la flexibilidad de los salarios nominales no es un remedio adecuado para eliminar el desempleo involuntario. 
    
    El sexto y último libro (Breves consideraciones sugeridas por la Teoría General) cuenta con tres capítulos en los que desarrolla la teoría del ciclo económico y reabre debates pasados que se creían cerrados\footnote{Keynes considera que las innovaciones contenidas en la Teoría General permiten revisar las conclusiones obtenidas por los economistas clásicos y neoclásicos sobre algunos temas. como el mercantilismo o las leyes contra la usura.} en economía. También expresa opiniones en otros temas de su época, como la distribución de la renta o el papel del Estado en la economía.
\end{specialblock}

Citando a KEYNES: \textit{He llamado a mi teoría una teoría general, por ello me refiero a que estoy preocupado por el comportamiento del sistema económico en su conjunto}. En otras palabras, considera que era necesario estudiar el desempleo involuntario en un marco de equilibrio general (aunque él no uso esta expresión), ya que su origen debe ser buscado en otros sectores de la economía y no sólo en el mercado de trabajo.\footnote{Por ejemplo, en su apéndice al Capítulo 19, donde criticó a PIGOU, KEYNES escribió:\\
\textit{Mantengo que el salario real no está determinado principalmente por ajuste salarial sino por otras fuerzas del sistema, en particular la relación entre el esquema de la eficiencia marginal del capital y el tipo de interés}.} 

La decisión de KEYNES de adoptar una perspectiva de interdependencia no debe ser interpretada como una adhesión al enfoque de equilibrio general walrasiano.\footnote{En aquella época, las ideas de WALRAS no eran apreciadas en Cambridge y, KEYNES no pensaba que la teoría walrasiana pudiera ser de utilidad en su proyecto. CLOWER cita un extracto de una carta que envío KEYNES a GEORGESCU-RODAN en 1934: \textit{All the same, I shall hope to convince you some day that Walras’s theory and all the others along those lines are little better than nonsense}.\\
Como consecuencia, KEYNES adopta una aproximación más ensayística, influenciado por otros economistas de la escuela de Cambridge (como MARSHALL) que conceden una mayor importancia al razonamiento verbal y se aleja de miembros de la escuela de Lausanne (como WALRAS o PARETO) que otorgaron un papel protagonista al instrumental matemático.} KEYNES consideraba que el camino adecuado era adoptar un enfoque marshalliano y tratar de generalizarlo realizando las pertinentes modificaciones para defender sus tesis. 

Por otro lado, pese al surgimiento de una corriente de pensamiento en Cambridge que justificaba la existencia de desempleo involuntario en la existencia de competencia imperfecta,\footnote{JOAN ROBINSON publicó en 1933 su obra \textit{The Economics of Imperfect Competition}.} KEYNES decide usar un marco de competencia perfecta (posiblemente porque asociara la competencia imperfecta con colusión y sindicatos, y él quería centrar el tiro en algo más profundo).


\subsection{\textit{Baseline}: modelo neoclásico (brevemente)}

\begin{specialblock}{DISCLAIMER.- \textit{Modelo neoclásico}}
    El modelo neoclásico como tal no existe, sino que es un modelo que creamos aregando aportaciones de distintos autores pertenecientes a esta escuela. Además, como hemos mencionado el pensamiento económico de KEYNES está sujeto a distintas interpretaciones, por lo que lo que expondremos a continuación solo será una valoración y una formalización personal para lo cual usaremos gráficos y fórmulas que no necesariamente fueron utilizadas por KEYNES en sus obras. 
    
    KEYNES no llegó a formalizar todas sus ideas antes de morir en 1946. Además, su obra era difícil de interpretar en algunas cuestiones. Aunque sus principales proposiciones y conclusiones sí eran claras, no era fácil integrar algunas partes fundamentales de su obra con otras. El propio KEYNES, en el último capítulo de su Teoría General invitaba a otros economistas a que desarrollaran su teoría. Su invitación fue recogida por economistas de diversas tendencias como JOHN HICKS, ALVIS HANSEN, ROY HARROD, OSCAR LANGE o ABBA LERNER que se dedicaron al trabajo de interpretar, sistematizar y completar el modelo keynesiano. Esta corriente, conocida como síntesis neoclásica encontró una amplia aceptación que la llevó a convertirse en la nueva ortodoxia.    
\end{specialblock}

En el esquema de pensamiento neoclásico, los agentes de una economía cerrada participan en tres mercados: 
\begin{lnum}
    \item \textit{El mercado de fondos prestables} (mercado financiero). En el que las empresas demandan el ahorro de los hogares para realizar inversiones, a cambio de un tipo de interés real;\footnote{
    Entendemos por realizar inversiones a comprar bienes de capital, que son aquellos que sirven para producir otros bienes (maquinaria, edificios, etc.)\\
    El mercado de fondos prestables tiene un papel fundamental en el sistema economico, al determinar un tipo de interés real de equilibrio que resulta en un nivel de ahorro idéntico a la inversión. La igualdad entre ahorro e inversión garantiza que toda la producción encuentre salida en el mercado, es decir, que no exista una escasez generalizada de demanda que deje productos sin vender.\\
    \textbf{Analíticamente:} Si denominamos a $Y$ a la producción (renta). $DA$ a la demanda agregada; $C$ al consumo, $I$ a la inversión y $S$ al ahorro. La renta se destina al consumo y al ahorro: $Y = C + S$. La demanda agregada es el consumo más la inversión: $DA= C + I$. La condición necesaria para que toda producción encuentre salida en el mercado, es decir, para que $Y = DA$, es que todo el ahorro de los hogares sea destinado a invertir: $S = I$.\\
    El cumplimiento de esta condición es conocido bajo el nombre de la Ley de Say, que fue objeto de un intenso debate entre autores clásicos como RICARDO (favorable) y MALTHUS (crítico), que se saldó con la victoria del primero y no volvió a reabrirse con intensidad hasta la publicación de la Teoría General de Keynes.} y,
    \item \textit{El mercado de trabajo}. En el que las empresas y los hogares intercambian horas de trabajo a cambio de un salario real;\footnote{La vigencia de la Ley de Say garantiza a las empresas la venta de todo producto ofrecido al precio de equilibrio (mercado de bienes/dinero). En consecuencia:
    \begin{lnum}
        \item La demanda de trabajo vendrá determinada por su productividad marginal, estando dispuestas a contratar trabajadores hasta el punto en el que la productividad marginal del trabajo sea idéntica al salario real. [Nótese que este resultado se deriva directamente de la Ley de Say: si se diese el caso de que algunos productos no encontrasen salida en el mercado, las empresas despedirían a algunos trabajadores, incluso aunque su productividad marginal fuese superior a su salario real; ya que no tiene sentido contratar trabajadores para producir unos bienes que no se van a vender, ni siquiera en el extremo de que trabajasen gratuitamente.]; y, por otro lado,
        \item Los trabajadores ofrecerán tiempo de trabajo hasta que el salario real se iguale a la desutilidad marginal del trabajo (más bien, que el salario real se iguale a la relación marginal de sustitución entre el consumo y el ocio, que solo coincide con la desutilidad marginal del trabajo si la función de utilidad es cuasilineal) [Se ha optado por utilizar la expresión desutilidad marginal del trabajo para facilitar la exposición, al ser la expresión que utiliza Keynes en la Teoría General].
    \end{lnum}
    La demanda y la oferta de trabajo determinan, conjuntamente, un salario real que equilibra el volumen de empleo ofrecido y demandado, donde la productividad marginal del trabajo coincide con la desutilidad marginal del trabajo. Los volúmenes de trabajo y capital (que consideraremos constate, por ser de mayor interés para KEYNES el estudio de los ciclos y no del crecimiento) generan un nivel de producción agregada al que denominamos oferta agregada o producción (que las empresas venden en el mercado de bienes/dinero). Como sabemos, empleo y producciónse relacionan mediante la Función de Producción.}
    \item \textit{El mercado de bienes/dinero}. Donde todos los agentes económicos intercambian bienes a cambio de otros bienes, utilizando el dinero como medio de pago.\footnote{En este mercado, las empresas venden sus productos a otras empresas (bienes de capital) y/o a los hogares (bienes de consumo) a cambio de dinero. Las empresas utilizan ese dinero para pagar rentas monetarias (salarios e intereses) a los hogares y para adquirir bienes de capital de otras empresas.\\
    A su vez, los hogares utilizan el dinero recibido para comprar los bienes de consumo ofrecidos por las empresas, o por otros hogares. En este mercado se intercambia, por tanto, un volumen fijo de producción u oferta agregada a cambio de un stock fijo de dinero u oferta monetaria. Esta oferta monetaria o stock de dinero es la cantidad de dinero en circulación, cuya cuantía fija la determina una autoridad monetaria central (un banco central), siguiendo las normas de algún sistema monetario. Por ejemplo, bajo el patrón oro, el dinero que la autoridad monetaria pone a circular en la economía debe estar respaldado por sus reservas de oro (en alguna ratio, porcentaje).\\
    Nótese, por tanto, que \textbf{el circuito monetario comienza y termina en este mercado}. }
\end{lnum}   

\textbf{Importante.-} Será en este último mercado donde el sector real de la economía (mercado de trabajo, mercado de fondos prestables y oferta agregada) se encuentre con el sector monetario (oferta monetaria y demanda de dinero), dando lugar a un nivel general de precios de equilibrio. Es decir, en el mercado de bienes/dinero, la demanda de bienes es al mismo tiempo oferta de dinero y viceversa. Por tanto, el nivel de precios que resulta de este mercado equilibra la demanda de bienes con la oferta agregada y la demanda de dinero con la oferta monetaria (simultaneamente).

\subsubsection{Exceso de demanda de dinero en el Modelo Neoclásico: Mecanismo de Ajuste}
Si hubiera un exceso de oferta de bienes (esto es, un exceso de demanda de dinero), los agentes económicos tratarían de vender su exceso de bienes para paliar su escasez de demanda de dinero y la competencia entre los agentes económicos por vender sus bienes bajaría el precio de los mismos hasta que el exceso de oferta de bienes (exceso de demanda de dinero) desapareciese. 

Podemos entender mejor estos efectos partiendo de la Ecuación de Fisher (Teoría Cuantitativa del Dinero, 1911), que proporciona la siguiente ecuación: 

$$
\begin{aligned}
\underbrace{MV}_{\substack{\text{\scriptsize Valor monetario de}\\\text{\scriptsize las compras de bienese}}} =\underbrace{PY}_{\substack{\text{\scriptsize Valor monetario de}\\\text{\scriptsize las ventas de bienese}}}\quad 
&\text{donde,}
\begin{cases}
    M\equiv\quad \text{Oferta monetaria}\\[0.5em]
    V\equiv\quad \text{Velocidad de circulación}\quad \substack{
    \text{\scriptsize (Recordar, se creía relativamente constante)}\\
    \text{\scriptsize (Cantidad de veces que la misma unidad monetaria}\\
    \text{\scriptsize se utiliza para realizar un intercambio)}}\\[0.5em]
    Y\equiv\quad \text{La producción}\\[0.5em]
    P\equiv\quad \text{El nivel general de precios}\\[0.5em]
\end{cases}\\[2em]
\Longrightarrow \quad &M^d=\frac{PY}{V}\qquad \text{[2]}
\end{aligned}$$ 

Despejado $Y$, entonces obtendríamos qué cantidad de bienes podrán adquirir el conjunto de los agentes (teniendo en cuenta que el valor monetario de sus compras viene determinado por $MV$, y para un nivel de precios $P$). Es decir, obtenemos la demanda agregada de bienes: 

$$Y^d=\frac{MV}{P}\qquad \text{[3]}$$

Por lo tanto, en equilibrio $M^d = M$ y $Y^d = Y$. Puede comprobarse fácilmente la equivalencia entre un exceso de demanda de dinero y un exceso de oferta de bienes:

$$M^d=\frac{PY}{V}>M\qquad \Longrightarrow\qquad Y>\frac{MV}{P}=Y^d$$

Y viceversa:

$$M^d=\frac{PY}{V}<M\qquad \Longrightarrow\qquad Y<\frac{MV}{P}=Y^d$$

Siendo la producción $Y$ y la oferta monetaria $M$ dos variables determinadas fuera del mercado de bienes/dinero, y suponiendo $V$ una variable muy estable y determinada por otros factores independientes (tecnología de pagos, etc.), se deduce que todo desequilibrio en el mercado de bienes/dinero se corrige únicamente a través de variaciones en el nivel general de precios $P$.\footnote{Nótese que las variaciones del nivel general de precios $P$ vienen compañadas necesariamente de variaciones de los salarios nominales en la misma proporción. En efecto, lo que se negocian en el mercado son salarios reales. Mientras no cambie el equilibrio del mercado de trabajo, es decir, siempre que no cambie la productividad marginal del trabajo o la desutilidad marginal del trabajo, los salarios reales permanecerán inalterados, y toda variación de precios vendrá acompañada de una variación de salarios nominales en la misma proporción.} Por tanto, la flexibilidad de $P$ juega un papel fundamental en el modelo neoclásico, al ser condición necesaria para el cumplimiento de la Ley Say. 

Si ante un desequilibrio en el mercado de bienes/dinero, el nivel general de precios no se ajusta, podría perfectamente haber productos sin vender, es decir, una escasez generalizada de demanda. Esto implica que, si bien el tipo de interés real es la variable clave que equilibra el ahorro y la inversión, esto solo será posible si los precios son perfectamente flexibles. Por tanto, como se presuponen perfectamente flexibles, el tipo de interés real garantiza el cumplimiento de la Ley de Say, esto es, que bajo ninguna circunstancia puede acontecer una escasez generalizada de demanda de bienes.\footnote{En estas circunstancias, las empresas producirán bienes y contratarán trabajadores hasta el punto en el que la productividad marginal del trabajo sea iguale a la desutilidad marginal del trabajo. Como resultado, la economía siempre se encontrará en situación de pleno empleo.} 

\subsubsection{Atesoramiento en el Modelo Neoclásico}
Llegados a este punto, cabe preguntarse qué ocurriría en el modelo neoclásico si se produjese un deseo repentino por parte de los agentes económicos de atesorar dinero, es decir, de mantener sus ahorros en activos líquidos (efectivo, depósitos a la vista, etc.), en lugar de acudir al mercado de fondos prestables para ceder estos ahorros a cambio de un tipo de interés. En este caso, podría parecer inevitable que tuviese lugar un exceso de ahorro sobre la inversión\footnote{Ejemplo (planteamiento keynesiano): Imaginemos que un hogar ingresa $Y$ unidades monetarias que destina al consumo ($C$) y al ahorro ($S$).\\En el modelo neoclásico, este hogar acudiría al mercado de fondos prestables a prestar un volumen de ahorro $S$ a las empresas. El tipo de interés real de equilibrio garantizaría que este nivel de ahorro $S$ fuese similar a la inversión (1) de las empresas. Imaginemos ahora que este hogar decide no acudir al mercado de fondos prestables a prestar sus ahorros, sino guardar estos ahorros en dinero (por ejemplo, en efectivo). Parece inevitable, en este caso, que se produzca una reducción de la oferta de fondos prestables. Ello debería dar lugar a un incremento del tipo de interés real, y en consecuencia, a una reducción de la inversión.\\
Nótese que esta reducción de la inversión no se corresponde con una reducción del ahorro, sino con el atesoramiento del mismo. El atesoramiento habrá provocado una reducción de la inversión (a través de un aumento del tipo de interés real), sin que al mismo tiempo se haya reducido el ahorro. El resultado aparente de esta situación es la aparición de un exceso de ahorro sobre la inversión: esto es, la aparición de una escasez generalizada de demanda agregada, provocando la ruptura de la Ley de Say.}, lo que se \textit{podría} traducir en un aumento del tipo de interés real (es decir, las variables nominales, en este caso el atesoramiento, impactan sobre las variables reales, en este caso el tipo de interés y por tanto, la inversión). Sin embargo, lo que ocurre según el modelo neoclásico es distinto:
\begin{lnum}
    \item El atesoramiento no es más que un incremento de la demanda de dinero, que provoca \textit{ceteris paribus} un exceso de demanda de dinero;
    \item Para paliar este exceso de demanda de dinero, los agentes económicos necesitan ofrecer más bienes en el mercado, pero dado que ello no es posible (la oferta agregada de bienes es fija), lo que harán será demandar menos bienes, dando lugar \textit{ceteris paribus} a un exceso de oferta de bienes; por tanto,
    \item El exceso de demanda de dinero y el exceso de oferta de bienes se corrigen simultáneamente a través de una reducción del nivel general de precios; y,
    \item Como la reducción del nivel general de precios aumenta el poder adquisitivo de los fondos prestables (por el abaratamiento en términos salariales y de los bienes de capital);
    \item Aunque una parte del ahorro se haya atesorado, la parte que sí se ofrece en el mercado es capaz de financiar el mismo volumen de inversión real que antes del atesoramiento\footnote{Ejemplo: Imaginamos que los agentes ahorran 50 unidades monetarias, que ceden a las empresas en el mercado de fondos prestables a cambio de un tipo de interés real. Las empresas utilizan estos ahorros para realizar inversiones por un valor de 50 unidades monetarias. Si ahora los agentes económicos decidiesen atesorar 10 unidades monetarias, tan solo acudirían al mercado de fondos prestables con las 40 unidades monetarias restantes. Sin embargo, este atesoramiento de dinero provocaría una caída de los precios de los bienes que permitiría a las empresas realizar con 40 unidades monetarias las mismas inversiones que antes realizaban con 50 unidades monetarias. Como resultado, el atesoramiento no reduce la inversión. Por supuesto, esta caída de precios no aumentaría la riqueza de la sociedad en su conjunto: si tras esa bajada de precios los agentes decidiesen aprovechar para desatesorar las 10 unidades monetarias (por ejemplo, comprando bienes de inversión o de consumo), lo único que ocurriría es que subirían de nuevo los precios y volveríamos exactamente a la situación inicial. En definitiva, el atesoramiento, de existir, solo afecta a los precios y nunca a las variables reales. En realidad, los efectos del atesoramiento sobre el sistema económico son equivalentes a una reducción de la cantidad de dinero en circulación. El [\textbf{Anexo I}] recoge, paso a paso, una explicación gráfica de esta situación.};
    \item En consecuencia, ni el tipo de interés real de equilibrio ni la inversión se verán afectados, garantizando, de nuevo, el cumplimiento de la Ley de Say. 
\end{lnum}

\paragraph{Representación gráfica}
La representación gráfica del modelo neoclásico se puede hacer a través de los siguientes seis gráficos. 

\imagenfit{Atesoramiento_Neoclasicos.png}{Representación gráfica del Equilibrio Neoclásico y los efectos del \textit{atesoramiento}}{Apuntes ICEX-CECO. Tema 3.A.4}

Los gráficos representan lo siguiente:
\begin{lnum}
    \item El primero de ellos representa el mercado de trabajo, que determina un nivel de empleo de equilibrio.
    \item El gráfico inferior transforma ese nivel de empleo en un volumen de producción, a través de la función de producción.
    \item El siguiente gráfico es instrumental y solo sirve para llevar la variable Y del ej e de ordenadas al eje de abscisas;
    \item En el gráfico superior, se muestra el mercado de bienes. El nivel de producción u oferta agregada se determinó en el mercado de trabajo, por lo que aparece en el mercado de bienes como una variable dada, e independiente del nivel general de precios. La demanda agregada de bienes proviene de la ecuación cuantitativa del dinero, y tiene pendiente negativa;
    \item  A la derecha, tenemos el mercado de dinero. La demanda de dinero también proviene de la ecuación cuantitativa del dinero, y tiene pendiente positiva. El mercado de bienes y el mercado de dinero son, en realidad, un mismo mercado que proporciona un nivel general de precios que equilibra simultáneamente la demanda agregada con la producción y la demanda de dinero con la oferta monetaria; por último,
    \item El último gráfico se corresponde con el mercado de fondos prestables, que equilibra el ahorro y la inversión a través del tipo de interés real.
\end{lnum}

Partiendo de un equilibrio inicial $M_0^d$:
\begin{lnum}
    \item Un deseo repentino de atesoramiento por parte de los agentes económicos daría lugar a una reducción de la velocidad de circulación del dinero, pues una parte del dinero disponible dejaría de utilizarse repentinamente para hacer transacciones; 
    \item Esto reduce la pendiente de la demanda de dinero ($M_1^d$) y desplaza hacia la izquierda la demanda de bienes ($Y_1^d$);\footnote{En realidad, más que un desplazamiento se produciría un cambio en la curvatura de la función. En la demanda de bienes, el precio se aproxima a infinito a medida que la producción se aproxima a cero y viceversa.\\
    Las variaciones de la velocidad de circulación del dinero no modifican esta propiedad de la función, pero sí alejan o acercan los valores de la función respecto del origen a través de la curvatura de la función. En este caso, la reducción de la velocidad de circulación del dinero aproximaría todos los pares $(P,Y)$ hacia el origen, pero dicho desplazamiento hacia la izquierda no sería paralelo.
    } 
    \item Como resultado, para el nivel de precios inicial ($P_0$) habrá un exceso de demanda de dinero y una escasez de demanda de bienes;
    \item Los oferentes de bienes tendrán que competir bajando los precios. A medida que esta bajada de precios se produce, los agentes económicos necesitarán menos dinero para realizar transacciones;
    \item La reducción de los precios reducirá en la misma proporción la $PMg_L$ y la desutilidad marginal del ocio;
    \item Lo que conducirá a una reducción de los salarios nominales en la misma proporción, manteniendo inalterado el salario real; 
    \item En el nuevo equilibrio ($P_1$), desaparecen simultáneamente los excesos de oferta de bienes y de demanda de dinero y las variables reales de la economía (salario real, empleo, producción, ahorro, inversión y tipo de interés real).
\end{lnum}

El ejemplo anterior nos muestra como en el modelo neoclásico existe una separación total entre los aspectos monetarios del sistema económico (la demanda de dinero y la oferta monetaria) y la economía real (el salario real, el tipo de interés real, la producción y el empleo). Nada de lo que ocurra en la parte monetaria del sistema afectará a las variables reales (al largo plazo), por ello, el dinero es neutral (a largo plazo). La neutralidad del dinero se podría observar también estudiando los efectos de un incremento en la cantidad de dinero (pero haciendo una distinción fundamental entre Tasa de interés natural y Tasa de interés nominal, citando a WICKSELL). En este caso, la oferta monetaria se desplazaría hacia la derecha, al igual que la demanda agregada. Fruto de estos desplazamientos se produciría, para los precios iniciales, un exceso de oferta de dinero y un exceso de demanda de bienes. Ambos excesos se corregirían simultáneamente a través de un incremento del nivel general de precios, sin que ninguna variable real sufriese alteración alguna.

\subsection{Formalización del modelo keynesiano: sistema económico y política económica}

\subsubsection{Teoría General del Tipo de Interés}
La diferencia más inmediata entre el pensamiento de Keynes y el modelo neoclásico anterior tiene que ver con las propiedades del dinero. \footnote{Como hemos visto, en el modelo neoclásico, el dinero es simplemente un medio de pago cuya función es facilitar los intercambios de mercancías, evitando los altos costes de transacción del sistema de trueque. Los agentes económicos no tienen ningún incentivo a atesorar dinero, puesto que éste no proporciona ningún rendimiento ni utilidad intrínseca. En cualquier caso, ya hemos visto que, si por algún motivo hubiera atesoramiento, lo único que ocurriría en el sistema económico es una reducción generalizada de los precios. }

KEYNES considera que el dinero cuenta con una serie de características especiales que justifican que los agentes económicos prefieran atesorarlo en lugar de emplearlo en otros usos. El dinero, a diferencia de cualquier otra mercancía, permite a su poseedor conservar riqueza libre de toda incertidumbre.\footnote{En línea con las ideas expresadas en su Tratado sobre la Probabilidad. Keynes entiende por incertidumbre aquella situación en la que se desconocen las probabilidades de que ocurran los difer entes escenarios que pueden tener lugar. ni se puede ordenar a los posibles escenarios en función de su probabilidad.} 

En un escenario de incertidumbre, los títulos financieros que compran los hogares en el mercado de fondos prestables (acciones, bonos, etc) están sujetos a diversos riesgos. Incluso aquellos títulos menos arriesgados (por ejemplo, deuda pública a corto plazo) están expuestos al riesgo de tipo de interés, es decir, a la posibilidad de que su precio fluctúe a medida que se producen sorpresas en el tipo de interés.\footnote{Ejemplo: Un hogar podría comprar en el año $t$ un bono que reporta $X$ unidades monetarias en el año $t+1$. El precio de dicho bono en $t$ vendrá determinado por $\frac{X}{1+i}$, donde $i$ es el tipo de interés expresado en tantos por uno. Puede verse claramente que existe una relación inversa entre el precio del bono y el tipo de interés anual. Por tanto, el poseedor de este bono estará expuesto, durante el tiempo que lo mantenga en su cartera (en este caso, un año si no lo vende antes de su vencimiento), a fluctuaciones en el precio del bono, provocadas por las fluctuaciones del tipo de interés.} 

KEYNES indica claramente en el Capítulo 13 de la Teoría General que la incertidumbre que rodea a la evolución futura de los tipos de interés es condición necesaria para que los agentes económicos decidan atesorar dinero: 

\begin{adjustwidth}{2cm}{0pt}
\raggedright 
\textit{Pero dado que la tasa de interés nunca es negativa, ¿por qué preferiría alguien guardar su riqueza en una forma que rinde poco o ningún interés a conservarla en otra que sí lo da (suponiendo, claro está, por el momento, que el riesgo de pérdida es igual para un saldo bancario que para un bono)? (...) Hay una condición necesaria sin la cual no podría haber preferencia por la liquidez como medio de conservar riqueza. Esta condición necesaria es la incertidumbre respecto al futuro de la tasa de interés (...) porque si pudiera preverse con certeza las que dominen en todo tiempo en el futuro, todas las tasas venideras podrían inferirse de las presentes para las deudas de diversos plazos, las que se ajustarían al conocimiento de las tasas futuras}\\
\raggedleft{\authorfont{KEYNES}, “\textit{Teoría General del Empleo, el Interés y el Dinero}”, 1936}
\end{adjustwidth}

De no existir incertidumbre sobre los tipos de interés en el futuro, el riesgo desaparecería y no habría razón para atesorar dinero (siempre que exista en el mercado algún bono o título libre de riesgo de impago, como la deuda pública a corto plazo). Sin embargo, dado que los tipos de interés en el futuro son inciertos, los agentes económicos tendrán incentivos para mantener una parte de su riqueza en activos líquidos. Vuelve a remarcar esta idea en el Capítulo 15 de la Teoría General: 

\begin{adjustwidth}{2cm}{0pt}
\raggedright 
\textit{En una sociedad estática o en que, por cualquier otra razón, nadie sienta la menor incertidumbre acerca de las futuras tasas de interés (...), la propensión a atesorar, será siempre cero en equilibrio}\\
\raggedleft{\authorfont{KEYNES}, “\textit{Teoría General del Empleo, el Interés y el Dinero}”, 1936}
\end{adjustwidth}

No obstante, la existencia de atesoramiento como medio de protección frente a la incertidumbre no permite, por sí misma, rechazar las conclusiones del modelo neoclásico. Como hemos visto, el modelo neoclásico es capaz de explicar que, de existir atesoramiento, el único efecto que tendría lugar en el sistema económico es una reducción del nivel general de precios, sin alterar el equilibrio real del sistema económico. La verdadera innovación de Keynes consiste precisamente en rechazar esta proposición. 

Para ello, presupone que en una economía monetaria con incertidumbre los agentes económicos toman sus decisiones de consumo y demanda de dinero en dos etapas sucesivas:
\begin{lnum}
    \item En primer lugar, deciden qué porciones de su renta destinan al consumo y al ahorro; y,
    \item Posteriormente, deciden qué porciones de su riqueza resultante (tras añadir el ahorro) se conservan en activos no líquidos (bonos, acciones, etc.) y en activos líquidos (efectivo, depósitos a la vista, etc.).
\end{lnum} 

Al hacer este supuesto, los agentes ya no solo demandan dinero para hacer transacciones en el mercado de bienes, sino también para conservar una parte de su riqueza en activos líquidos. La demanda de dinero ya no será una proporción $\frac{1}{V}$ de la renta nominal, sino una proporción de la riqueza de los agentes económicos. Es decir, la verdadera innovación de KEYNES consiste en la sustitución del mercado de bienes/dinero por el de bonos/dinero. 

En el esquema de KEYNES también habrá tres mercados, pero su configuración será diferente a la del modelo neoclásico:
\begin{lnum}
    \item En primer lugar, existirá un mercado de bienes, en el que las empresas venden sus productos a los hogares, que devuelven a cambio las rentas (salarios e intereses) que previamente recibieron de las empresas.
    \item Habrá un mercado de trabajo, en el que se intercambian horas de trabajo a cambio de un salario;\footnote{Como veremos después, el mercado de trabajo del modelo keynesiano no determina ni el volumen de empleo ni el salario real. Las empresas y los trabajadores negocian un salario nominal que solo a fecta de manera directa al nivel general de precios de la economía. Lo que ocurra en este mercado solo afecta a las variables reales (producción, empleo, salario real, etc.) de manera indirecta e incierta.} y,
    \item Habrá un mercado de bonos/dinero, en el que los agentes económicos intercambian sus activos no líquidos (bonos)\footnote{Por simplificar, desarrolla su teoría asumiendo que todos los mercados financieros se resumen en un único mercado en el que se intercambian bonos a cambio de dinero.} a cambio de activos líquidos (dinero). 
\end{lnum} 

De esta forma, toda oferta de bonos será al mismo tiempo demanda de dinero y viceversa. El hecho de que algunos agentes económicos estén dispuestos a vender bonos para obtener dinero responde, precisamente, a su deseo de atesorar dinero. En una situación de incertidumbre, los agentes pueden preferir conservar parte de su riqueza en dinero, incluso aunque el rendimiento del dinero sea nulo. 

El mercado resultante de bonos/dinero genera un precio de equilibrio que iguala simultáneamente la demanda de bonos con su oferta, y la demanda de dinero con la oferta monetaria. El precio de los bonos guarda una relación única e inversamente proporcional con el tipo de interés nominal.  De este modo, si el precio de los bonos (el tipo de interés nominal) fuera superior al de equilibrio, habría un exceso de oferta de bonos que presionaría a la baja el precio de los bonos hasta restaurar el equilibrio, y viceversa. 

La sustitución del mercado de bienes/dinero por el mercado de bonos/dinero modifica por completo las consecuencias del atesoramiento, invalida la teoría neoclásica del tipo de interés e imposibilita el cumplimiento de la Ley de Say. En este nuevo esquema, las variaciones en la demanda de dinero no se corresponden con variaciones en la demanda de bienes (como en el modelo neoclásico), sino con variaciones en la oferta de bonos.\footnote{En el modelo neoclásico, cuando un agente económico necesita dinero, vende algunos bienes en el mercado para conseguir ese dinero. Según Keynes, cuando un agente económico necesita dinero. vende algunos bonos para conseguir ese dinero.} En consecuencia, las decisiones de atesoramiento no provocan cambios en el nivel general de precios de los bienes, sino en el precio de los bonos. Los cambios en el precio de los bonos implican necesariamente cambios en el tipo de interés nominal, que a su vez provocan cambios en el tipo de interés real. Las decisiones de atesoramiento dan lugar, por tanto, a variaciones en el tipo de interés real, que generan efectos reales (reducción de la inversión y aumento del ahorro) en el sistema económico, y no meramente nominales (cambios en el nivel general de precios) como ocurría en el modelo neoclásico. El nuevo planteamiento de Keynes rechaza, por tanto, lo que él llama la Teoría clásica del tipo de interés, y da lugar a lo que nombra como la Teoría General del Tipo de interés. 

A diferencia del enfoque neoclásico, Keynes no cree que el tipo de interés real se determine en el mercado de fondos prestables,\footnote{Este mercado no existe como tal en el enfoque keynesiano.} ni que sea el premio por ahorrar (o dicho de otra forma, el coste de oportunidad del consumo presente), pues considera que el tipo de interés real se determina como la diferencia entre un tipo de interés nominal determinado en el mercado de bonos/dinero y una inflación esperada exógena. El tipo de interés nominal, por su parte, será el premio por renunciar a la conservación de la riqueza en activos líquidos y vendrá determinado por la demanda de dinero y la oferta monetaria (el coste de oportunidad de la liquidez). 

Bajo este nuevo enfoque, las fuerzas que determinan el tipo de interés real ya no guardan relación directa ni con el ahorro ni con la inversión, por lo que no hay ningún motivo para confiar en un equilibrio automático entre estas dos variables:

\begin{adjustwidth}{2cm}{0pt}
\raggedright 
\textit{Debiera parecer evidente que la tasa de interés no puede ser la recompensa al ahorro o a la espera como tales; porque si un hombre atesora sus ahorros en efectivo no gana interés, aunque ahorre lo mismo que antes. Por el contrario, la mera definición de tasa de interés nos dice, en muchas palabras, que la tasa de interés es la recompensa por privarse de liquidez durante un periodo determinado; porque dicha tasa no es, en sí misma, más que la inversa de la proporción que hay entre una suma de dinero y lo que se puede obtener por desprenderse del control del dinero a cambio de una deuda durante un periodo determinado de tiempo.\\
La tasa de interés no es el "precio" que pone en equilibrio la demanda de recursos para invertir con la buena disposición para abstenerse del consumo presente. Es el "precio " que equilibra el deseo de conservar la riqueza en forma de efectivo, con la cantidad disponible de este último}\\
\raggedleft{\authorfont{KEYNES}, “\textit{Teoría General del Empleo, el Interés y el Dinero}”, 1936}
\end{adjustwidth}

La teoría general del tipo de interés defendida por Keynes implica necesariamente la ruptura con la Ley de Say. Suponiendo que no existe inflación, el tipo de interés nominal determinado en el mercado de dinero (izquierda) será similar al tipo de interés real, que determina el ahorro y la inversión (derecha).\footnote{El tipo de interés del mercado de dinero (izquierda) es nominal mientras que el tipo de interés del mercado de fondos prestables (derecha) es real. Por sencillez, asumimos que no existe inflación, de forma que ambos coincidan. Si añadiesemos la inflación, simplemente tendríamos un tipo de interés más elevado en el gráfico de la izquierda que en el de la derecha. siendo la inflación la diferencia entre ambos. Las conclusiones, no obstante, serían idénticas: el tipo de interés real no iguala necesariamente el ahorro y la inversion.} 

\imagenfit{RupturaSAY_TipoInteres.png}{La Teoría General del Tipo de Interés}{Apuntes ICEX-CECO. Tema 3.A.4}

Como veremos más adelante, el tipo de interés real resultante no tiene por qué igualar el ahorro y la inversión. 

Además, completa la Teoría General del Tipo de interés con un análisis minucioso de la demanda de dinero o preferencia por la liquidez. La curva de demanda de dinero o la curva de preferencia por la liquidez, nos indica la cantidad de riqueza que los hogares desean mantener en activos líquidos (efectivo, depósitos a la vista, etc.) para cada posible tipo de interés. La pendiente negativa de la curva de demanda de dinero indica que, a medida que aumenta el tipo de interés (esto es, el coste de oportunidad de mantener riqueza en dinero) se reduce la demanda de dinero, pues los agentes económicos preferirán guardar más riqueza en activos no líquidos (bonos). Para ello, distingue tres clases de preferencia por la liquidez:
\begin{lnum}
    \item \textit{Motivo transacción}. En primer lugar, los agentes económicos mantienen una parte de su riqueza en dinero para poder realizar operaciones corrientes. El principal determinante de la demanda de dinero por motivo de transacción será la renta nominal ($PY$), por lo que toda variación de producción agregada o del nivel general de precios vendrá acompañada de una variación, para cada tipo de interés, de la demanda de dinero en el mismo sentido;
    \item \textit{Motivo especulación}. Los agentes demandan dinero por la creencia de que el precio actual de los bonos es demasiado alto con respecto al futuro, o dicho de otra forma el tipo de interés es demasiado bajo con respecto al futuro. La demanda de dinero por motivo de especulación puede servir para aprovechar las variaciones en el precio de los bonos (es decir, del tipo de interés) para obtener ganancias de capital, este sería el caso de los traders profesionales\footnote{Ejemplo: un especulador profesional cuenta con unos recursos que invierte en diferentes activos. Por sencillez, imaginamos que solo opera con bonos de renta fija. El especulador tratará de comprar bonos al menor precio posible para venderlos al mayor precio posible. Las variaciones del precio de los bonos vendrán determinadas por el tipo de interés, por lo que:
    \begin{la}
        \item Cuando el tipo de interés sea demasiado alto (según las perspectivas del especulador), el precio de los bonos será demasiado bajo y decidirá comprar bonos (reducir su demanda de dinero por motivo de especulación);
        \item Cuando el tipo de interés sea demasiado bajo, el precio de los bonos será demasiado alto y por tanto buen momento para vender los bonos (demandar dinero por motivo de especulacion).
    \end{la}
    Este tipo de operaciones que en la vida real no se realizan solo con bonos sino con todo tipo de activos, es lo que entendemos por '\textit{demanda de dinero por motivo de especulación}.} o para evitar el riesgo del tipo de interés a la hora de conservar riqueza, es decir, inversores que buscan reducir su exposición a la incertidumbre del mercado.\footnote{Ejemplo: Un particular podría destinar la totalidad de sus ahorros a la compra de bonos con vencimiento $t+2$, con la intención de obtener un rendimiento de sus ahorros. Sin embargo, el ahorrador sabe que es posible que necesite vender los bonos en $t+ 1$ para cubrir cualquier necesidad imprevista. Si el tipo de interés subiese en $t+1$, el valor de los bonos habrá caído, provocando una pérdida. El ahorrador entonces puede mantener todos sus ahorros en dinero para evitarlo, es decir, demandar dinero por motivo de especulación si considera que la subida del tipo de interés en $t+1$ es probable.} La demanda de dinero por motivo de especulación dependerá de la diferencia entre el tipo de interés presente y las expectativas sobre el tipo de interés futuro. La variedad de las opiniones sobre dichas expectativas determinará la sensibilidad de este componente de la preferencia por la liquidez con respecto al tipo de interés.\footnote{En efecto, si todos los inversores tienen las mismas expectativas, una variación del tipo de interés hará que todos ellos reaccionen de manera similar, demandando más o menos dinero en la misma dirección. El resultado será una demanda de dinero por motivo de especulación muy elástica. Por el contrario, si hay diferencias de opiniones, la demanda de dinero por motivo de especulación será más inelástica.}
    \item \textit{Motivo precaución}. Permite a los agentes económicos mantener un exceso de dinero para satisfacer cualquier imprevisto y su importancia dependerá de la facilidad de liquidación de los activos alternativos.\footnote{La demanda de dinero por motivo precaución dependerá de lo líquidos que sean los activos alternativos al dinero:
    \begin{la}
        \item Si los mercados de activos (bonos, acciones, etc.) son poco líquidos, la demanda de dinero por motivo de precaución aumentará, pues los agentes económicos buscarán evitar los costes de transformar activos poco líquidos en dinero, y para ello demandarán más dinero por motivo de precaución.
        \item Si por el contrario, los mercados de activos son muy líquidos, los agentes no necesitarán demandar demasiado dinero por motivo de precaución, pues ante cualquier imprevisto, podrán transformar sus activos en dinero fácilmente.
    \end{la}}
\end{lnum} 

Como hemos visto, la nueva concepción keynesiana del dinero y de la preferencia por la liquidez sustituye la Teoría clásica del tipo de interés por una nueva Teoría General del Tipo de interés, que conecta los aspectos monetarios (demanda de dinero y oferta monetaria) del sistema económico con la economía real (tipo de interés real, ahorro e inversión).  En la concepción keynesiana del sistema económico no serán válidas:
\begin{lnum}
    \item Ni la dicotomía clásica;\footnote{La dicotomía clásica establece que la economía monetaria y la economía real no interfieren entre sí, de forma que las variables nominales son totalmente independientes de las variables reales.} 
    \item Ni la neutralidad del dinero;\footnote{La neutralidad del dinero es una proposición neoclásica que establece la cantidad de dinero en circulación tiene un efecto neutral sobre las variables reales del sistema económico. Igualmente. las variaciones en la demanda de dinero tan solo modifican el nivel general de precios.} 
    \item Ni la Ley de Say.
\end{lnum} 

\subsubsection{Teoría del Empleo: Principio de la Demanda Efectiva}
El abandono de la Ley de Say abre la posibilidad a un exceso o escasez de demanda generalizada. Por tanto, la existencia de trabajadores dispuestos a trabajar a un salario igual o inferior a la productividad marginal del trabajo ya no será suficiente para que dichos trabajadores estén empleados: si la producción adicional que generan estos trabajadores no se consiguiese vender, nunca tendría sentido su contratación. 

El volumen de empleo, por tanto, no vendrá determinado por la oferta y la demanda de trabajo (es decir, por la intersección de la desutilidad marginal del trabajo y su productividad marginal), sino por el \textit{Principio de la Demanda Efectiva}, al que Keynes dedica la totalidad del Capítulo 3 de la Teoría General, y lo formaliza a través de dos ecuaciones:
\begin{lnum}
    \item \textit{Función de Oferta Agregada}. Proporciona los ingresos monetarios mínimos que el conjunto de las empresas exige obtener para estar dispuestas a contratar un determinado volumen de empleo ($N$), tomando como dados: (i) el salario nominal ($W$); (ii) la función de próducción; y, (iii) un grado de competencia en el mercado de bienes. Keynes no especificó qué forma tendría esta función, pero sí aclaró que tendría pendiente positiva, de forma que el conjunto de las empresas solo estaría dispuesto a contratar más trabajadores si, \textit{ceteris paribus} los ingresos monetarios esperados por el conjunto de las empresas aumentan;\footnote{Este es uno de los casos en los que Keynes no terminó de explicar exactamente lo que quería decir, abriendo el paso a varias posibles interpretaciones. Algunos autores post-keynesianos trataron de estudiar las propiedades de esta función, representándola con forma convexa. Si el volumen de empleo es $N$, el salario nominal de un empleado es $W$ y el mark-up óptimo de las empresas es una proporción $\mu$ en tantos por uno sobre los costes variables, entonces la Función de Oferta Agregada sería: $P^O=(1+\mu)NW$. Donde $\mu$ a su vez dependerá de otras variables como el grado de monopolio, la función de producción, etc. Si representamos la Función de Oferta Agregada como una función convexa, asumimos que el mark-up óptimo depende positivamente del volumen de empleo $N$.}  y,
    \item \textit{Funcion de la Demanda Agregada}. Proporciona el valor esperado de los ingresos monetarios que el conjunto de las empresas espera obtener para cada posible nivel de empleo ($N$), tomando como dado el salario nominal ($W$). La Función de Demanda Agregada será la suma del gasto en consumo ($C$) e inversión ($I$) del conjunto de la economía. Dedicó gran parte de la Teoría General a explicar los determinantes de la Función de Demanda Agregada, pero tampoco especificó claramente qué forma tendría esta función aunque al igual que la Función de Oferta Agregada, tendría pendiente positiva. Esto implica que cuanto mayor sea el nivel de empleo de la economía, mayores serán las ventas que esperan obtener las empresas.\footnote{Según la interpretación de algunos autores post-keynesianos, esta función sería cóncava. El motivo de ello es que la propensión a consumir es decreciente con el nivel de empleo y renta (aunque lo justifican por diversos motivos). En efecto, deja claro en el Capítulo 10 de la Teoría General que:\\\textit{La propensión marginal a consumir no es constante para todos los niveles de ocupación, y es probable que ofrezca, por regla general, una tendencia a disminuir a medida que la ocupación crece}} 
\end{lnum}

\paragraph{Equilibrio: OA-DA}
La intersección de la Función de Demanda Agregada y la Función de Oferta Agregada proporciona el nivel de empleo de la economía, que hará coincidir las expectativas de ventas de las empresas con el nivel mínimo exigido por las mismas para contratar un determinado volumen de empleo.

\imagenfit{PpioDemandaEfectiva.png}{Principio de la Demanda Efectiva}{Apuntes ICEX-CECO. Tema 3.A.4}

Este punto de equilibrio será la Demanda Efectiva (DE).\footnote{KEYNES no incluyó ningún gráfico que ilustrase el Principio de la Demanda Efectiva. La representación de la $OA$ como una función convexa y la $DA$ como una función cóncava se encuentra en la obra de autores post-keynsianos como WEINTRAUB y DAVIDSON.} 

El rechazo de Keynes a la Ley de Say se ilustra muy bien con este gráfico, pues solo en el punto de DE, el ahorro y la inversión del conjunto de la economía se esperan iguales. Para cualquier punto a la izquierda del punto DE, se espera un exceso de demanda de bienes, porque el gasto agregado que realizarían los agentes superaría el nivel mínimo de ingresos que exigen las empresas. Por el contrario, a la derecha del punto DE, el gasto esperado del conjunto de la sociedad será inferior a los ingresos mínimos que exigen las empresas. 

Según Keynes, el cumplimiento de la Ley de Say requiere que las Funciones de Oferta y Demanda agregada sean idénticas, es decir, que estén perfectamente superpuestas, de forma que la Oferta y la Demanda Agregadas estén en equilibrio para todo posible nivel de empleo. 

Si este fuese el caso, las empresas maximizarían sus beneficios contratando trabajadores hasta que el salario real se igualase a la productividad marginal del trabajo. Si los trabajadores pudiesen ofrecer trabajo hasta que el salario real se iguale a la desutilidad marginal del trabajo, el pleno empleo estaría garantizado.

Por el contrario, si ambas funciones no están superpuestas, solo habrá un volumen de empleo de equilibrio, que será generalmente inferior al de pleno empleo. En efecto, puede observarse que, aunque de alguna forma se consiguiese reducir el salario real hasta llevarlo a su nivel de pleno empleo, las empresas no estarían dispuestas a contratar trabajadores hasta tal nivel, dado que sus ventas esperadas no serían suficientes para emplear tal volumen de trabajadores. Keynes lo expresa muy claramente en el Capítulo 3 de la Teoría General:

\begin{adjustwidth}{2cm}{0pt}
\raggedright 
\textit{La insuficiencia de la demanda efectiva frenará el proceso de la producción, aunque el producto marginal de la mano de obra exceda todavía en valor a la desutilidad marginal del trabajo }\\
\raggedleft{\authorfont{KEYNES}, “\textit{Teoría General del Empleo, el Interés y el Dinero}”, 1936}
\end{adjustwidth}

Bajo este enfoque se rechaza, por tanto, la teoría neoclásica del empleo, según la cual, la productividad marginal del trabajo y la desutilidad marginal del trabajo son las fuerzas que determinan el salario real y el empleo agregado de la economía. 

En el enfoque de Keynes, la Demanda Efectiva determina el empleo y el empleo determina el salario real. De hecho, el salario real puede ser superior al de equilibrio y provocar desempleo involuntario, ya que el carácter involuntario del desempleo tiene que ver con que el salario real esté determinado, en última instancia, por la Demanda Efectiva, y no por las negociaciones que lleven a cabo trabajadores y empresas en el mercado de trabajo. En efecto, no hay nada que los trabajadores puedan hacer para reducir su salario real, pues la única manera de lograr tal cosa es a través de un incremento de la Demanda Efectiva, algo que a priori está fuera de su alcance.\footnote{Los trabajadores podrían aceptar una rebaja generalizada en sus salarios nominales. Como veremos más adelante, la reducción de salarios nominales tendrá un efecto incierto sobre la Demanda Efectiva, y por tanto, sobre el nivel de empleo y el salario real (según KEYNES). Éste autor no niega que una reducción de salarios nominales pueda reducir el desempleo, pero lo considera poco probable.}

\paragraph{Propensión a consumir} 
Como hemos visto, la demanda agregada juega un papel fundamental en la Teoría del Empleo de Keynes y el estudio de sus determinantes ocupa gran parte de la Teoría General. 

La Función de Demanda Agregada se compone del gasto agregado en consumo y del gasto agregado en inversión (en economía cerrada). Dentro de ambos componentes se encuentran los gastos que realizan tanto el sector público como el sector privado. 

\textbf{Definición.-} La propensión a consumir será la relación funcional existente entre el consumo y sus determinantes objetivos y subjetivos, a la que KEYNES dedica prácticamente todo el Libro III. 

Dentro de los factores objetivos encontramos elementos como: (i) la renta\footnote{Recordemos que Keynes mide todas las variables en unidades de salario, es decir, en volumen de empleo. De esta forma, la función de consumo o propensión a consumir nos relaciona cuantos salarios se gastan para cada volumen de empleo: si hay $X$ personas trabajando, el consumo agregado será de $Y$ salarios.} total de la economía; (ii) cambios imprevistos en el valor del capital; (iii) el tipo de interés; (iv) la política fiscal; y, (v) las expectativas sobre el ingreso futuro. Dentro de los cuales el más relevante será la renta (según él).\footnote{Por el contrario, el tipo de interés será poco relevante:\\
\begin{adjustwidth}{2cm}{0pt}
\raggedright 
\textit{La conclusión más importante sugerida por la experiencia es, según creo, que la influencia a corto plazo de la tasa de interés sobre los gastos individuales hechos con un determinado ingreso es secundaria y carece relativamente de importancia, excepto. quizá, cuando se presenten cambios excepcionalmente grandes (...)}\\
\raggedleft{\authorfont{KEYNES}, “\textit{Teoría General del Empleo, el Interés y el Dinero}”, 1936}
\end{adjustwidth}} 

Dentro de los factores subjetivos se encuentran elementos diversos que van desde la previsión frente a la vejez hasta el orgullo o la pura avaricia. 

En general, Keynes considera que el consumo es una función bastante estable de la renta, es decir, que todos los demás factores apenas varían o tienen un escaso impacto; por eso, un concepto fundamental en la teoría del consumo propuesta por Keynes será la \textit{propensión marginal a consumir} (la derivada parcial del consumo con respecto a la renta, la sensibilidad del consumo a la renta). Keynes denominó \textit{Ley psicológica fundamental} a la proposición de que la propensión marginal a consumir estaba comprendida entre cero y uno. 

\imagenfit{Estimaciones_PMgCAG.png}{Propensiones Marginales estimadas para Estados Unidos y Alemania}{\textit{La función de consumo en la Economía Española}, BARBANCHO (aprox. 1960). Ministerio de Agricultura, Pesca y Alimentación}

\imagenfit{Estimaciones_PMgCAG_España.png}{Propensiones Marginales estimadas para España (1944-1959)}{\textit{La función de consumo en la Economía Española}, BARBANCHO (aprox. 1960). Ministerio de Agricultura, Pesca y Alimentación}

\paragraph{Determinantes de la inversión: Incentivo a invertir}
La Teoría General reserva su Libro IV a explicar los determinantes de la inversión. Afirma que las empresas realizaran inversiones hasta que el tipo de interés\footnote{En la Teoría General, Keynes compara la eficiencia marginal del capital y el tipo de interés en términos nominales. Se podría no obstante, comparar ambas variables en términos reales, dando los mismos resultados.} se iguale a lo que él denominó la eficiencia marginal del capital.

\textbf{Definición.-} (Eficiencia marginal del capital) La tasa de descuento que lograría igualar el valor presente de la serie de anualidades dada por los rendimientos del bien de capital, en todo el tiempo que dure, a su precio de oferta.\footnote{Esta definición es idéntica o muy similar a lo que hoy entendemos como Tasa Interna de Rentabilidad (TIR).La Tasa Interna de Rentabilidad (TIR) indica la rentabilidad de un proyecto. Matemáticamente, es el tipo de descuento que iguala el Valor Actualizado Neto (VAN) de una inversión a cero.} 

La curva de eficiencia marginal del capital relaciona el tipo de interés con la inversión total de la economía, y tendrá pendiente negativa (es decir, a mayor tipo de interés, menor inversión): el tipo de interés y la curva de eficiencia marginal del capital serán los determinantes de la inversión. 

Cabe destacar, en lo referido la curva de eficiencia marginal del capital, que ésta se basa en expectativas sobre los rendimientos futuros de las inversiones. Para él, estas expectativas están sujetas a la existencia de incertidumbre, entendiendo como tal la imposibilidad de asignar probabilidades a los distintos resultados posibles. De esta forma, los inversores operan en base a lo que denomina \textit{animal spirits}, es decir, del optimismo espontáneo que sintiesen por cualquier motivo. Ahora bien, debemos aclarar algo muy importante, según el autor:
\begin{la}
    \item Antes de que apareciese la bolsa de valores, las inversiones las realizaban empresarios comprometidos con una determinada actividad a largo plazo. Las inversiones no eran líquidas, por lo que los empresarios trataban de estimar con el máximo rigor posible los rendimientos futuros de las inversiones. La eficiencia marginal del capital estaba determinada, por tanto, por las expectativas de empresarios informados y comprometidos con una determinada actividad económica; más tarde,
    \item La aparición de la bolsa de valores modifica esta situación. Gracias a este nuevo mercado de activos, las inversiones de las empresas se convierten en activos más líquidos, que se pueden transformar rápidamente en dinero ante cualquier necesidad eventual de los inversores. En este nuevo escenario, el valor de las empresas y de las inversiones ya no vendrá determinado por las expectativas de los empresarios que ponen en marcha y gestionan las empresas, sino por los inversores que trafican en la bolsa de valores.  
\end{la}

Dicho de otra forma, tras la aparición de la bolsa de valores, las expectativas que determinan la curva de eficiencia marginal del capital serán las de los traders que participan en el mercado de valores. Por tanto:
\begin{la}
    \item Si los participantes de la bolsa son optimistas, la curva de eficiencia marginal del capital aumentará y las acciones serán más caras, generando incentivos a la inversión.\footnote{Keynes parece anticiparse a lo que más tarde se conocería como la $Q$ de Tobin: si una empresa vale $X$ unidades monetarias en el mercado de valores, nadie compraría en el mercado de bienes los activos reales necesarios para crear una empresa similar, si el precio de dichos activos fuese superior a $X$. Por el contrario, si el precio de dichos activos es inferior a $X$, todos querrían comprar los activos reales necesarios para crear una empresa idéntica y venderla inmediatamente en el mercado de valores por $X$ unidades monetarias.}; y,
    \item De la misma forma, un pesimismo generalizado en el mercado de valores reducirá la eficiencia marginal del capital y deprimirá el precio de las acciones, desincentivando así la inversión
\end{la}

\begin{adjustwidth}{2cm}{0pt}
\raggedright 
\textit{[L]as revaluaciones diarias de la bolsa de valores (...) ejercen inevitablemente influencia decisiva sobre la tasa de inversiones corrientes: porque no tiene sentido crear una nueva empresa incurriendo en un gasto mayor que aquel al que se puede comprar otra igual ya existente, mientras que hay un incentivo para gastar en un nuevo proyecto lo que podría parecer una suma extravagante, si puede venderse en la bolsa de valores con una ganancia inmediata. Por eso ciertas clases de inversiones se rigen por el promedio de las expectativas de quienes trafican en la bolsa de valores, tal como se manifiesta en el precio de las acciones, más bien que por las expectativas genuinas del empresario profesional.}\\
\raggedleft{\authorfont{KEYNES}, “\textit{Teoría General del Empleo, el Interés y el Dinero}”, 1936}
\end{adjustwidth}

Keynes se pregunta, a continuación, ¿Cómo se manejan, entonces, en la práctica, estas revaluaciones tan importantes de las inversiones existentes que se realizan cada día y aun cada hora? 

Según su parecer, la mayor parte de los inversores no analizan demasiado los fundamentos de las empresas en las que invierten, sino que se limitan a estimar una variación del precio a corto plazo.\footnote{En terminología actual, se podría decir que Keynes creía que la mayoría de los inversores utilizaban el análisis técnico frente al análisis fundamental a la hora de especular en el mercado de valores.} En efecto, apunta que para un inversor cualquiera: \textit{no es sensato pagar 25 por una inversión cuyo rendimiento probable se cree que justifica un valor de 30, si al mismo tiempo se supone que el mercado lo estimará en 20 tres meses después.} El comportamiento de los inversores, por tanto, se desliga de sus verdaderas creencias sobre la capacidad de las empresas de generar ingresos a largo plazo, y se orienta simplemente a tratar de anticipar los movimientos inmediatos de los precios de las acciones. Para ello, los inversores tratarán de anticiparse a lo que hagan los demás.\footnote{Keynes explica este comportamiento de manera muy didáctica en el mismo Capítulo 12 de la Teoría General, haciendo una analogía entre el mercado de valores y los concursos de belleza: 

\begin{adjustwidth}{2cm}{0pt}
\raggedright 
\textit{La inversión por profesionales puede compararse a esos concursos de los periódicos en que los concursantes tienen que seleccionar las seis caras más bonitas entre un centenar de fotografias, ganando el premio aquel competidor cuya selección corresponda más aproximadamente al promedio de las preferencias de los competidores en conjunto; de tal manera que cada concursante ha de elegir, no los semblantes que él mismo considere más bonitos, sino los que crea que serán más del agrado de los demás concursantes, todos los cuales observan el problema desde el mismo punto de vista (...) [H]emos alcanzado el tercer grado en el que dedicamos nuestra inteligencia a anticipar lo que la opinión promedio espera que sea la opinión promedio.}\\
\raggedleft{\authorfont{KEYNES}, “\textit{Teoría General del Empleo, el Interés y el Dinero}”, 1936}
\end{adjustwidth}} El hecho de que las expectativas estén basadas en percepciones tan subjetivas y cambiantes convierte la eficiencia marginal del capital en la principal responsable del ciclo económico (según KEYNES). En efecto, mientras que la propensión a consumir es una función bastante estable, la inversión fluctuará considerablemente ante cambios imprevistos en las expectativas. Así lo expresa así en el Capítulo 22 de la Teoría General: 

\begin{adjustwidth}{2cm}{0pt}
\raggedright 
\textit{En particular, encontraremos que las fluctuaciones en la propensión a consumir, en el estado de la preferencia por la liquidez y en la eficiencia marginal del capital han desempeñado su parte. Pero sugiero que el carácter esencial del ciclo económico y, especialmente, la regularidad de la secuencia de tiempo y de la duración que justifica el que lo llamemos ciclo, se debe sobre todo a cómo fluctúa la eficiencia marginal del capital.}\\
\raggedleft{\authorfont{KEYNES}, “\textit{Teoría General del Empleo, el Interés y el Dinero}”, 1936}
\end{adjustwidth}

Por tanto, las variaciones de la eficiencia marginal del capital son las causantes principales del ciclo, aunque apunta a que el efecto de los cambios en la eficiencia marginal del capital se amplifica a través del consumo. 

En efecto, las variaciones de la eficiencia marginal del capital dan lugar a variaciones en la inversión, que modifican la Demanda Agregada y con ella, la producción y el empleo. Sin embargo, el ciclo no termina aquí, ya que las variaciones de la renta (producción) y el empleo propician a su vez variaciones del consumo en el mismo sentido, amplificando el efecto inicial de las variaciones de la inversión.\footnote{Por ejemplo, ante un aumento de la inversión en una unidad monetaria, las empresas aumentarán la producción inicialmente en una unidad monetaria, que se destinará a pagar rentas a los hogares. Los hogares gastarán parte de ese aumento de la renta en bienes de consumo (concretamente, la propensión marginal a consumir), que provocará un nuevo aumento de la producción, y por consiguiente de la renta. Este proceso se repite hasta agotarse tras infinitas iteraciones.}
Este efecto (que ya había sido introducido previamente por R. KHAN), es lo que se conoce como efecto multiplicador y nos indica la derivada parcial de la producción del conjunto de la economía con respecto a la inversión total.\footnote{Recordemos que el gasto en consumo e inversión incluyen tanto el gasto privado como el público. De este modo, el multiplicador nos indica también cuanto aumentaría la producción total de la economía ante un incremento del gasto público.}

El efecto multiplicador dependerá de la propensión marginal a consumir y de otros factores. Si ignoramos todos los factores ajenos a la propensión marginal a consumir, el efecto multiplicador se puede calcular de acuerdo a la siguiente expresión, en la que $EM$ indica el efecto multiplicador y $c$ es la propensión marginal a consumir:

$$EM=\frac{1}{1-c}$$

En el análisis de KHAN se identifican otros factores que afectan al multiplicador, como:\footnote{\footnote{En los primeros cálculos de KHAN, el efecto multiplicador en una economía abiérta estaría comprendido entre dos y tres. A día de hoy, este cálculo se considera excesivamente elevado. Los modelos DSGE utilizados actualmente en macroeconomía no suelen encontrar un efecto multiplicador muy superior a uno. Las estimaciones empíricas realizadas con modelos SVAR y otras técnicas aumentan el rango de posibles valores del multiplicador, pero nunca en niveles tan elevados.}}
\begin{lnum}
    \item La respuesta del tipo de interés ante un aumento de la producción, ya que el multiplicador se reduciría si el aumento de la producción supone un incremento del tipo de interés-; o,
    \item El grado de apertura de la economía (el multiplicador sería menor en una economía abierta, al satisfacerse parte del aumento del consumo con importaciones).
\end{lnum}

\paragraph{Determinación de los salarios: Efectos en la economía real}
Como sabemos, el modelo neoclásico anticipa que en presencia de desempleo involuntario la competencia entre los trabajadores por los puestos de trabajo supondría una reducción de los salarios nominales que, \textit{ceteris paribus}, reducirían los salarios reales \textbf{(importante)}:
\begin{lnum}
    \item La reducción de los salarios nominales no afectaría directamente a los precios, ya que estos vienen determinados por la ecuación cuantitativa de Fisher: $MV=PY$; pero,
    \item Si la cantidad de dinero y su velocidad de circulación se mantienen inalteradas (como se presupone), los precios caerían solo en la medida en que crezca la producción;
    \item Pero, 
    \begin{lnum}
        \item Dado que la producción solo puede aumentar si aumenta el empleo; y,
        \item Que el empleo solo puede aumentar si baja el salario real
    \end{lnum}
    Se deduce necesariamente que, si el mecanismo de ajuste es tal y como dicen los neoclásicos, 
    \item Los salarios nominales deberán haber caído en proporción superior a los precios, lo que se traduce en una reducción de los salarios reales que acaba con el desempleo involuntario.
\end{lnum}

\textsc{Salarios reales}

Keynes hace dos objeciones a esta teoría:\footnote{En última instancia, como hemos visto, KEYNES defiende la idea de que el equilibrio macroeconómico agregado puede alcanzarse en presencia de desempleo involuntario.}
\begin{lnum}
    \item Observa que los salarios nominales son rígidos a la baja y que los trabajadores impedirán su reducción incluso cuando exista desempleo involuntario (a menudo organizados en sindicatos);\footnote{En realidad, esta observación no es una aportación de Keynes, ya que los economistas neoclásicos ya explicaban el fenómeno del desempleo argumentando que su única causa posible es la negativa de los trabajadores a aceptar un ajuste del salario real en el mercado de trabajo (a través del salario nominal).} sin embargo, argumenta que la causa del desempleo involuntario no puede ser la oposición de los trabajadores a reducir el salario real porque cuando dicha reducción se produce a través de un incremento del nivel general de precios, los trabajadores no suelen oponer resistencia en forma de huelgas o protestas. La causa real de la rigidez a la baja en los salarios nominales se encuentra, en el deseo de los trabajadores de mantener su posición relativa en la distribución salarial del conjunto de los trabajadores y no en impedir un ajuste a la baja del salario real del conjunto de la economía.\footnote{Según Keynes, aunque lo relevante para un trabajador sea siempre el salario real, lo que negocia con su empleador será exclusivamente el salario nominal. Esto se debe a que en un mercado de trabajo heterogéneo, el trabajador no conoce con exactitud cuál debería sería su salario justo o de equilibrio: su productividad marginal de trabajo; por tanto, opta por fijarse en lo que ganan el resto de trabajadores y negocia su propio salario en función del salario de los demás.\\
    Por ejemplo, un trabajador cualificado podría estimar que su salario debería ser al menos un 35\% superior al de los trabajadores no cualificados. De este modo, este trabajador cualificado nunca aceptará una reducción de su salario nominal hasta que no hayan bajado previamente los salarios nominales de los trabajadores no cualificados, provocando que los salarios nominales sean rígidos a la baja.\\
    Dado que el nivel general de precios es idéntico para todos los trabajadores, las variaciones del mismo nunca afectarán a la posición relativa de un trabajador en el conjunto de la distribución salarial y, en consecuencia, el nivel de precios será irrelevante en la negociación del salario nominal. El hecho de que los trabajadores negocien salarios nominales en vez de reales se conoce habitualmente como \textit{ilusión monetaria}. En cualquier caso, habrá desempleo involuntario siempre que el salario real sea superior a la desutilidad marginal del trabajo, independientemente del salario nominal negociado en el mercado de trabajo.\\
    \begin{adjustwidth}{2cm}{0pt}
    \raggedright 
    \textit{Desde el momento que existe movilidad imperfecta del trabajo y que los salarios no tienden a producir igualdad precisa de ventajas netas para diferentes ocupaciones, cualquier individuo o grupo de individuos que consienta una reducción de sus salarios nominales en relación con otros, sufrirá una disminución relativa de sus salarios reales, cosa que basta para justificar su resistencia a ella. Por el contrario, sería impracticable oponerse a toda reducción de los salarios reales debida a un cambio en el poder adquisitivo del dinero, que afecta a todos los trabajadores por igual; y de hecho, por lo general no se opone resistencia a esta clase de fenómenos, a menos que sean extremos.}\\
    \raggedleft{\authorfont{KEYNES}, “\textit{Teoría General del Empleo, el Interés y el Dinero}”, 1936}
    \end{adjustwidth}\\
    Como corolario, esta idea rechaza implícitamente la visión marxista de la acción sindical como un elemento de lucha de clases entre capitalistas y trabajadores por la determinación del salario real del conjunto de la clase trabajadora. En realidad, la acción sindical es el terreno de lucha entre los sindicatos de distintas empresas o sectores para aumentar su salario real relativo. Según Keynes, el salario real vendrá dado únicamente por el volumen de empleo (determinado por la Demanda Efectiva), siendo irrelevante la acción sindical:
    \begin{adjustwidth}{2cm}{0pt}
    \raggedright 
    \textit{El efecto de la unión de un grupo de trabajadores consiste en proteger su salario real relativo. El nivel general de los salarios reales depende de otras tijeras del sistema económico.}\\
    \raggedleft{\authorfont{KEYNES}, “\textit{Teoría General del Empleo, el Interés y el Dinero}”, 1936}
    \end{adjustwidth}}
\end{lnum}

\textsc{Salarios nominales}

Tras explicar las causas de la rigidez nominal de los salarios, Keynes destina el Capítulo 19 de la Teoría General a explicar los efectos esperados de una reducción de los salarios nominales sobre el empleo. 

Su conclusión es clara: la flexibilidad de los salarios nominales no constituye un mecanismo de ajuste del mercado de trabajo, es decir, no necesariamente elimina el desempleo involuntario, pudiendo mantenerlo inalterado o incluso incrementarlo. 

El efecto de una reducción de los salarios nominales sobre el empleo y los salarios reales es incierto, según KEYNES, pues tanto el empleo como el salario real vendrán determinados en última instancia por la Demanda Efectiva, es decir, por la propensión a consumir, el tipo de interés y la eficiencia marginal del capital: 

\begin{adjustwidth}{2cm}{0pt}
    \raggedright 
    \textit{De este modo, la baja de los salarios nominales no tenderá a aumentar la ocupación durante mucho tiempo, excepto en virtud de sus repercusiones, ya sea sobre la propensión a consumir de la comunidad en su conjunto, sobre la curva de la eficiencia marginal del capital, o sobre la tasa de interés. No hay más modo de analizar el efecto de una reducción de los salarios nominales que el de examinar sus posibles efectos sobre estos tres elementos}\\
    \raggedleft{\authorfont{KEYNES}, “\textit{Teoría General del Empleo, el Interés y el Dinero}”, 1936}
\end{adjustwidth}

Keynes encuentra efectos positivos y negativos de una reducción de salarios nominales sobre los determinantes de la Demanda Efectiva:
\begin{la}
    \item \textit{Efectos positivos.} Reconoce que toda reducción de salarios nominales:
    \begin{lnum}
        \item \textit{Reducción del tipo de interés real}. Reducirá el coste marginal de producción de los bienes, dando lugar a una reducción del nivel general de precios y ésta reducción de precios reducirá la demanda de dinero por motivo de transacción, provocando una caída del tipo de interés nominal y, dada una inflación esperada, del tipo de interés real, aumentando la inversión.
        \item \textit{Mejora de las expectativas de la rentabilidad de las inversiones}. Otro efecto positivo sobre la inversión sería la posible mejora de las expectativas de los inversores, que podrían hacer juicios más optimistas sobre la rentabilidad de las inversiones tras la caída de los salarios nominales. 
        \item \textit{Aumento de la demanda externa}. En el caso de una economía abierta, la reducción de salarios y precios aumentaría también la demanda externa, gracias al incremento de la competitividad frente al exterior. 
    \end{lnum} 
    \item \textit{Efectos negativos}. Keynes recalca que una reducción de salarios nominales y precios dará lugar a:
    \begin{lnum}
        \item \textit{Efectos redictributivos negativos y reducción de la propensión a consumir}. Una alteración en la distribución de la renta con efectos perniciosos ya que, según KEYNES, la caída de los salarios y precios supondrá una transferencia de renta real, de quienes perciben salarios a quienes tienen un ingreso fijo en términos monetarios (bonistas, etc.). Esta redistribución de recursos tendería a reducir la propensión a consumir media del conjunto de la comunidad.
        \item \textit{Reducción de la inversión por encarecimiento de la deuda}. La caída de los precios daría lugar a un incremento del valor de las deudas, reduciendo la capacidad de endeudamiento de los empresarios y aumentando las quiebras, provocando una caída de la inversión.
        \item \textit{Reducción de la inversión por expectativas deflacionarias}. La caída de precios en un periodo podría dar lugar a expectativas de deflación en los periodos siguientes, aumentando así el tipo de interés real y deprimiendo la inversión.
        \item \textit{Encarecimiento de las importaciones}. En una economía abierta, la caída de los precios internos supondría un aumento del valor real de las importaciones, compensando total o parcialmente el efecto positivo de la ganancia de competitividad exterior sobre la demanda externa.\footnote{Keynes parece anticiparse a lo que posteriormente sería conocido como \textit{condición de Marshall-Lerner}, que establece el requisito que debe cumplirse para que una ganancia de competitividad-precio incremente el saldo de la balanza comercial.} 
    \end{lnum} 
\end{la}

El efecto final de una reducción de salarios y precios dependerá, por tanto, de la intensidad de los efectos anteriores. En una situación de depresión económica, KEYNES estima que los efectos de la deflación son especialmente nocivos, por lo que no recomendará abordar la crisis con una reducción de salarios y precios, sino con políticas monetarias y fiscales expansivas. 

Especialmente nocivo sería el caso de una economía deprimida con deflación en el que el tipo de interés nominal hubiera bajado ya hasta el 0\%. En este caso, la política monetaria sería inútil y la deflación de salarios y precios sería especialmente perniciosa, siendo la política fiscal la única vía de salida a la crisis. En estas circunstancias, rechazará la flexibilidad salarial como remedio natural para alcanzar el pleno empleo.

Una vez vistos todos los principales elementos de la Teoría General, se puede consultar en los Anexos 2.1 y 2.2. una ilustración gráfica del \textit{modelo keynesiano} en su conjunto.

\subsubsection{El papel del Estado y la Política económica}
Keynes justifica la intervención del sector público en la economía para corregir lo que para él serán los principales fallos del sistema capitalista: 
\begin{lnum}
    \item El desempleo involuntario; y,
    \item La desigualdad excesiva en la distribución de la renta y la riqueza.
\end{lnum}

\begin{adjustwidth}{2cm}{0pt}
    \raggedright 
    \textit{Los principales problemas e inconvenientes de la sociedad económica en que vivimos son su incapacidad para procurar la ocupación plena y su arbitraria y desigual distribución de la riqueza y los ingresos}\\
    \raggedleft{\authorfont{KEYNES}, “\textit{Teoría General del Empleo, el Interés y el Dinero}”, 1936}
\end{adjustwidth}

Defiende la intervención del sector público en la economía para aumentar la demanda agregada en la medida suficiente para alcanzar el pleno empleo, siendo el primer economista en aportar una visión sistematizada de las políticas de demanda. 

A diferencia de los economistas neoclásicos, no abogará por la flexibilidad de los salarios, sino por el mantenimiento de un nivel estable de salarios nominales que garantice la estabilidad de precios. En Capítulo 19 afirma lo siguiente (Capítulo 24): 

\begin{adjustwidth}{2cm}{0pt}
    \raggedright 
    \textit{Suponer que la política de salarios flexibles es un auxiliar correcto y adecuado de un sistema que en conjunto corresponde al tipo del laissez-faire, es lo opuesto a la verdad}\\
    (Y continua), \textit{[O]pino ahora que el mantenimiento de un nivel general estable de salarios nominales es, en general, la política más aconsejable para un sistema cerrado; al tiempo que la misma conclusión será válida para un sistema abierto, a condición de que pueda lograr el equilibrio con el resto del mundo por medio de fluctuaciones en los tipos de cambio.}\\
    \raggedleft{\authorfont{KEYNES}, “\textit{Teoría General del Empleo, el Interés y el Dinero}”, 1936}
\end{adjustwidth}

Keynes considera que en presencia de desempleo involuntario la única ventaja clara que puede ocasionar una reducción de los salarios nominales es la reducción del tipo de interés. Sin embargo, opina que dicho efecto se puede conseguir igualmente a través de una política monetaria expansiva: 

\begin{adjustwidth}{2cm}{0pt}
    \raggedright 
    \textit{Podemos, (...), producir precisamente los mismos efectos sobre la tasa de interés reduciendo los salarios, al mismo tiempo que dejamos invariable la cantidad de dinero, que aumentando la cantidad (de dinero) y dejando sin variación el nivel de salarios.}\\
    \raggedleft{\authorfont{KEYNES}, “\textit{Teoría General del Empleo, el Interés y el Dinero}”, 1936}
\end{adjustwidth}

Dado que las reducciones de salarios nominales pueden generar efectos adversos, Keynes afirma que no tiene sentido preferir una política de salarios flexibles frente a la utilización activa de la política monetaria. Sin embargo, Keynes tampoco es especialmente optimista acerca de la efectividad de la política monetaria, ya que considera que la demanda de dinero puede ser muy elástica y estar afectada por cambios bruscos e impredecibles, que impidan controlar el tipo de interés. 

Además, considera que los cambios en los tipos de interés podrían ser insuficientes para contrarrestar las grandes fluctuaciones en las expectativas de los inversores y empresarios. Por tanto, abogará por la utilización de la política fiscal, sobre todo a través de la inversión pública:\footnote{No lo afirma directamente en la Teoría General, pero se deduce que considera más efectiva la puesta en marcha de inversiones públicas que las reducciones de impuestos a la hora de incrementar la demanda agregada. Ello se debe a que parte de la renta disponible que generan las bajadas de impuestos será ahorrada, dada la ley psicológica fundamental.}

\begin{adjustwidth}{2cm}{0pt}
    \raggedright 
    \textit{Por mi parte, soy ahora un poco escéptico respecto al éxito de una política puramente monetaria dirigida a influir sobre la tasa de interés. Espero ver al Estado, que está en situación de poder calcular la eficiencia marginal de los bienes de capital a largo plazo sobre la base de la conveniencia social general, asumir una responsabilidad cada vez mayor en la organización directa de las inversiones, ya que probablemente las fluctuaciones en la estimación que haga el mercado de la eficiencia marginal de las diferentes clases de capital, (...) , serán demasiado grandes para contrarrestarlas con alguna modificación factible de la tasa de interés}\\
    \raggedleft{\authorfont{KEYNES}, “\textit{Teoría General del Empleo, el Interés y el Dinero}”, 1936}
\end{adjustwidth}

En definitiva, cree en la eficiencia asignativa del mercado, pero aboga por la intervención en aras de alcanzar el pleno empleo:

\begin{adjustwidth}{2cm}{0pt}
    \raggedright 
    \textit{En lo que ha fallado el sistema actual ha sido en determinar el volumen del empleo efectivo y no su dirección}\\
    \raggedleft{\authorfont{KEYNES}, “\textit{Teoría General del Empleo, el Interés y el Dinero}”, 1936}
\end{adjustwidth}

\subsection{Representación gráfica del Modelo \textit{keynesiano}}

\subsubsection{Efecto del atesoramiento} 

La representación gráfica del modelo keynesiano se puede hacer a través de los siguientes ocho gráficos. 

\imagenfit{Atesoramiento-Keynes.png}{Efecto del atesoramiento en el modelo Keynesiano}{Apuntes ICEX-CECO. Tema 3.A.4}

Los gráficos representan lo siguiente:
\begin{lnum}
    \item El primero de ellos muestra el mercado de bienes, que detennina el nivel de empleo;
    \item El gráfico inferior representa el mercado de trabajo en ténninos nominales. En este mercado tenemos una demanda nominal de trabajo que coincide con el valor de la productividad marginal,\footnote{El valor de la productividad marginal del trabajo sería el ingreso marginal de las empresas que obtendrían las empresas al vender una unidad adicional, multiplicado por la productividad marginal fisica del trabajo.} y un salario nominal que consideramos rígido;\footnote{Este salario nominal está determinado por la negociación de trabijadores y sindicatos en el mercado de trabajo.}
    \item El gráfico inmediatamente inferior nos muestra el mercado de trabajo expresado en términos reales. En este mercado aparece la demanda real de trabajo, la productividad marginal física del trabajo, y la oferta de trabajo;
    \item El siguiente gráfico inferior es la función de producción, que relaciona el empleo con la producción;
    \item A la derecha, tenemos un gráfico auxiliar que solo sirve para traspasar la variable Y del eje de ordenadas al eje de abscisas;
    \item El gráfico inmediatamente superior nos indica el nivel general de precios,\footnote{La curva que mostramos es otra forma de oferta agregada, pero es importante no confundirla con la del primer gráfico.
    } que tiene pendiente positiva hasta alcanzar el pleno empleo, convirtiéndose en vertical en el nivel de producción de pleno empleo;\footnote{En este modelo, el nivel de precios no se detenninará directamente por la ecuación cuantitativa de Fisher, sino que vendrá dado por el coste marginal de producción (suponiendo competencia imperfecta), que a su vez se calcula como el salario nominal entre la prouctividad marginal física del trabajo:
    $$P=CMg=\frac{W}{PMg_L}$$
    El nivel de precios tendrá pendiente positiva hasta que el nivel de producción $Y$ alcance el correspondiente nivel de producción de pleno empleo. Ello se debe a que a medida que aumenta la producción, se incrementa el número de trabajadores, reduciéndose así la productividad marginal. El resultado es que a medida que crece la producción, sube el coste marginal, y por tanto los precios. Una vez alcanzado el pleno empleo, es imposible seguir aumentando la producción, porque no hay acuerdo posible entre empresas y trabajadores para aumentar el volumen de empleo (la productividad marginal del trabajo sería inferior al salario que exigen los trabajadores), por lo que la curva de precios será vertical a partir de ese punto.
    }
    \item El siguiente gráfico superior representa el ahorro y la inversión en función del tipo de interés real, que por sencillez vamos a suponer similar al tipo de interés nominal (es decir, suponemos que la inflación esperada es cero). La curva de inversión representa la eficiencia marginal del capital y se desplaza con los cambios en las expectativas de los inversores y de los empresarios. La curva de ahorro se desplaza con los cambios en el nivel de producción, $Y$ (renta); y,
    \item Por último, contamos con el mercado de dinero/bonos a la derecha. Este mercado determina el tipo de interés. El tipo de interés equilibrará la preferencia por la liquidez con la oferta monetaria.
\end{lnum}

Partiendo de un equilibrio inicial con pleno en $0$, suponemos que se produce un aumento repentino de la preferencia por la liquidez; por ejemplo, porque los agentes económicos esperan que el precio de los bonos va a caer de manera inminente, por el motivo que sea. Ante esta expectativa de caída de precios en el mercado de bonos:
\begin{lnum}
    \item Los agentes se anticipan y venden parte de sus bonos en el periodo actual. lo que conduce a un desplazamiento de la demanda de dinero a la derecha, que nos llevaría al equilibrio $1$ en el mercado de dinero;
    \item La subida de tipos de interés supondría una reducción de la inversión, que se materializaría en un desplazamiento de la $DA$ hacia abajo, dando lugar a:
    \begin{lnum}
        \item Una caída del empleo; y, por consiguiente,
        \item De la producción.
    \end{lnum} 
    \item La caída de la producción supondría una caída de la demanda de dinero, llevando al mercado de dinero al equilibrio $2$.
\end{lnum}   

En el equilibrio 2,\footnote{Por sencillez gráfica, se omiten todos los pasos intermedios entre el equilibrio $1$ y el equilibrio $2$.
} la $DA$ se habrá desplazado hacia abajo con respecto al equilibrio $0$. Se puede ver que la caída total de las ventas es inferior al desplazamiento hacia abajo de la curva debido al efecto multiplicador:
\begin{lnum}
    \item El volumen de empleo cae, dando lugar a un incremento de la productividad marginal del trabajo, que reduce los precios;
    \item La reducción de precios produce un desplazamiento hacia la izquierda de la demanda nominal de trabajo y un incremento del salario real;
    \item El ahorro y la inversión se equilibran de nuevo en $2$, pero este equilibrio se produce porque cae la renta $Y$, y con ella, el ahorro ($S$). 
\end{lnum}

Es decir, en el modelo keynesiano, la variable que equilibra el ahorro y la inversión no es el tipo de interés sino la renta ($Y$).

El resultado del atesoramieno en este modelo será la aparición de desempleo involuntario. Vemos que:
\begin{lnum}
    \item Las variables monetarias (demanda de dinero y oferta monetaria) afectan a las variables reales a través del tipo de interés, rompiendo, por tanto, la neutralidad del dinero y la dicotomía clásica;
    \item El salario real es independiente del salario nominal que negocian los trabajadores y, por tanto, hay variaciones del salario real y del empleo sin necesidad de que haya habido un incremento del poder sindical o nuevas regulaciones de salario mínimo, etc;
    \item El nivel de precios se determina directamente como función del coste marginal; y,
    \item La ecuación cuantitativa de Fisher se sigue cumpliendo, pero cambia el comportamiento de las variables. En este ejemplo anterior, caen la producción y el empleo, pero la oferta monetaria se mantiene constante. Ello implica necesariamente que la velocidad de circulación del dinero no es estable y no depende solo de la tecnología de pagos, sino que es afectada por las expectativas y otras variables. En este modelo, toda variación de la demanda agregada que no se produzca por una variación en la oferta monetaria (política fiscal, \textit{animal spirits}, demanda de dinero, etc.) afecta a la velocidad de circulación del dinero. Se rechaza, por tanto, que esta variable sea estable y venga determinada por factores como la tecnología de pagos, etc; y,
    \item Por ultimo, es muy destacable de este modelo que el ahorro y la inversión no determinan simultáneamente el tipo de interés. Por el contrario, el tipo de interés se determina en el mercado de bonos/dinero. El tipo de interés determina la inversión, la inversión determina la renta y la renta determina el ahorro. La inversión, por tanto, determina el ahorro. Ahorro e inversión se igualan con variaciones en la renta y no del tipo de interés. Se incunple, por tanto, la Ley de Say, pues puede haber perfectamente un exceso de ahorro sobre inversión que provoque una escasez de demanda generalizada. Esta escasez solo se corregirá cuando la producción y el ahorro caigan lo suficiente.
\end{lnum}

\subsubsection{Efecto de un descenso de los salarios nominales}


\clearpage
\section{La economía post-keynesiana}
\puntosclave{

Las contribuciones de autores como HICKS y HANSEN permiten reconciliar las ideas de Keyes con la economía neoclásica, dando lugar a la denominada \textit{Síntesis Neoclásica} [\authorfont{Ver Tema 3.A.5}]. Este nuevo enfoque constituyó la ortodoxia macroeconómica entre los años 50's y los años 70's. 

Tras la publicación de la Teoría General, algunos seguidores y colaboradores suyos, la mayoría discípulos en la Universidad de Cambridge, rechazan frontalmente las contribuciones de la síntesis neoclásica. Éstos se centrarán en extender lo que ellos consideraron la verdadera interpretación fidedigna de las nociones de la Teoría General a otros campos de la economía, como la Teoría del crecimiento económico, la distribución de la renta y el funcionamiento del sistema financiero. A principios de los años 70's, estas contribuciones críticas se clasificaron como \textit{escuela postkeynesiana}. Algunos autores catalogados dentro de esta escuela son J. ROBINSON, N. KALDOR, M. KALECKI, P. SRAFFA y H. MINSKY.

En la década de los años 60's y 70's, hubo otros autores descontentos con la falta de
microfundamentos de la sínstesis neoclásica. La separación entre la microeconomía
neoclásica y macroeconomía de la síntesis neoclásica impedía recoger con exactitud la
contribución de KEYNES. Autores como R. CLOWER, A. LEIJONHUFVUD, E.
MALINVAUD y J. BÉNASSY forman parte de este grupo de economistas, conocidos
posteriormente como Keynesianos del Desequilibrio.
}

La publicación de la Teoría General dio lugar a un intenso debate y análisis. PIGOU, maestro de KEYNES en Cambridge, no tardó en elaborar una dura crítica en su artículo \textit{La Teoría General del Empleo de Mr. J. M. Keynes}, 1936. En noviembre del mismo año, el Quarterly Journal of Economics dedicó varios artículos a comentar la Teoría General, con contribuciones tan relevantes como las de TAUSSING, LEONTIEFF y ROBERTSON. 

Los desarrollos posteriores a la publicación de la Teoría General son cuando menos homogeneos. Por esta razón, su tratamiento esquemático resulta muy complejo, aunque podemos utilizar la siguiente triple división canónica:\footnote{Cabría también hablar de los Keynesianos estructuralistas o desarrollistas, corriente afín pero no auto-definida como seguidora de las tésis de KEYNES. Sus reprentantes e ideas principales van encaminadas al Sur Global, siendo muy importante en América Latina}
\begin{la}
    \item Keynesianismo ortodoxo / Síntesis neoclásica. Esta corriente tratará de compatibilizar a KEYNES con la visión neoclásica y es objeto de estudio en el [\authorfont{Ver Tema 3.A.5}], siendo sus principales exponentes SAMUELSON, HICKS, MODIGLIANI, KLEIN, HANSEN, TOBIN y otros;
   \item Los post-keynesianos. Corriente heterodoxo que considera una \textit{traición} la reinterpretación keynesiana que hacen los defensores de la síntesis y cuyos exponentes principales serán J. ROBINSON, KALDOR, SRAFFA, MINSKY, KHAN y DAVIDSON. A su vez, podemos encontrar tres vertientes diferentes:
   \begin{la}
       \item La Escuela de Cambridge post-keynesiana (o segunda Escuela de Cambridge);
       \item Los neorricardianos o sraffianos; y,
       \item Los institucionalistas americanos. 
   \end{la}
   \item La Nueva Economía Keynesiana. Corriente que aparece tras la crisis del keynesianismo ortodoxo por la estanflación y las críticas monetaristas y cuyas figuras principales serán varíadas a lo largo del tiempo (por ser una corriente viva y activa hoy en día). Será objeto de estudio en el [\authorfont{Ver Tema 3.A.7}].
\end{la}

Dentro de esta Tema desarrollaremos las princpales teorías de la Economía post-keynesiana, por ser las otras dos corrientes propias de otras partes del temario.  A finales de los años 30's, algunas voces alternativas comenzaron a cuestionar la interpretación conciliadora del pensamiento de KEYNES, negando que pudiese ser asimilado como un caso especial del pensamiento económico neoclásico. ROBINSON y KALDOR se encontraron dentro del grupo economistas que insistieron en esta línea, marcando la dirección de lo que posteriormente se convertiría en la economía postkeynesiana.\footnote{El término postkeynesiano no ha tenido el mismo significado a lo largo del tiempo. En los años 50's era empleado a menudo en un sentido cronológico, describiendo una amplia variedad de desarrollos posteriores a la Teoría General. En Estados Unidos, este término también comprendía a la síntesis neoclásica. La utilización del término postkeynesiano en un sentido doctrinal se encuentra por primera vez en el artículo \textit{La Reconstrucción de la Economía Política: una Introducción a la Economía Post-Keynesiana}, publicado por KREGEL en 1973.} Los autores de este escuela de pensamiento van a insistir en la inestabilidad instrínseca del sistema económico y van a extender el pensamiento de KEYNES a otras áreas de la economía.

\subsection{Breve historia y marco teórico común} 
Como su nombre indica, el origen de la teoría poskeynesiana se remonta al trabajo de KEYNES. Sin embargo, los poskeynesianos contemporáneos fueron fuertemente influenciados por el economista polaco KALECKI, a quien a menudo se le atribuye haber establecido en la década de 1930, al mismo tiempo que KEYNES, pero de manera independiente, los principios fundamentales de funcionamiento de una economía conducida por la demanda.\footnote{En efecto, según ROBINSON, que es con KALDOR y SRAFFA una de las grandes fuentes de inspiración de la teoría poskeynesiana, KALECKI produjo una versión más coherente que la de KEYNES. } Podemos distinguir algunos momentos cruciales del pensamiento post-keynesiao:
\begin{lnum}
    \item Aunque economistas como ROBINSON publicaron en los años de 1930 y 1940 lo que hoy podría llamarse trabajo poskeynesiano, fue con la publicación de su libro de 1956, \textit{Acumulación de capital}, que realmente se puede discernir verdaderamente una visión de conjunto que constituye una alternativa a la economía neoclásica dominante. \\En ese momento, esta visión, que proponía notablemente una teoría alternativa de la distribución del ingreso entre salarios y ganancias, conocida como teoría neokeynesiana o neocambridgeana.
    \item La controversia sobre el capital de los dos Cambridge de los años de 1960, la que cuestiona la relación entre la escasez de capital y su precio, es un segundo momento clave, en el que los economistas keynesianos afiliados a la Universidad de Cambridge en Inglaterra se dan cuenta de que su visión de la economía, de hecho, es una teoría distinta a la que prevalece entre los economistas keynesianos del Instituto Tecnológico de Massachusetts (ubicado en la ciudad de Cambridge, frente a Boston), y en casi cualquier otro lugar de EU;
    \item El tercer momento clave es la publicación del artículo de EICHNER y KREGEL (1975), en la prestigiosa revista\textit{ Journal of Economic Literature}, cuyo título contiene el nombre de teoría poskeynesiana, y este artículo afirmaba que se trata de un nuevo paradigma. 
\end{lnum}

También en la década de 1970 la teoría marxista experimentó un renacimiento, y fue cuando los keynesianos estadounidenses que afirmaban ser una versión radical de Keynes se organizaron y crearon las dos primeras revistas con el objetivo de difundir el pensamiento poskeynesiano: \textit{Cambridge Journal of Economics}, y \textit{Journal of Post Keynesian Economics}. Posteriormente, aparecerán otras revistas con objetivos similares, se crearán cursos de teoría poskeynesiana, surgirán escuelas de verano en teoría poskeynesiana, se multiplicarán las conferencias y trabajos tanto en Europa, Estados Unidos y en América Latina. 

La realidad es que la teoría poskeynesiana, como todas las teorías económicas heterodoxas, se encuentra en una situación precaria y muy minoritaria, debido a la hegemonía e intolerancia ejercida desde la década de 1980 por los defensores de la teoría neoclásica en los departamentos de la ciencia económica. 

\subsubsection{Las Características de la economía}
\paragraph{Para los heterodoxos en general}
Como todas las teorías heterodoxas,\footnote{\textbf{Reformular.-} Los lectores que no estén familiarizados con la teoría poskeynesiana, pero que deseen saber más, pueden acceder al libro introductorio en francés de Lavoie et al. (2021) o a la versión en inglés de King (2015). Los lectores más decididos pueden consultar dos colecciones de textos en francés, la editada por Piegay y Rochon (2003) que se centra en el dinero, y la editada por Berr et al. (2018) que trata todos los aspectos de la teoría poskeynesiana y que presenta a los fundadores de esta escuela depensamiento, como se evoca en la sección anterior. El libro de Lavoie (2014) es una recopilación avanzada, y en inglés de tesis y modelos poskeynesianos, recomendada para estudiantes de posgrado. También hay que destacar dos números especiales de la revista quebequense L’Actualité économique, dedicada a la teoría económica, un número doble publicado en 1982 y dirigido por Jacques Henry y Mario Seccareccia, y otro en 1992, dirigido por Marc Lavoie.} la economía poskeynesiana se basa en un conjunto de supuestos que la distingue de la economía dominante, la economía neoclásica.  Estos supuestos constituyen los conceptos esenciales de una escuela de pensamiento y preceden a la constitución de los supuestos y teorías que van a ser elaboradas.\footnote{Estos presupuestos son las creencias metafísicas que gobiernan un paradigma o un programa de investigación. Diferentes economistas o metodólogos no entienden probablemente la identificación exacta de estos supuestos, pero afirmamos que se encuentran, casi de manera idéntica, en la pluma de numerosos autores heterodoxos. En consecuencia, vamos a identificarlos brevemente, antes de pasar a las características que son específicas de la teoría poskeynesiana en comparación con otras escuelas de pensamiento heterodoxas.} 

La teoría neoclásica y las escuelas heterodoxas se distinguen por cuatro rasgos metodológicos esenciales, en pares opuestos, a los que hay que añadir un rasgo político:
\begin{la}
    \item Respecto a los rasgos metodológicos:
    \begin{lnum}
        \item Epistemología realista (\textbf{realismo}) frente a una epistemología instrumental: los autores heterodoxos creen que las simplificaciones necesarias para el análisis económico deben, no obstante, ser descriptivas; es decir: deben representar el mundo como es, no un mundo imaginario. Este es el presupuesto del realismo de las hipótesis;
        \item \textbf{Holismo metodológico} frente al individualismo metodológico. Esto se reflejaperfectamente en las muchas paradojas macroeconómicas, la más conocida de las cuales es la paradoja del ahorro de KEYNES;
        \item \textbf{Racionalidad procedimental} frente a la racionalidad ilimitada. Los economistas heterodoxos en su conjunto rechazan la hiperracionalidad presente en los economistas neoclásicos, los cuales presumen que todos los agentes económicos saben cómo funciona realmente la economía;\footnote{Los heterodoxos se inspiran más bien en la visión de Herbert Simon, tal como fue elaborada en psicología por Gerd Gigerenzer (2009), según la cual las personas solo buscan soluciones satisfactorias, adoptando para ello normas, convenciones, reglas de comportamiento o incluso hábitos que les permitan tomar decisiones sin perder demasiado tiempo.} y,
        \item \textbf{Concepción de la economía de producción} frente a la concepción de la economía basada en el intercambio y la escasez. Podemos dar dos interpretaciones: (i) los heterodoxos no asumen que las economías operan a su capacidad total y con pleno empleo; y, (ii) las cuestiones esenciales se relacionan con el inicio (creación) de un excedente durante la producción y no con la asignación óptima por el intercambio de recursos existentes.
    \end{lnum}
    \item Respecto al rasgo político:\textbf{generalmente a favor de la intervención estatal por tendencia natural a la inestabilidad} frente a la capacidad de autorregulación natural de los mercados (en un contexto ideal).
\end{la}


\paragraph{Para los poskeynesianos en particular}
Ahora se puede preguntar, en qué se diferencia la teoría poskeynesiana de otras teorías heterodoxas. Nos basaremos en la opinión de Eckhard Hein (2017) para identificar cinco características clave:
\begin{lnum}
    \item \textit{Principio de demanda efectiva}.\footnote{\textbf{Recordar.-} El principio de demanda efectiva requiere que la producción se ajuste a la demanda. La economía está conducida por la demanda, no por las limitaciones que surgen de la oferta y las dotaciones existentes.\\
    Muchos economistas reconocen la validez de este principio, pero solo cuando se trata de corto plazo. Por el contrario, los economistas marxistas y neoclásicos, por ejemplo, siguen convencidos de que, a largo plazo, la economía se ve impulsada por limitaciones del lado de la oferta. Los economistas poskeynesianos se distinguen por su negativa a creer que los factores del lado de la oferta pueden ser una limitación, incluso a largo plazo.} Para los poskeynesianos, el principio de demanda efectiva se aplica en todo momento, y la inversión determina el ahorro causalmente. Por lo tanto, existe un número infinito de posibles equilibrios de largo plazo, que dependen de las restricciones impuestas por la demanda y las instituciones establecidas. Los factores del lado de la oferta finalmente se ajustarán, por cambios en los movimientos de la población o mediante la aceleración del progreso técnico.
    \item \textit{Tiempo}. Para los postkeyneasianos, la verdadera escasez es el tiempo. Estos autores distinguen entre el tiempo lógico, que no tiene profundidad, y el tiempo histórico, que es irreversible. El camino que se toma después de cualquier modificación es de suma importancia, porque la tendencia a largo plazo es solo el resultado de la sucesión de una serie de períodos cortos.\footnote{Los poskeynesianos plantean la necesidad de construir modelos dinámicos, que tengan en cuenta la evolución en el tiempo de las existencias de activos físicos, deudas y riqueza financiera, y que puedan explicar la reorganización de la estructura productiva. Este es el tiempo dinámico. El equilibrio que se alcanza durante un período prolongado no es independiente del camino recorrido. Estas ideas, antes consideradas poco conducentes a la formalización, están ahora en el corazón de los últimos desarrollos matemáticos no lineales construidos en torno a las nociones de histéresis, dependencia del camino tomado, irreversibilidad y efectos de bloqueo, que con mayor frecuencia implican la existencia de equilibrios múltiples.}
    \item \textit{Economía monetizada}. Una economía de producción monetaria que excluye las relaciones de trueque y requiere que las transacciones se realicen en moneda en la unidad de cuenta determinada por el estado, con contratos expresados en esta unidad de cuenta y mantener activos reales en su forma monetizada: activos financieros. La moneda y el crédito juegan un papel fundamental; están integrados en el proceso de producción desde el principio.\footnote{Por tanto, el papel de los bancos es fundamental porque brindan los avances que requieren las empresas productoras para iniciar la producción o para incentivar a los hogares a consumir. Por ello, los poskeynesianos conceden gran importancia a los flujos de crédito, que ayudan a explicar la evolución de la producción nacional, así como a los stocks de deuda, que ayudan a explicar las caídas repentinas de la renta y las crisis financieras. Los llamados modelos coherentes de flujo de valores (SFC), propuestos por GODLEY y LAVOIE (2007) y ahora extremadamente populares entre los economistas poskeynesianos, son particularmente adecuados para analizar estas economías monetizadas. Estos modelos integran flujos financieros y stocks y transacciones entre diferentes sectores, teniendo en cuenta los flujos reales de actividad económica. Su estructura asegura que no haya agujeros negros, con cada transacción sujeta a un asiento contable cuádruple. Permiten introducir las consecuencias financieras, por ejemplo, sobre deudas, de las distintas transacciones asociadas a los flujos reales.}
    \item \textit{La incertidumbre fundamental o radical}. Evidentemente, está ligada a la de la época histórica y de racionalidad razonable, dotadas de un conocimiento limitado. En el tiempo histórico, el futuro no puede ser idéntico al presente o al pasado: el mundo no es ergódico, lo que significa que promedios y fluctuaciones observadas en el pasado no pueden repetirse de manera idéntica para cada período de tiempo.\footnote{Esto explica por qué los poskeynesianos, aunque se dedican a estudios econométricos, siguen siendo escépticos en cuanto a la validez o generalidad de los resultados empíricos así obtenidos. A diferencia de los economistas neoclásicos que creen, como LUCAS, el líder de los nuevos clásicos, que en el contexto de la incertidumbre radical el razonamiento económico ya no tiene ningún valor, los economistas poskeynesianos piensan que, excepto en tiempos de crisis, la incertidumbre crea un elemento de continuidad, ya que los agentes o instituciones difícilmente modificarán su comportamiento ante fluctuaciones de todo tipo, precisamente por su vacilación ante información insuficiente o considerada insuficientemente fiable.}
    \item \textit{La importancia a las cuestiones relacionadas con la distribución del ingreso y la riqueza, y los conflictos que estas pueden generar}. Por supuesto, estas cuestiones de distribución y conflictos también se encuentran entre autores institucionalistas y marxistas. El hecho es que la cuestión de la distribución funcional del ingreso ya estaba en el corazón de la economía poskeynesiana desde sus inicios en la década de 1950.\footnote{Los poskeynesianos también consideran que la crisis financiera de finales de la década de 2000 puede ser atribuida a los cambios dramáticos en la distribución del ingreso, en particular la creciente participación en las ganancias y la creciente desigualdad en la distribución de los salarios. Además, muchos poskeynesianos consideran que la política monetaria restrictiva de los años ochenta jugó un papel importante en el aumento de la desigualdad, favoreciendo a los rentistas (SECCARECCIA, 2019). Los conflictos redistributivos juegan un papel importante, tanto en la determinación de la tasa de inflación como en la determinación de la actividad económica.}
    \item (Algunos autores) \textit{Teoría microeconómica propia}. Los economistas poskeynesianos rechazan la teoría marginalista de los precios y argumentan que los costos marginales son esencialmente constantes, lo que implica costos unitarios decrecientes hasta la capacidad total. Según los poskeynesianos, los ajustes a los cambios en la demanda en una economía de mercado ocurren principalmente a través de cantidades y no a través de precios, pues éstos varían cuando cambian los costos unitarios normales, no cuando cambia la demanda.\footnote{Además, objetan el predominio de los efectos sustitutivos neoclásicos y, en cambio, ponen sobre todo énfasis los efectos ingreso.}
\end{lnum}

\subsection{Escuela de Cambridge} 
A partir de la teoría de Keynes la Universidad de Cambridge se convertiría en un importante centro de debates económicos. Uno de esos debates se produjo en torno al concepto de función de producción.\footnote{WICKSELL y WICKSTEED formalizaron la Función de producción de Cobb-Douglas, que representa las relaciones en un sistema con dos factores de producción y un producto (que puede ser consumido o ahorrado como inversión).} 

\subsubsection{La controversia del capital}

\textsc{Cambridge. Inglaterra}

JOAN ROBINSON en \textit{La Función de Producción y la Teoría de Capital} (1953) evidenció la circularidad que la concepción del capital tenía dentro del marco de la economía ortodoxa. El debate mismo comenzó con la siguiente observación de ROBINSON:

\begin{adjustwidth}{2cm}{0pt}
    \raggedright 
    \textit{Al estudiante de teoría económica se le enseña a escribir $Y = f (L, K)$, donde $L$ es una cantidad de trabajo, $K$ una cantidad de capital y $Y$ una tasa de producción de las mercancías. Se le instruye asumir todos los trabajadores como iguales, y a medir $L$ en horas-hombre de trabajo, se le dice algo sobre el problema de números índices en la elección de una unidad de producción, y se le apresura al siguiente asunto, en la esperanza de que se olvide de preguntar en qué unidades se mide $K$. Antes de que él lo pregunte, se ha convertido en un profesor, y así los hábitos de pensamiento descuidado se transmiten de una generación a la siguiente}\\
    \raggedleft{1956}
\end{adjustwidth}

A raíz de esta observación, ROBINSON se pregunta cómo se mide el factor capital de la función de producción que utilizan los neoclásicos. La heterogeneidad de los bienes de capital impide su agregación a través de una medición física razonable. SRAFFA ya se había anticipado a esta cuestion en 1936, cuando en una carta dirigida a ROBINSON argumentaba que: \textit{(S)i uno mide el trabajo y la tierra en términos de hombres o de acres, el resultado posee un significado claro, (...) En cambio, si mides el capital en toneladas, el resultado pura y simplemente no tiene sentido}. 

Los neoclásicos agregaban el capital a través de su valor monetario, que puede calcularse como el valor presente descontado de todos los rendimientos futuros de un determinado bien de capital. ROBINSON argumenta que dicho cálculo requiere conocer anticipadamente el tipo de interés, que según la Teoría Neoclásica, vendrá determinado por la $PMg_K$. Sin embargo, el cálculo de la $PMg_K$ requiere conocer anticipadamente el valor del capital, dando lugar a una \textbf{circularidad irresoluble}.\footnote{Otra forma de aproximarse a esta conclusión de ROBINSON y SRAFFA es la siguiente:
Los neoclásicos afirmaban que un mercado competitivo fuerza a los productores a reducir al mínimo los costes de producción. Pero, dice ROBINSON:
\begin{lnum}
    \item Los costes de producción son los precios de los inputs, como los bienes de capital, que toman su valor de los productos finales;
    \item Si el precio de los productos finales determina el precio del capital; entonces,
    \item Es un argumento circular decir que el precio del capital determina el de los productos finales.
\end{lnum}

Otra forma de abordar el problema la presentó SRAFFA, que comienza notando que el problema de la medición y concepción del capital se remonta a la economía clásica. Para los clásicos el capital, dado que la inversión es equivalente al ahorro y este a su vez es consumo pospuesto, es matemáticamente equivalente a la diferencia entre todo lo que se produce y todo lo que se consume en el proceso durante un periodo dado. Pero esa diferencia es también la plusvalía o \textit{nueva riqueza} producida en un periodo dado. El problema, del que los clásicos eran conscientes, aparece cuando notamos que la inversión depende de la tasa de ganancia. Como es obvio, hay en el argumento una circularidad que, si bien necesariamente no lo invalida, lo convierte en poco útil desde la perspectiva analítica: \textbf{no podemos determinar los precios de mercado sin conocer la tasa de ganancia, pero la tasa de ganancia no puede ser determinada sin conocer los precios de mercado.} La tasa de ganancia depende de los precios de mercado, pero estos dependen a su vez de la tasa de ganancia. \\
De esta forma, no podemos valuar las inversiones y, consecuentemente, no se puede determinar el valor del capital. Tampoco lo podemos derivar lógicamente del valor del trabajo (como siendo el valor restante que queda del monto total del producto después que se pagan salarios) dado que el trabajo y el capital no son homogéneos entre sí. Así, a menos que esa derivación sea utilizada simplemente por convención o algún otro mecanismo, no hay ninguna razón aparente por la cual la diferencia monetaria entre el monto de los salarios y el del producto deba ser el costo del capital.\\
Como el establecimiento de precios dura tiempo, el precio de los bienes se determina antes de que se vendan. Para ROBINSON, una economía capitalista se caracteriza por dos sistemas de precios:
\begin{la}
    \item Uno es el de los bienes como producto, que depende en gran medida de os salarios y los beneficios; y,
    \item El segundo, el precio de los activos de capital, que depende de las expectativas de beneficios futuros, la capitalización actual y las tasas de los flujos de los beneficios. 
\end{la}

La traslación de esta Controversia del capital al campo de la distribución la plasmó SRAFFA en 1960 con \textit{Producción de mercancías por medio de mercancías}, explicando que las relaciones tecnológicas son la base de la producción de bienes y servicios. Por tanto: \textbf{no son los ajustes de mercado, si no la sustituibilidad entre salarios y beneficios, el conflicto laboral y la intervención del gobierno.}\\
En este libro también expuso un método para resolver el arduo problema con el que se enfrentó Ricardo sin éxito: la medida invariable del valor. Para ello, ideó una mercancía-patrón abstracta, formada por una cesta de bienes adecuadamente elegida. Esta mercancía-patrón era la que se elaboraba con la combinación de capital y trabajo promedio de la utilizada en la economía. No obstante, esta mercancía patrón no conservaba invariable su valor a lo largo de un amplio periodo de tiempo, porque cambiaba en cuanto un avance tecnológico variase el promedio de la relación capital-trabajo. Por eso se les conoce como \textit{neorricardianos}}

En consecuencia, no existe una forma adecuada de medir en la misma escala la diversidad de bienes de capital, invalidando la utilización de una función de producción agregada y la Teoría Neoclásica de la distribución de la renta.\footnote{Recuérdese que esta Teoría se caracterizaba por defender que la distribución iba asociada a las aportaciones productivas de los diferentes factores. Es decir, la distribución venía definida por la Productividad Marginal de los diferentes factores productivos.
} Para entender mejor el argumento de ROBINSON (y SRAFFA), supongamos una economía con $L_0$ unidades de trabajo y $K_0$ unidades de capital físico, que dan lugar a una producción de $Y_0$ unidades. Las $K_0$ unidades de capital fisico, dado un tipo de interés nominal $i_0$ , tendrían unvalor monetario $M_0$. Si ahora el tipo de interés cambiase a $i_1$, el valor monetario de las mismas $K_0$ unidades de capital fisico sería $M_1$. En consecuencia, si utilizamos el valor monetario del capital como medida para agregar el capital (cosa que defiendía la Teoría ortodoxa), tendríamos la siguiente inconsistencia: 

$$Y_0=f(M_0,L_0)=f(M_1, L_0)$$

En consecuencia, no podemos utilizar el valor monetario del capital para agregar este factor de producción, echando por tierra toda teoría que se base en el uso de una Función de Producción Agregada. La crítica de ROBINSON y SRAFFA tiene importantes raminificaciones en la teoría de la distribución de la renta, ya que los conceptos de productividad marginal del capital o del trabajo pierden todo su sentido.

\textsc{Cambridge, Masachusetts. Estados Unidos}

Los economistas neoclásicos no ven ningún gran problema en esta inconsistencia formulada por éstos economistes. 

Según ellos es suficiente agregar el valor en dinero de todos los bienes con el fin de obtener el monto total de capital. Esto es, obviamente, una simplificación pero una simplificación permisible a niveles introductorios y como mecanismo para estudiar relaciones generales, no una representación fiel de la realidad. Obviamente, esto no resuelve realmente el problema de como se generaliza al nivel nacional, lo que finalmente llevó a la sugerencia, por SAMUELSON y otros, que el capital puede ser estudiado utilizando las llamadas Funciones Surrogadas de Producción (también conocidas como Funciones de producción per cápita), que buscan relacionar la productividad de un trabajador cualquiera con la cantidad de capital que utilice, con el precio de este último determinado de acuerdo al mercado competitivo.

En la actualidad generalmente se considera que la sugerencia de las Funciones Surrogadas de Producción no es teóricamente satisfactoria pero funcionan bien a nivel empírico. 

\subsubsection{Equilibrios múltiples}
La controversia del capital se vuelve crucial para el desarrollo de la Teoría post-keynesiana en todas sus ámbitos de investigación (crecimiento, producción, etc.), tomando una posición central el concepto de dependencia histórica y la importancia de factores no puramente económicos (es decir, externamente al área formal de la disciplina).

\begin{specialblock}{Distribución en la Teoría Neoclásica}
    En la Teoría Neoclásica se supone que existe una función de producción agregada, con capital y trabajo como factores, y que el capital se puede medir en una sola unidad ($K$) independiente de los precios. De manera general, eso permite dibujar una curva bien comportada (Función de Producción neoclásica de buen comportamiento, por ejemplo) en el que capital y trabajo serán sustituibles (en diferentes grados) y esta decisión se tomará basándose en el precio de cada factor productivo:
    \begin{la}
        \item Cuanto más alto es el salario (más caro el trabajo), las empresas sustituyen trabajo por capital; 
        \item Cuanto más alto es el tipo de interés (más caro el capital), las empresas sustituyen capital por trabajo.
    \end{la} 

    De ahí sale una relación monótona entre tipo de interés y intensidad de capital y, por tanto, un único equilibrio de distribución y de técnicas elegidas.
    
\end{specialblock}

La polémica de Cambridge mostró que, en general, esto no funciona: el fenómeno de la resustitución de técnicas y \textit{reverse capital deepening}\footnote{La intensificación inversa del capital es la propiedad por la cual puede resultar eficiente asociar una tasa de interés más baja (más alta) con un capital por trabajador más bajo (más alto). Esta propiedad es incompatible con la Teoría Neoclásica de que, en virtud del principio de sustitución, las técnicas de producción más intensivas en capital se convertirán en óptimas a medida que se reduzca la tasa de interés. La intensificación inversa del capital es un ejemplo importante de las aparentes paradojas asociadas a los efectos indirectos en una economía de producción. Implica que la elección técnica no puede considerarse una función monótona del tipo de interés y cuestiona las implicaciones políticas generalizadas de la visión tradicional.
} implica que, para ciertos sistemas de producción, la relación entre tipo de interés y elección de técnicas no es monótona, y que la misma técnica puede ser óptima a tipos de interés muy distintos. Eso abre la puerta, lógicamente, a más de una combinación coherente de precios, distribución y técnicas: ya no hay garantía de unicidad ni de estabilidad del equilibrio lo que nos conduce a la existencia de \textit{equilibrios múltiples}.

Si varias configuraciones de técnicas y de distribución son compatibles con el mismo conjunto de datos tecnológicos, no tiene sentido hablar de la posición de equilibrio como si fuese un estado natural hacia el que la economía converge. KALDOR, ROBINSON y SRAFFA (al igual que muchos otros post-keynesianos) insistieron justamente en esto: los equilibrios son dependientes de la trayectoria, es decir, de la historia del sistema; el camino influye en el destino, y no hay ningún punto predeterminado al que el sistema tienda de manera automática.\footnote{Eso encaja con la lectura keynesiana de la demanda efectiva, incluso en el largo plazo, el nivel de actividad y de inversión depende de expectativas, instituciones y distribución del ingreso, no de un mecanismo de mercado que empuje a la plena utilización de recursos.
}

En estas condiciones, la elección real de técnicas ya no puede explicarse como un simple resultado de la variación de precios relativos en un modelo abstracto. Para los autores postkeynesianos, la elección de técnicas queda determinada por la historia concreta del sistema: el patrón previo de inversión, la estructura de propiedad, las instituciones laborales, las políticas industriales, los cambios en la demanda efectiva, los avances tecnológicos discretos, etc. No es algo que se pueda deducir de una Función de Producción y dos precios de factores; es un proceso histórico, conflictivo, donde la técnica es un paquete de organización social, tecnología y expectativas sobre la rentabilidad futura que se hereda. 

\imagenfit{Path-Dependancy.png}{Diagrama cíclico de técnicas de producción estado-dependientes}{\href{https://slidemodel.com/templates/path-dependence-slide-template-powerpoint/three-phase-slide-path-dependence-presentation/}{Three Phase Slide for Path Dependence Presentation}}

Podemos ver la influencia de esta concepción relativista o estado-depediente de la Función de Producción y, en especial, de la combinación Técnico-Eficiente de factores productivas en las técnicas no-paramétricas de estimación de la Frontera de Producción [\authorfont{Ver Tema 3.A.12}]:

\imagenfit{DEA_Estimations.png}{Comparación entre (a) la expansión radial y (b) la DEA basada en holguras.}{\href{https://www.mdpi.com/1999-4893/17/3/89}{\textit{On the Use of Data Envelopment Analysis for Multi-Criteria Decision Analysis}. PASCOE (2024)}}

En el ámbito del cremiento económico, ROBINSON expone un sencillo modelo de crecimiento económico, expuesto en \textit{La acumulación de capital}. Consideremos un sistema económico de dos sectores: el de consumo ($C$) y el de inversión ($I$). La suma de ambas macromagnitudes es la renta nacional, o producción anual ($Y$); así, la ecuación de la renta es: $Y=C+I$. La renta de pleno empleo se tipifica a la unidad (YPE=1) y de ahí se tiene una ecuación de balance: $C+I=1$, o recta de balance en su representación gráfica. Además de éstos, hay otro supuesto básico: el consumo es una función lineal y creciente de la inversión; la función de consumo es: $C= C_0+\alphaI$, siendo $C_0$ y $\alpha$ variables exógenas. 

\imagenfit{Robinson_Crecimiento.png}{Modelo de creciminto a largo plazo. ROBINSON, 1956}{E. ESCARTÍN, \textit{Apuntes de Historia del Pensamiento Económico. Tema 27}}

El senda de crecimiento posible ($C_2DFC_1$), es decir, combinación de $(C,I)$, está acotado por un conjunto de restricciones:\footnote{La barrera de inflación no es propiamente una restricción pero nos indica que una vez alcanzado el nivel óptimo (pleno empleo), cualquier aumento de la inversión solo conducirá a un aumento de la inflación, ya que la economía real no puede crecer más (está en pleno empleo)}
\begin{lnum}
    \item Restricción de balance, que determina el pleno empleo;
    \item Restricción de supervivencia, que determina la cantidad mínima de consumo necesario;
    \item La restricción de capital, que determina el límite a partir del cual el consumo de supervivencia y el consumo autónomo ($C_0$) agregado no pueden superar sin que exista una cantidad mínima de inversión
\end{lnum}

\paragraph{La Hipótesis de la Inestabilidad Financiera: MINSKY} 
Otro referente de la escuela postkeyenesiana es el estadounidense MINSKY, quien en 1986 publicó el libro \textit{Estabilizando una Economía Inestable}. En esta obra desarrolla la \textit{Hipótesis de la Inestabilidad Financiera}, según la cual, la fragilidad intrínseca del sistema financiero será la principal responsable de los ciclos económicos en una economía capitalista. 

\imagenfit{CicloMinsky.png}{(Micro) Decisiones de inversión en la Empresa (Financiación Interna/Financiación Externa) | (Macro) Ciclo de MINSKY en la economía real por decisiones de sobreinversión}{(Micro) \href{https://en.wikipedia.org/wiki/Minsky_moment}{Minsky moment (Wikipedia.org)} | (Macro) \href{https://www.tandfonline.com/doi/abs/10.2753/PKE0160-3477340301}{\textit{Post-Keynesian modelling: where are we,
and where are we going to?}. C. SEBASTIEN et. al, 2010}}

El ciclo de MINSKY puede caraterizarse de la siguiente forma:
\begin{la}
    \item \textit{A nivel microeconómico.}\footnote{El precio de demanda de los bienes de capital, $P_K$, depende de los flujos de caja previstos por la empresa. El precio de oferta de estos bienes, $P_I$, viene dado por un margen sobre los costes laborales unitarios. El importe de la inversión que una empresa puede financiar internamente está limitado por sus flujos de caja previstos, lo que da lugar a la ecuación $P_II_{int}=\Pi^e$ (se supone que la hipérbola representa beneficios constantes).} Tan pronto como una parte de la inversión se financia con fondos externos, $I_x$, aparecen dos riesgos:
    \begin{la}
        \item Por un lado, el riesgo del prestamista recae sobre el sistema bancario y aumenta el precio de oferta de los bienes de capital ($P_K$); y,
        \item Por otro lado, a medida que aumenta el endeudamiento, las empresas se vuelven más frágiles, ya que tienen que gestionar beneficios esperados inciertos y gastos financieros ciertos. Por lo tanto, el riesgo del prestatario disminuye el precio de demanda de los bienes de capital ($P_K$) cuando se incrementa el importe de la financiación externa, es decir, el endeudamiento.
    \end{la} 
    Durante un auge, aparece la fragilidad financiera porque el riesgo del prestatario y el riesgo del prestamista se vuelven imprudentemente bajos, lo que aumenta la proporción de financiación externa con respecto a la financiación interna y la carga de la deuda. En este punto, MINSKY aplica al nivel macroeconómico los resultados anteriores obtenidos a nivel microeconómico;
    \item \textit{A nivel macroeconómico.} Partiendo de un momento de estabilidad, el sistema financiero comienza a deteriorarse cuando el exceso de confianza propicia una relajación de riesgos que desemboca en un sobreendeudamiento general. Tarde o temprano, los riesgos excesivos asumidos en la fase de auge se traducen en impagos y liquidaciones masivas de activos que provocan una crisis financiera. Ante esta situación, las instituciones financieras reaccionan reduciendo el crédito, trasladando así la crisis a la economía real.
\end{la}

Algunas voces consideran que tanto la Gran Depresión del 29 como la Gran Recesión de 2008 validan empíricamente esta hipótesis.\footnote{Para un desarrollo más elaborado [\authorfont{Ver Anexo \ref{anx:minsky}}]}

\subsubsection{Desequilibrio}
Los economistas postkeynesianos no fueron los únicos keynesianos en criticar la teoría macroeconómica de la síntesis neoclásica. La separación abrupta de la economía entre microeconomía y macroeconomía daba lugar a una inconsistencia teórica, que además impedía explicar parte de la evidencia empírica observada. Los orígenes esta línea de investigación se suelen establecer en 1965 con la publicación de la \textit{Contrarrevolución Keynesiana},\footnote{Con este título, CLOWER critica la contrarrevolución científica que significó la síntesis neoclásica. En aquellos momentos, la teoría macroeconómica de la síntesis neoclásica se asimilaba bajo el término \textit{paradigma keynesiano}. CLOWER y otros autores consideraban que este \textit{paradigma keynesiano} no tenían mucho que ver con KEYNES.} de R. CLOWER, y con la publicación en 1968 del libro de A. LEIJONHUFVUD, \textit{Sobre la Economía Keynesiana y la Economía de Keynes}. 

Según estos autores, la verdadera esencia de la teoría keynesiana es el estudio de los efectos que provoca el desequilibrio de unos mercados sobre el comportamiento de los agentes en otros mercados. 

En la teoría neoclásica walrasiana, este escenario de desequilibrio está descartado gracias a la existencia del subastador walrasiano, que impide la realización de intercambios en cualquier mercado antes de haber alzancado un vector de precios de equilibrio general.\footnote{En la teoría neoclásica walrasiana se estudian las propiedades del vector de precios de equilibrio general, esto es el vector de precios que equilibra simultáneamente todos los mercados. Véase [\authorfont{Ver Tema 3.A.21}]} Los keynesianos del desequilibrio, por el contrario, se centran en estudiar el comportamiento real de una economía descentralizada sin dicho subastador, en la que los agentes económicos realizan intercambios con precios que no se corresponden con los de equilibrio general. 

Los estudios de estos autores concluyen que un mercado puede encontrarse en equilibrio o en desequilibrio independiente de que el precio fijado para tal mercado se corresponda o no con el precio que prevalecería en equilibrio general. En estas circunstancias, puede ocurrir perfectamente que exista desempleo involuntario aunque el salario real coincida con su valor de equilibrio general, o que el mercado de trabajo esté en equilibrio con un salario real distinto al que prevalecería en equilibrio general. Dicho de otra forma: el problema del desempleo involuntario (o de cualquier mercado) es el resultado de un problema de coordinación, surgido por la ausencia de un subastador en una economía descentralizada. 

En estas circunstancias, el problema del desempleo y el desequilibrio obedece, más que a la rigidez de los precios y los salarios, a la falta de coordinación a la hora de modificar dichos precios y salarios en aras de alcanzar el equilibrio general. 

Siguiendo la publicación de BARRO y GROSSMAN (1971), se expone un modelo de Desequilibrio General de la Producción y el Empleo con Exceso de Oferta, que explica estas nociones. En el Anexo 4. se puede consultar la formalización de las seis ecuaciones básicas del modelo, representadas en el Gráfico 3. 

\imagenfit{DesequilibrioGeneral.png}{El Modelo de Desequilibrio General con Exceso de Oferta}{Apuntes ICEX-CECO. Tema 3.A.4}

El Gráfico muestra el mercado de trabajo (izquierda) y el mercado de bienes de consumo (derecha). El mercado de trabajo relaciona la demanda y la oferta de trabajo con el salario real. En este modelo, la oferta de trabajo recibirá el nombre de \textit{oferta nocional de trabajo}, representada por $L^s$. Por el lado de la demanda, encontramos la tradicional demanda de trabajo con pendiente negativa $L^d$, que en este caso recibe el nombre de \textit{demanda nocional de trabajo}. La demanda nocional de trabajo relaciona la cantidad demandada de trabajadores para cada posible salario real, suponiendo que las empresas consiguen vender la totalidad de su producción al precio vigente de los bienes que producen. La demanda nocional de trabajo y la oferta nocional de trabajo proporcionan un salario real $W^*$ de equilibrio general walrasiano, que da lugar a un volumen de empleo $L^*$. El equilibrio general walrasiano será entonces el punto A. 

Supongamos ahora que los bienes de consumo se están vendiendo a un precio excesivamente alto, $P_1$, superior al precio de equilibrio general walrasiano $P^*$. En este caso, las empresas no conseguirán vender toda la producción asociada al volumen de empleo $L^*$. En su lugar, venderán una cantidad de bienes que tan solo requiere la contratación $L_1$ trabajadores. La línea vertical $L_1$ será la \textit{demanda efectiva de trabajo} asociada al precio $P_1$. Para cada precio, habrá una demanda efectiva de trabajo, en concreto:
\begin{la}
    \item Cuanto menor sea el precio del bien, mayores serán las ventas de bienes de consumo y por tanto, mayor será la demanda efectiva de trabajo. El salario real no tendrá ningún efecto sobre la demanda efectiva, aunque habrá un límite superior por encima del cual las empresas nunca contratarán un determinado volumen de empleo, que es prcisamente la demanda nocional de trabajo, es decir, su productividad marginal (como en el Equilibrio General); de la misma forma,
    \item Cuanto mayor sea el precio del bien, menores serán las ventas y, por tanto, la demanda efectiva de trabajo.
\end{la} 

El mercado de bienes muestra también una \textit{oferta nocional de bienes de consumo} $Y^s$, que tendrá pendiente negativa con respecto al salario real:
\begin{lnum}
    \item A medida que sube el salario real;
    \item Se reduce la demanda nocional de trabajo; y, por tanto,
    \item Obtenemos un menor nivel de producción; dando lugar a,
    \item Una reducción del bien de consumo.
\end{lnum}

La \textit{demanda nocional de bienes de consumo} tendrá pendiente positiva con el salario real ($Y^d$), porque:
\begin{lnum}
    \item La subida del salario real produce:
    \begin{la}
        \item Un efecto renta en los consumidores. A mayor salario real, más capacidad de compra y más consumo; y,
        \item Un efecto sustitución. Al subir el salario, los consumidores deciden renunciar a tiempo de ocio a cambio de un mayor consumo; y, por tanto,
    \end{la}
    \item Se incrementa la demanda de consumo de los hogares.
\end{lnum}

En efecto, la subida de precios reduce el valor real de los saldos monetarios mantenidos por los hogares, reduciendo su capacidad de compra. Por este motivo, la demanda nocional de consumo asociada a $P_1$ inferior, para cada salario real $W$, a la que prevalecería con los precios de equilibrio general walrasiáno $P^*$. Recordemos ahora, que con el precio $P_1$, las empresas no demandan el nivel de empleo $L^*$, sino $L_1$. Dicho nivel de empleo determina, junto con el nivel de precios $P_1$, la \textit{demanda efectiva del bien de consumo}, $Y_1$. Esta demanda efectiva será independiente del salario real, encontrando un límite inferior en la demanda nocional del bien de consumo. Por tanto, para un $P_1 > P^*$:
\begin{la}
    \item Ante una reducción del salario nominal. La demanda efectiva de trabajo $L_1$ producirá desempleo involuntario aunque el salario real sea $W$. Tras la reducción del salario nominal, sin reducir el precio del bien (es decir, si bajamos tan solo el salario nominal) nos movemos al punto C, donde no existe desempleo involuntario, pero existe un exceso de oferta del bien de consumo: la demanda efectiva del bien sería inferior a la oferta nocional del bien; por otro lado,
    \item Ante un aumento del salario nominal. Manteniendo constante el precio (es decir, si subimos tan solo el salario nominal), nos iriamos al punto D, en el que existe desempleo involuntario, pero el mercado de bienes está en equilibrio, ya que la demanda efectiva del bien se iguala a su oferta nocional.
\end{la}

Como vemos, la ausencia de un subastador que modifique el salario real y el precio del bien de manera coordinada puede perpetuar una situación de desequibrio en alguno de los mercados.\footnote{En el punto B, el salario real es el adecuado, pero el precio del bien es excesivamente elevado. En esta situación, un subastador reduciría el precio del bien, pero al mismo tiempo se aseguraría de reducir el salario nominal en la misma proporción, manteniendo inalterado el salario real. En ausencia de un subastador, los agentes interactuan en los mercados de manera descentralizada sin esperar a que se alcacance el par de precios de equilibrio general walrasiano $(P^*, W^*)$. Esta ausencia de coordinación podría llevar a que, partiendo del punto B, las empresas bajasen el precio del bien, provocando al mismo tiempo un aumento del salario real, manteniendo el desempleo involuntario.
} Estas situaciones de \textit{equilibrio-desequilibrio} presentan ciertas particularidades y los conocemos como:
\begin{la}
    \item \textit{Desempleo keynesiano}. En el punto B existe exceso de oferta en ambos mercados y la causa no es un salario real excesivo sino una demanda insuficiente;
    \item En el punto C existe un exceso de oferta en el mercado de bienes que convive con equilibrio en el mercado de trabajo; y,
    \item \textit{Desempleo clásico}. En el punto D, el mercado de bienes está equilibrado y el mercado de trabajo presenta exceso de oferta, cuya causa se encuentra en un salario real excesivo.
\end{la}

Los keynesianos del desequilibrio tuvieron una gran influencia en la década de los 70's, especialmente en Estados Unidos y en Francia. Sin embargo, el surgimiento de la Nueva Macroeconomía Clásica terminó por eclipsar estas contribuciones.\footnote{En efecto, el mismo BARRO acabó engrosando las filas de la Nueva Macroeconomía Clásica}.

\clearpage
\section{Conclusión}

\subsection{Recapitulación}



\subsection{Extensiones y relación con otras partes del Temario}



\subsection{Opinión}

\subsubsection{Idea final (Salida o cierra)}

\clearpage
\pagenumbering{gobble}
\MyBibliography
\MyBibItem{M. LAVOIE y M. SECCARECCIA}{La economía poskeynesiana, ¿un
pensamiento heterodoxo desconocido?}{2025}{}


% --- Anexos (no aparecen en el índice) ---
\clearpage
\anexo{\textbf{Preguntas Test}}
%\input{\TemaCode/questions.tex} 


\clearpage
\anexo{MICHAEL KALECKI: Demanda Efectiva}{\label{anx:kalecki}}
En 1933, Kalecki publicó un artículo titulado \textit{Sobre el ciclo económico} en el que se anticipó a la teoría de la demanda efectiva de Keynes. Por esta razón, algunos de sus seguidores consideran que la escuela postkeynesiana debería llamarse, en realidad, escuela kaleckiana. 

De formación marxista, Kalecki rechaza la teoría neoclásica de la distribución de la renta, según la cual, los factores productivos se remuneran de acuerdo a su productividad marginal; y considera que la participación del trabajo y el capital en la renta depende del grado de monopolio de los capitalistas. 

Además, sustituye la función de producción neoclásica, con rendimientos marginales decrecientes del trabajo ($\downarrow PMG_L$), por una función con rendimientos constantes (esto es así porque las empresas suelen operar con exceso de capacidad instalada). Al descartar los rendimientos marginales decrecientes, Kalecki niega la noción neoclásica, aceptada por Keynes, de que todo aumento del empleo debe venir acompañado de una reducción del salario real. 

Según Kalecki, la distribución de la renta es un elemento central en la determinación de la demanda efectiva, y por tanto, de la producción y el empleo. Kalecki asume que los trabajadores consumen prácticamente toda la renta que perciben, mientras que los capitalistas solo gastan una parte.\footnote{De este supuesto se obtiene la famosa noción kaleckiana de que \textit{los capitalistas ganan lo que gastan y los trabajadores gastan lo que ganan}.\\ El supuesto de que los trabajadores gastan toda su renta lleva a Kalecki a defender que un aumento del salario real incrementaría el empleo. En su enfoque, un cambio en la distribución de la renta a favor de los trabajadores aumentaría la demanda agregada y con ella el empleo agregado de la economía.} En consecuencia, la demanda de trabajo dependerá positivamente del salario real, justo al contrario que en el esquema neoclásico. Este atributo kaleckiano de la demanda de trabajo se conoce como \textit{la paradoja de los costes de Kalecki}\footnote{\textit{Lo que es ventajoso para un empresario aislado no lo es necesariamente para el conjunto de los empresarios considerados como clase}. Kalecki, \textit{Ensayos escogidos sobre dinámica de la economía capitalista} (1971).\\
Para su desarrollo formal [\authorfont{Ver Anexo \ref{}]}}

Actualmente, en aras de simplificar, se podría decir que los trabajos poskeynesianos sobre la economía real se inspiran más en Kalecki, mientras que los trabajos en finanzas y dinero está más inspirado en Keynes, aunque los escritos de Kalecki sobre el dinero también son bastante compatibles con teorías monetarias poskeynesianas contemporáneas.

\clearpage
\anexo{ROY HARROD: Demanda Efectiva y crecimiento económico}\label{anx:harrod}
Debemos también destacar la influencia de ROY HARROD en el pensamiento económico post-keynesiana por ser pionero en la formalización de un modelo de crecimiento incorporando la demanda efectiva como variable relevante. Este será un punto de choque frontal entre la corriente post-keynesiana y la síntesis neoclásica, por considerar los segundos que la Demanda Efectiva no tiene cabida en el largo plazo (he aquí la noción de \textit{síntesis neoclásica}, pues conciliarán las ideas de KEYNES y la Teoría neoclásica de forma intertemporal: KEYNES c/p y Neoclásicos l/p). Harrod formó parte del entorno académico de Cambridge, gran amigo de KEYNES,\footnote{Se conocieron en el \textit{King's College} alrededor de 1920.\\
Tras la muerte de su amigo, en 1946, Harrod y Austin Robinson (marido de JOAN ROBINSON) escribieron un extenso obituario de Keynes para The Economic Journal. A instancias de Geoffrey Keynes (hermano de J.M. KEYNES), Harrod emprendió entonces la tarea de redactar una biografía importante de éste: \textit{The Life of John Maynard Keynes}, que fue publicada con gran reconocimiento en 1951, en un momento en que la mayoría de los familiares y amigos de Keynes aún estaban vivos.} y fue uno de los primeros economistas en formalizar las ideas de la Teoría General, por ello es considerado como un de los principales autores proto-post-keynesiano (junto a KALECKI).

En 1939, Harrod publicó su \textit{Ensayo de la Teoría Dinámica}, en el que extendía la preocupación de Keynes por la inestabilidad macroeconórnica al campo del crecimiento económico. Frente a los autores de la Síntesis Neoclásica que defenderían la introducción de los supuestos neoclásicos para los modelos de largo plazo. Harrod lleva el modelo keynesiano al largo plazo simplemente introduciendo dinámica en un modelo con rigideces. Harrod va a distinguir tres tasas de crecimiento:
\begin{lnum}
    \item La tasa \textit{efectiva} a la que crece la economía;
    \item La tasa \textit{natural}, a la que necesita crecer la economía para mantener el pleno empleo, y que vendría determinada por la suma del crecimiento de la población activa y la tecnología; y,
    \item La tasa \textit{garantizada}, que garantiza el cumplimiento de las expectativas de las empresas.
\end{lnum} 

La tasa garantizada viene determinada por la \textit{Ecuación Fundamental}:\footnote{El Para una demostración detallada de como se llega a esta ecuación [\authorfont{Ver Anexo }\ref{}]}

$$g^*=\frac{s}{c}\qquad \text{donde,}\;
\begin{cases}
    g\equiv\quad \text{Tasa de crecimiento garantizada};\\[0.5em]
    s\equiv\quad \text{Tasa de ahorro:}\; s=\frac{S}{Y}\;;\\[0.5em]
    c\equiv\quad \text{Ratio Marginal Capital-Producto:}\; c=\frac{\partial K}{ \partial F}\\[0.5em]
\end{cases}$$

Por lo tanto, el modelo de HARROD se plantea como sigue:
\begin{la}
    \item Si la tasa efectiva de crecimiento es superior a la tasa garantizada, $g$, las empresas habrán producido más de lo previsto y, dada una ratio $c$, tendrán que incrementar sus inversiones, dando lugar al llamado \textit{efecto acelerador} de la inversion:
    \begin{lnum}
        \item El incremento de la inversión incrementa la producción, por el aumento de la demanda; y,
        \item El incremento de la producción incrementa la inversión, por el deseo de las empresas de mantener una determinada relación capital-producto.
    \end{lnum}
    Ello da lugar a un círculo virtuoso (conceptualmente similar al del efecto multiplicador), que aumentaría la tasa de crecimiento hasta alcanzar a la tasa natural.
    \item Si la tasa efectiva de crecimiento cayese por debajo de la tasa garantizada, la economía entraría en una espiral recesiva que sólo acabaría si si cambiase el valor de los parámetros que determinan la tasa garantizada, por algún motivo (por ejemplo, el ahorro); y,
    \item Si la tasa efectiva de crecimiento fuese igual a su tasa garantizada, se encontraría en el \textit{filo de la navaja}, es decir, a un paso de la expansión o de la recesión.
\end{la}

El modelo de Harrod también permite la existencia de desempleo involuntario persistente, siguiendo a KEYNES. En efecto, incluso aunque la economía esté creciendo de manera estable a la tasa garantiza, nada garantiza que ésta coincida con la tasa de crecimiento natural. Todo ello justifica que la política macroeconómica trate de asegurar un nivel adecuado de demanda.\footnote{Las conclusiones de Harrod fueron rebatidas posteriormente por los economistas neoclásicos a través de la elaboración de modelos de crecimiento diferentes. En 1956, SOLOW, gran exponente de la síntesis neoclásica, publicó su importantísimo artículo titulado \textit{Una Contribución a la Teoría del Crecimiento Económico}, en el que inicia la familia de modelos de crecimiento basados en el uso de funciones agregadas de producción. SOLOW demuestra que, utilizando el mismo modelo que Harrod, pero modificando la tecnología (empleando una función de producción en la que los factores de producción son sustituibles), la economía seguirá una senda de crecimieno equilibrada hasta alcanzar el estado estacionario; en el cual, el crecimiento sería constante. De esta forma, el modelo de Harrod solo sería válido en el caso extremo de que los factores de producción fueran perfectamente complementarios. La respuesta postkeynesiana a los modelos neoclásicos de crecimiento vino de la mano de J. ROBINSON y SRAFFA, que se enfrentaron a SOLOW y SAMUELSON en la \textit{Controversia de Cambridge}, que analizamos un poco más adelante. Éste es uno de los debates más encedidos en el Pensamiento Económico del s.XX.} 

\clearpage
\anexo{Hipótesis de la Inestabilidad Financiera: MINSKY}\label{anx:minsky} 
MINSKY le da mucha importancia al impacto macroeconómico que tienen las decisiones de inversión y financiación de los agentes económicos, distinguiendo tres posibles tipos de finanzas: 
\begin{lnum}
    \item \textit{Finanzas de cobertura}. Son aquellas en las que los retornos de las inversiones de los inversores endeudados son capaces, en cada periodo, de pagar los intereses de su deuda viva y devolver parte del principal. Por ejemplo, si un inversor compra un inmueble y lo pone en alquiler, se encontrará realizando finanzas de cobertura si la renta que cobra cada mes es suficiente para pagar los intereses del préstamo que pidió para financiar su compra y devolver una parte de la deuda;
    \item \textit{Finanzas especulativas.} Son aquellas en las que los retornos de las inversiones solo permiten cubrir los intereses de las deudas adquiridas para financiar dichas inversiones. En este caso, el inversor necesitará refinanciar la deuda durante el tiempo necesario hasta que las inversiones realizadas se revaloricen lo suficiente como para devolver el principal del préstamo. Por ejemplo, un inversor que se endeuda para crear una start-up cuyos beneficios de explotación tan solo cubren el coste financiero de la deuda, se encontrará realizando finanzas especulativas. En efecto, el inversor está especulando con que tarde o temprano la start-up se revalorizará lo suficiente como para venderla a un precio que permita devolver el principal del préstamo y obtener una ganancia. Hasta ese momento, tendrá que refinanciar permanentemente los pagos del principal que vayan venciendo;
    \item \textit{Finanzas Ponzi}. Un inversor se encuentra realizando un juego de Ponzi si los rendimientos de sus inversiones ni siquiera cubren el pago de interés. En este caso, el inversor deberá aumentar su endeudamiento exponencialmente. 
\end{lnum}

Para Minsky, la estabilidad macroeconómica dependerá del peso que tengan las modalidades anteriores en el total de las inversiones de la economía. Si dominan las finanzas de cobertura, la economía tenderá a ser estable, pero a medida que ganan peso las finanzas especulativas y las finanzas Ponzi, el sistema económico tendrá más probabilidad de sufrir una crisis financiera. 

Para el autor, el peso de cada tipo de finanzas es endógeno y depende de la fase del ciclo económico:
\begin{lnum}
    \item En épocas de recesión, las instituciones financieras solo financian operaciones de cobertura. 
    \item En la fase de recuperación, los mercados financieros se estabilizan y las instituciones financieras comienzan a otorgar financión para inversiones más arriesgadas, propiciando un incremento del peso de las finanzas especulativas y Ponzi. 
    \item Cuando las finanzas especulativas y de tipo Ponzi alcanzan un nivel crítico, el conjunto de los inversores será muy vulnerable a una caída en el valor de sus activos, aunque sea modesta.
\end{lnum} 

Es justo en este momento cuando nos encontramos en lo que P. McCULLEY denominó \textit{Momento Minsky} (1998), en el contexto de la crisis financiera de Rusia. 

En el Momento Minsky, cualquier caída repentina del precio de los activos puede desembocar en un aumento de la morosidad. De producirse dicha situación, las instituciones financieras endurecerían las condiciones crediticas y retirarían financiación para proyectos de inversión arriesgados. Este corte crediticio reducirá las inversiones, y con ellas, el precio de los activos, aumentando de nuevo la morosidad. Este proceso puede continuar hasta hacer colapsar el sistema financiero. La restricción del crédito afectará a la economía real, provocando un aumento del desempleo y una caída de la producción. Sin embargo, gracias a esta restricción del crédito las finanzas de cobertura volverán a aumentar su peso, aumentando resiliencia financiera del conjunto de los inversores. La estabilización del sistema financiero será, paradójicamente, la causa del siguiente ciclo.

\end{document}